% Leçon 21 — Vocabulaire et compréhension de texte : démocratie — Chinois Tle A — Cours signé OnBuch+
% Programme officiel MINESEC. Compiler avec : tectonic ZH21-vocabulaire-democratie.tex
\newcommand{\DOCMATIERE}{Chinois}
\newcommand{\DOCNIVEAU}{Tle A}
\newcommand{\DOCMODULE}{Module 3 : Citoyenneté et ouverture au monde}
\newcommand{\DOCLECON}{ZH21}
\newcommand{\DOCTITRE}{Leçon 21 — Vocabulaire et compréhension de texte : démocratie}
% OnBuch+ — préambule « cours signés » ENRICHI (compilé avec tectonic / XeLaTeX)
% Système visuel « Pop » d'OnBuch+ : fond crème (jamais blanc), bordures encre
% épaisses, ombres « dures » décalées, titres Archivo Black en capitales.
% Version MATHÉMATIQUES : graphes (pgfplots/tikz), boîtes pédagogiques
% (définition, propriété, méthode, attention, le savais-tu, exercice résolu,
% exercice, TP, bilan), page de garde et signature OnBuch+ ancrée en bas.
%
% Macros attendues dans main.tex AVANT \input{preamble} :
%   \DOCMATIERE  \DOCNIVEAU  \DOCMODULE  \DOCLECON (numéro)  \DOCTITRE
\documentclass[11pt]{article}

\usepackage{fontspec}
\usepackage[french]{babel} % français = langue principale ; le chinois via \zh{..} / textebox (xeCJK)
\usepackage{xeCJK}
\usepackage{pifont} % \ding{51} \ding{55} utilisés par les modèles
\usepackage[a4paper,margin=2cm,top=3.4cm,bottom=2.8cm,headheight=1.4cm,footskip=1.3cm]{geometry}
\usepackage{amsmath,amssymb,amsfonts}
\usepackage{mathtools} % \xmapsto, \coloneqq... souvent utilisés par les modèles
\usepackage{cancel} % simplifications barrées (bilans de forces)
% \abs{z} (module, valeur absolue) et \norm{u} (norme), utilisés par les modèles
\providecommand{\abs}{}\let\abs\relax\DeclarePairedDelimiter{\abs}{\lvert}{\rvert}
\providecommand{\norm}{}\let\norm\relax\DeclarePairedDelimiter{\norm}{\lVert}{\rVert}
\usepackage[table]{xcolor} % [table] : fournit aussi \rowcolors (alternance de lignes)
\usepackage{pagecolor}
\usepackage{tikz}
\usetikzlibrary{arrows.meta,positioning,calc,shapes.geometric,shapes.misc,shapes.arrows,decorations.pathmorphing,decorations.pathreplacing,decorations.markings,patterns,shadows,mindmap,trees,fit,backgrounds,matrix,intersections,angles,quotes}
\usepackage{pgfplots}
\usepackage[version=4]{mhchem} % équations nucléaires \ce{^{235}_{92}U}
\usepackage[european,siunitx]{circuitikz} % schémas électriques
\pgfplotsset{compat=1.18}
\usepgfplotslibrary{fillbetween,groupplots} % aires entre courbes (calcul intégral)
\usepackage{enumitem}
\usepackage{fancyhdr}
% [most] charge toutes les bibliothèques tcolorbox (listings, minted, raster,
% documentation...) et gonfle inutilement la mémoire TeX sur un document long
% (~300 boîtes) : on ne charge que ce qui est réellement utilisé.
\usepackage[skins,breakable]{tcolorbox}
\usepackage{graphicx}
\usepackage{colortbl}
\usepackage{booktabs}
\usepackage{tabularx}
\usepackage{diagbox} % case d'en-tête diagonale des tableaux à double entrée
% Colonnes extensibles centrée / gauche / droite (C, L, R), souvent utilisées par les modèles
\newcolumntype{C}{>{\centering\arraybackslash}X}
\newcolumntype{L}{>{\raggedright\arraybackslash}X}
\newcolumntype{R}{>{\raggedleft\arraybackslash}X}
\usepackage{array}
\usepackage{multicol}
\usepackage{multirow}
\usepackage{siunitx}
% parse-numbers=false : accepte aussi \SI{2,5 \times 10^{-3}}{...}
\sisetup{locale=FR,output-decimal-marker={,},group-minimum-digits=4,parse-numbers=false}
\DeclareSIUnit\ohm{\ensuremath{\symup{\Omega}}} % U+2126 absent de la police
\DeclareSIUnit\division{div} % réglages d'oscilloscope (ms/div, V/div)
\DeclareSIUnit\tr{tr} % tours (vitesse de rotation, ex. \SI{1500}{\tr\per\minute})
\DeclareSIUnit\molar{M} % concentration molaire (ex. \SI{3}{\molar}, \SI{1}{\micro\molar})
\DeclareSIUnit\an{an} % année (le modèle confond parfois avec \year, un compteur TeX) — ex. \SI{5}{\kilo\meter\per\an}
\DeclareSIUnit\FCFA{FCFA} % franc CFA (ex. \SI{500}{\FCFA\per\kilo\gram})
\DeclareSIUnit\degreeBrix{\textdegree Brix} % teneur en sucre d'un sirop/jus (réfractométrie)
\DeclareSIUnit\month{mois} % le modèle utilise parfois \month, un compteur TeX, comme unité de durée
\usepackage{qrcode}
% \degree : commande fréquemment utilisée par les modèles pour un angle ou une
% température en dehors de \SI{}{} (non fournie par les paquets chargés).
\providecommand{\degree}{^{\circ}}
% \qed (fin de démonstration), utilisé par les modèles sans amsthm
\providecommand{\qed}{\hfill$\square$}
% \barre : double barre des valeurs interdites dans un tableau de variations (tkz-tab)
\providecommand{\barre}{\ensuremath{\|}}
% environnement proof (amsthm non chargé) utilisé par les modèles
\ifdefined\proof\else\newenvironment{proof}[1][Démonstration]{\par\noindent\textit{#1.}\ }{\qed\par}\fi
\usepackage{lastpage}
\usepackage{hyperref}
\hypersetup{hidelinks}

% ============================================================
% Palette « Pop » OnBuch+ — mêmes hex que la classe P (plus_ui.dart)
% ============================================================
\definecolor{poporange}{HTML}{F59321}
\definecolor{popdark}{HTML}{E07A0C}
\definecolor{popcream}{HTML}{FFF6E8}
\definecolor{popink}{HTML}{1C1714}
\definecolor{poppurple}{HTML}{7A5AE0}
\definecolor{popgreen}{HTML}{2E9E6A}
\definecolor{poppink}{HTML}{E0457B}
\definecolor{popblue}{HTML}{2F7DE1}
\definecolor{popgold}{HTML}{C98A12}
\definecolor{popmuted}{HTML}{8A7F76}
\definecolor{popline}{HTML}{EADFD2}
\definecolor{popgray}{HTML}{8A7F76}
\colorlet{popgrey}{popgray}
% Alias tolérants (le modèle nomme parfois les couleurs différemment)
\colorlet{popwhite}{white}
\colorlet{popblack}{popink}
\colorlet{popred}{poppink}
\colorlet{popyellow}{popgold}
% Teintes claires (fonds de tableaux/encadrés) — le modèle nomme parfois
% les couleurs avec un suffixe L (ex. popblueL) sans les avoir définies.
\colorlet{poporangeL}{poporange!20}
\colorlet{popdarkL}{popdark!20}
\colorlet{poppurpleL}{poppurple!20}
\colorlet{popgreenL}{popgreen!20}
\colorlet{poppinkL}{poppink!20}
\colorlet{popblueL}{popblue!20}
\colorlet{popgoldL}{popgold!20}
\colorlet{popredL}{popred!20}
\colorlet{popgrayL}{popgray!20}
% Teintes claires pour fonds de tableaux / zones
\colorlet{popblueL}{popblue!10!white}
\colorlet{poporangeL}{poporange!14!white}
\colorlet{popgreenL}{popgreen!12!white}
\colorlet{poppurpleL}{poppurple!12!white}
\colorlet{poppinkL}{poppink!12!white}
\colorlet{popgoldL}{popgold!14!white}

\pagecolor{popcream}

% ============================================================
% Typographie
% ============================================================
\defaultfontfeatures{Ligatures=TeX}
\setmainfont{PlusJakartaSans-Medium.ttf}[Path=../../../fonts/,BoldFont=PlusJakartaSans-Bold.ttf]
% Symboles, API et emojis absents de la police principale : repli sur DejaVu Sans (caractères actifs)
\newfontface\symfont{DejaVuSans.ttf}[Path=../../../fonts/]
\makeatletter
\newcommand\symactive[1]{\catcode#1=13 \begingroup\lccode`\~=#1 \lowercase{\endgroup\def~}{{\symfont\char#1}}}
\newcommand\symblank[1]{\catcode#1=13 \begingroup\lccode`\~=#1 \lowercase{\endgroup\def~}{}}
\newcommand\symhyph[1]{\catcode#1=13 \begingroup\lccode`\~=#1 \lowercase{\endgroup\def~}{\mbox{-}}}
\newcount\symc
\newcommand\symrange[2]{\symc=#1 \@whilenum\symc<#2 \do{\expandafter\symactive\expandafter{\the\symc}\advance\symc by 1 }}
\newcommand\symblankrange[2]{\symc=#1 \@whilenum\symc<#2 \do{\expandafter\symblank\expandafter{\the\symc}\advance\symc by 1 }}
\makeatother
\symrange{"0250}{"02FF}\symactive{"2713}\symactive{"2714}\symactive{"2717}\symactive{"2718}\symactive{"2610}\symactive{"2611}\symactive{"2612}
\symrange{"2600}{"27BF}
\catcode"2705=13 \begingroup\lccode`\~="2705 \lowercase{\endgroup\def~}{{\symfont\char"2713}}\symblank{"26BD}\symblank{"26A1}
\symrange{"2460}{"257F}\symactive{"223C}\symactive{"2248}\symactive{"2260}
\symactive{"01F8}\symactive{"01F9}\symactive{"1E3E}\symactive{"1E3F}
\symhyph{"2011}\symhyph{"2E1A}\symblankrange{"1F300}{"1FAFF}\symblank{"FE0F}
% Caractères chinois : WenQuanYi Zen Hei (sous-ensemble GB2312 + pinyin), gras/italique simulés
\setCJKmainfont{WenQuanYiZenHei-subset.ttf}[Path=../../../fonts/,AutoFakeBold=3,AutoFakeSlant=0.2]
\xeCJKsetup{CJKspace=false}
% Pinyin : voyelles à caron (ǎ ǐ ǒ ǔ ǖ ǘ ǚ ǜ) absentes de la police principale -> caractères actifs
\newfontface\pyfont{WenQuanYiZenHei-subset.ttf}[Path=../../../fonts/]
\newcommand\pyactive[1]{\catcode#1=13 \begingroup\lccode`\~=#1 \lowercase{\endgroup\def~}{{\pyfont\char#1}}}
\pyactive{"01CD}\pyactive{"01CE}\pyactive{"01CF}\pyactive{"01D0}\pyactive{"01D1}\pyactive{"01D2}\pyactive{"01D3}\pyactive{"01D4}\pyactive{"01D5}\pyactive{"01D6}\pyactive{"01D7}\pyactive{"01D8}\pyactive{"01D9}\pyactive{"01DA}\pyactive{"01DB}\pyactive{"01DC}
% Maths en sans-serif (Fira Math) avec chiffres et lettres droites de Plus
% Jakarta, pour que les formules se fondent dans le texte.
\usepackage[math-style=ISO,bold-style=ISO]{unicode-math}
\setmathfont{FiraMath-Regular.otf}[Path=../../../fonts/]
\setmathfont{PlusJakartaSans-Medium.ttf}[Path=../../../fonts/,range={up/num,up/{Latin,latin},\mathup}]
\usepackage{icomma} % virgule décimale sans espace en mode math
% Caractères mathématiques Unicode absents des polices de texte (Plus Jakarta,
% Archivo Black) : rendus par leur équivalent mathématique, en texte comme en
% formule — sinon ils s'affichent en carré vide (ex. « Arithmétique dans ℤ »).
\usepackage{newunicodechar}
\newunicodechar{ℝ}{\ensuremath{\mathbb{R}}}
\newunicodechar{ℤ}{\ensuremath{\mathbb{Z}}}
\newunicodechar{ℂ}{\ensuremath{\mathbb{C}}}
\newunicodechar{ℕ}{\ensuremath{\mathbb{N}}}
\newunicodechar{ℚ}{\ensuremath{\mathbb{Q}}}
\newunicodechar{𝔻}{\ensuremath{\mathbb{D}}}
\newunicodechar{α}{\ensuremath{\alpha}}
\newunicodechar{β}{\ensuremath{\beta}}
\newunicodechar{γ}{\ensuremath{\gamma}}
\newunicodechar{δ}{\ensuremath{\delta}}
\newunicodechar{ε}{\ensuremath{\varepsilon}}
\newunicodechar{θ}{\ensuremath{\theta}}
\newunicodechar{λ}{\ensuremath{\lambda}}
\newunicodechar{μ}{\ensuremath{\mu}}
\newunicodechar{σ}{\ensuremath{\sigma}}
\newunicodechar{τ}{\ensuremath{\tau}}
\newunicodechar{φ}{\ensuremath{\varphi}}
\newunicodechar{ψ}{\ensuremath{\psi}}
\newunicodechar{ω}{\ensuremath{\omega}}
\newunicodechar{Σ}{\ensuremath{\Sigma}}
% majuscules produites par \MakeUppercase dans les titres (α-aminés -> Α)
\newunicodechar{Α}{\ensuremath{\alpha}}
\newunicodechar{Β}{\ensuremath{\beta}}
\newunicodechar{Γ}{\ensuremath{\Gamma}}
\newunicodechar{Δ}{\ensuremath{\Delta}}
\newunicodechar{Ε}{\ensuremath{\varepsilon}}
\newunicodechar{Θ}{\ensuremath{\Theta}}
\newunicodechar{Λ}{\ensuremath{\Lambda}}
\newunicodechar{Μ}{\ensuremath{\mu}}
\newunicodechar{Τ}{\ensuremath{\tau}}
\newunicodechar{Φ}{\ensuremath{\Phi}}
\newunicodechar{Ψ}{\ensuremath{\Psi}}
\newunicodechar{Ω}{\ensuremath{\Omega}}
\newunicodechar{↦}{\ensuremath{\mapsto}}
\newunicodechar{∈}{\ensuremath{\in}}
\newunicodechar{∉}{\ensuremath{\notin}}
\newunicodechar{∧}{\ensuremath{\wedge}}
\newunicodechar{⇔}{\ensuremath{\Leftrightarrow}}
\newunicodechar{⟺}{\ensuremath{\Longleftrightarrow}}
\newunicodechar{⇒}{\ensuremath{\Rightarrow}}
\newunicodechar{⟹}{\ensuremath{\Longrightarrow}}
\newunicodechar{∘}{\ensuremath{\circ}}
\newunicodechar{∪}{\ensuremath{\cup}}
\newunicodechar{∩}{\ensuremath{\cap}}
\newunicodechar{≡}{\ensuremath{\equiv}}
\newunicodechar{⊂}{\ensuremath{\subset}}
\newunicodechar{⊥}{\ensuremath{\perp}}
\newunicodechar{∃}{\ensuremath{\exists}}
\newunicodechar{∀}{\ensuremath{\forall}}
\newunicodechar{∛}{\ensuremath{\sqrt[3]{\,}}}
\newunicodechar{∜}{\ensuremath{\sqrt[4]{\,}}}
\newunicodechar{⃗}{\ensuremath{{}^{\to}}}
\newunicodechar{ⁿ}{\ifmmode^{n}\else\textsuperscript{\ensuremath{n}}\fi}
\newunicodechar{ˣ}{\ifmmode^{x}\else\textsuperscript{\ensuremath{x}}\fi}
\newunicodechar{ᵇ}{\ifmmode^{b}\else\textsuperscript{\ensuremath{b}}\fi}
\newunicodechar{ʳ}{\ifmmode^{r}\else\textsuperscript{\ensuremath{r}}\fi}
\newunicodechar{ᵘ}{\ifmmode^{u}\else\textsuperscript{\ensuremath{u}}\fi}
\newunicodechar{ʸ}{\ifmmode^{y}\else\textsuperscript{\ensuremath{y}}\fi}
\newunicodechar{ᵃ}{\ifmmode^{a}\else\textsuperscript{\ensuremath{a}}\fi}
\newunicodechar{ᵉ}{\ifmmode^{e}\else\textsuperscript{\ensuremath{e}}\fi}
\newunicodechar{ᵏ}{\ifmmode^{k}\else\textsuperscript{\ensuremath{k}}\fi}
\newunicodechar{ᵖ}{\ifmmode^{p}\else\textsuperscript{\ensuremath{p}}\fi}
\newunicodechar{ᴺ}{\ifmmode^{N}\else\textsuperscript{\ensuremath{N}}\fi}
\newunicodechar{ᶜ}{\ifmmode^{c}\else\textsuperscript{\ensuremath{c}}\fi}
\newunicodechar{ʰ}{\ifmmode^{h}\else\textsuperscript{\ensuremath{h}}\fi}
\newunicodechar{ᵠ}{\ifmmode^{q}\else\textsuperscript{\ensuremath{q}}\fi}
\newunicodechar{ₙ}{\ifmmode_{n}\else\textsubscript{\ensuremath{n}}\fi}
\newunicodechar{ₐ}{\ifmmode_{a}\else\textsubscript{\ensuremath{a}}\fi}
\newunicodechar{ᵢ}{\ifmmode_{i}\else\textsubscript{\ensuremath{i}}\fi}
\newunicodechar{ᵣ}{\ifmmode_{r}\else\textsubscript{\ensuremath{r}}\fi}
\newunicodechar{ₖ}{\ifmmode_{k}\else\textsubscript{\ensuremath{k}}\fi}
\newunicodechar{ᵦ}{\ifmmode_{\beta}\else\textsubscript{\ensuremath{\beta}}\fi}
\newunicodechar{ₓ}{\ifmmode_{x}\else\textsubscript{\ensuremath{x}}\fi}
\newfontfamily\popdisplay{ArchivoBlack-Regular.ttf}[Path=../../../fonts/]
\newfontfamily\poplogo{SpaceGrotesk-Bold.ttf}[Path=../../../fonts/]
\newfontfamily\popsemi{PlusJakartaSans-SemiBold.ttf}[Path=../../../fonts/]
\newfontfamily\popxbold{PlusJakartaSans-ExtraBold.ttf}[Path=../../../fonts/]
\newfontfamily\popxboldls{PlusJakartaSans-ExtraBold.ttf}[Path=../../../fonts/,LetterSpace=3.0]

\setlist[enumerate]{leftmargin=1.7em,itemsep=5pt,topsep=4pt}
\setlist[itemize]{leftmargin=1.7em,itemsep=4pt,topsep=4pt,label={\color{poporange}\textbullet}}
\setlength{\parindent}{0pt}
\setlength{\parskip}{5pt}
\renewcommand{\arraystretch}{1.25}

% pgfplots : style Pop par défaut
\pgfplotsset{
  every axis/.append style={axis line style={popink,line width=1pt},tick style={popink},
    grid style={popline,line width=0.5pt},label style={font=\small},tick label style={font=\footnotesize}},
  popcurve/.style={poporange,line width=1.8pt,no markers,smooth},
}
% Même style disponible dans un \draw TikZ simple (hors pgfplots)
\tikzset{popcurve/.style={poporange,line width=1.8pt}}
\tikzset{popgrid/.style={popline,very thin}}
\colorlet{popcurve}{poporange} % le nom du style est parfois utilisé comme couleur

% ============================================================
% Pill / squiggle
% ============================================================
\newcommand{\poppill}[3]{%
  \tikz[baseline=-0.55ex]\node[draw=popink,line width=1.2pt,rounded corners=2.6pt,%
    fill=#2,inner xsep=6.5pt,inner ysep=3pt]{%
    \popxboldls\fontsize{7}{7}\selectfont\color{#3}\MakeUppercase{#1}};%
}
\newcommand{\popsquiggle}[1]{%
  \tikz[baseline=-0.4ex]{
    \draw[#1,line width=2.6pt,line cap=round]
      plot [smooth] coordinates {(0,0.09) (0.34,0.01) (0.68,0.21) (1.02,0.01) (1.36,0.10)};
  }%
}
% Mot-clé surligné (marqueur orange sous le texte)
\newcommand{\cle}[1]{{\popxbold\color{popdark}#1}}

% ============================================================
% En-tête / pied
% ============================================================
\pagestyle{fancy}
\fancyhf{}
\renewcommand{\headrulewidth}{0pt}
\renewcommand{\footrulewidth}{0pt}
\lhead{\raisebox{-0.15ex}{\poplogo\fontsize{13}{13}\selectfont\color{popink}OnBuch\color{poporange}+}}
\rhead{\poppill{\DOCMATIERE\ \textbullet\ \DOCNIVEAU\ \textbullet\ Leçon \DOCLECON}{white}{popink}}
\lfoot{\popxbold\fontsize{7.6}{7.6}\selectfont\color{popmuted}\MakeUppercase{Cours signé OnBuch+}\ \textbullet\ {\popsemi\DOCTITRE}}
\rfoot{\tikz[baseline=-0.6ex]\node[rounded corners=8.5pt,fill=popink,minimum height=17pt,minimum width=17pt,inner xsep=5pt,inner sep=0]{\popxbold\fontsize{8}{8}\selectfont\color{popcream}\thepage\,/\,\pageref*{LastPage}};}

% ============================================================
% Titres
% ============================================================
\newcommand{\chaptitre}[2]{%
  \begin{tcolorbox}[enhanced,colback=poporange,colframe=popink,boxrule=2.6pt,arc=11pt,
      drop shadow=popink,left=15pt,right=15pt,top=12pt,bottom=11pt,before skip=3pt,after skip=18pt]
    \poppill{#2}{popink}{popcream}\par\vspace{9pt}
    {\raggedright\hyphenpenalty=10000\exhyphenpenalty=10000\popdisplay\fontsize{22}{24}\selectfont\color{popcream}\MakeUppercase{#1}\par}
  \end{tcolorbox}
}
% Section numérotée : pastille orange avec le numéro + titre Archivo
\newcounter{popsec}
\newcommand{\coursec}[1]{%
  \stepcounter{popsec}\setcounter{popsub}{0}%
  \par\vspace{16pt}\noindent
  \begin{minipage}{\linewidth}
  \tikz[baseline=-0.7ex]\node[draw=popink,line width=1.6pt,fill=poporange,rounded corners=4pt,minimum size=22pt,inner sep=0]{\popdisplay\fontsize{12}{12}\selectfont\color{popcream}\thepopsec};%
  \hspace{8pt}{\popdisplay\fontsize{15}{17}\selectfont\color{popink}\MakeUppercase{#1}}\par
  \vspace{4pt}\noindent{\color{poporange}\rule{36pt}{3.6pt}}
  \end{minipage}\par\nopagebreak\vspace{8pt}
}
\newcounter{popsub}
\newcommand{\courssub}[1]{%
  \stepcounter{popsub}%
  \par\vspace{9pt}\noindent{\popxbold\fontsize{11.5}{13}\selectfont\color{popink}{\color{poporange}\thepopsec.\thepopsub}\hspace{6pt}#1}\par\nopagebreak\vspace{4pt}
}

% ============================================================
% Boîtes pédagogiques — même recette : fond blanc, bordure épaisse colorée,
% ombre dure encre, badge plein. Titre optionnel [#1].
% ============================================================
\tcbset{popbase/.style={breakable,enhanced,colback=white,boxrule=2.3pt,arc=9pt,
  shadow={3pt}{-3pt}{0pt}{popink},left=14pt,right=14pt,top=11pt,bottom=11pt,before skip=11pt,after skip=15pt,
  fontupper=\popsemi\fontsize{10.6}{14.5}\selectfont\color{popink},
  fontlower=\popsemi\fontsize{10.6}{14.5}\selectfont\color{popink}}}
% Coche dessinée (le glyphe ✓ n'existe pas dans les polices du document)
\newcommand{\popcheck}[1]{\tikz[baseline=-0.5ex]\draw[#1,line width=1.6pt,line cap=round,line join=round](-0.13,0.0)--(-0.04,-0.09)--(0.13,0.1);}
\providecommand{\coche}{\popcheck{popgreen}}
\providecommand{\resultat}[1]{\par\smallskip\noindent\fcolorbox{popgreen}{popgreenL}{\parbox{\dimexpr\linewidth-2\fboxsep-2\fboxrule\relax}{\textbf{Résultat.} #1}}\par\smallskip}
\newcommand{\popboxhead}[3]{\poppill{#1}{#2}{white}\ifx\relax#3\relax\else\hspace{6pt}{\popxbold\fontsize{10.6}{12}\selectfont\color{popink}#3}\fi\par\vspace{7pt}}

\newtcolorbox{aretenir}[1][]{popbase,colframe=popgreen,before upper={\popboxhead{À retenir}{popgreen}{#1}}}
% Texte en chinois (document de lecture, dialogue, extrait) : cadre
% d'étude. Un titre optionnel (source, auteur) s'affiche dans l'en-tête.
\newtcolorbox{textebox}[1][]{popbase,colframe=popblue,before upper={\popboxhead{文本 · Texte}{popblue}{#1}},after upper={}}
% Marques ✓ / ✗ (forme correcte / fautive) souvent utilisées par les modèles
\providecommand{\cross}{\textcolor{poppink}{\ensuremath{\times}}}
\providecommand{\xmark}{\cross}\providecommand{\crossmark}{\cross}\providecommand{\cmark}{\ensuremath{\checkmark}}
% Ligne de réponse / justification à compléter (inventée par les modèles)
\providecommand{\justification}[1]{\par\noindent\textit{Justification~:} \dotfill\par}
% \textipa (package tipa non chargé) : la police ne couvre pas l'API -> simple passage
\providecommand{\textipa}[1]{#1}
% Soulignement (ulem non chargé) : \uline, \ul
\providecommand{\uline}[1]{\underline{#1}}\providecommand{\ul}[1]{\underline{#1}}
% \glossaire{\item ...} inventée par les modèles : liste de glossaire
\providecommand{\glossaire}[1]{\begin{itemize}#1\end{itemize}}
% \alert (beamer) inventée par les modèles : texte mis en évidence
\providecommand{\alert}[1]{\textbf{\textcolor{poppink}{#1}}}
% variantes ASCII d'ordinaux écrites en macro
\providecommand{\iere}{\textsuperscript{re}}\providecommand{\ieme}{\textsuperscript{e}}\providecommand{\ere}{\textsuperscript{re}}\providecommand{\eme}{\textsuperscript{e}}\providecommand{\er}{\textsuperscript{er}}\providecommand{\ie}{\textsuperscript{e}}
% Ordinaux français écrits en macro par les modèles : 1\ière, 3\ième
\newcommand{\ière}{\textsuperscript{re}}\newcommand{\ième}{\textsuperscript{e}}
% \exemple (inventée par les modèles) : étiquette « Exemple : »
\providecommand{\exemple}{\textbf{Exemple~:}\ }
% \square utilisable aussi hors mode math (cases à cocher)
\AtBeginDocument{\let\squareMath\square\renewcommand{\square}{\ensuremath{\squareMath}}}
% Mot ou phrase en chinois à l'intérieur d'un paragraphe français
\newcommand{\zh}[1]{\ifmmode\text{#1}\else{#1}\fi}
\let\eng\zh \let\esp\zh \let\all\zh % alias tolérés
% flèches d'intonation inventées par les modèles
\providecommand{\textsearrow}{}\renewcommand{\textsearrow}{\ensuremath{\searrow}}\providecommand{\textnearrow}{}\renewcommand{\textnearrow}{\ensuremath{\nearrow}}
% \fr{..} : mot français (inventé par les modèles)
\providecommand{\fr}[1]{\textit{#1}}
% zhbox : encadré de texte chinois inventé par les modèles
\newenvironment{zhbox}{\begin{quote}}{\end{quote}}
% \courssubsub : titre de niveau 3 inventé par les modèles
\providecommand{\courssubsub}[1]{\par\medskip\noindent\textbf{#1}\par\smallskip}
% \zhbf : texte chinois en gras inventé par les modèles
\providecommand{\zhbf}[1]{\textbf{#1}}
% \vs : « versus » inventé par les modèles
\providecommand{\vs}{\textit{vs}\ }
% \blank : case à compléter inventée par les modèles
\providecommand{\blank}[1][]{\underline{\hspace{1.6em}}}
\newtcolorbox{exemplebox}[1][]{popbase,colframe=poporange,before upper={\popboxhead{Exemple}{poporange}{#1}}}
\newtcolorbox{definition}[1][]{popbase,colframe=popblue,before upper={\popboxhead{Définition}{popblue}{#1}}}
\newtcolorbox{propriete}[1][]{popbase,colframe=poppurple,before upper={\popboxhead{Propriété}{poppurple}{#1}}}
\newtcolorbox{methode}[1][]{popbase,colframe=popgold,before upper={\popboxhead{Méthode}{popgold}{#1}}}
\newtcolorbox{attention}[1][]{popbase,colframe=poppink,before upper={\popboxhead{Attention}{poppink}{#1}}}
\newtcolorbox{experience}[1][]{popbase,colframe=popblue,colback=popblueL,before upper={\popboxhead{Expérience}{popblue}{#1}}}
% « Le savais-tu ? » : carte violette pleine (contexte, culture, Cameroun)
\newtcolorbox{savaistu}[1][]{popbase,colback=poppurple,colframe=popink,
  fontupper=\popsemi\fontsize{10.4}{14.2}\selectfont\color{white},
  before upper={\poppill{Le savais-tu\,?}{white}{poppurple}\ifx\relax#1\relax\else\hspace{6pt}{\popxbold\color{white}#1}\fi\par\vspace{7pt}}}
% Exercice résolu : énoncé en haut, \tcblower puis la solution sur fond crème
\newcounter{popexo}\newcounter{popexr}
\newtcolorbox{exoresolu}[1][]{popbase,colframe=poporange,colbacklower=popcream,
  segmentation style={popink,line width=1.2pt,dashed},
  before upper={\stepcounter{popexr}\popboxhead{Exercice résolu \thepopexr}{poporange}{#1}},
  before lower={\poppill{Solution}{popink}{popcream}\par\vspace{6pt}}}
% Exercice (non résolu en place) — la correction est dans la partie Corrigés
\newtcolorbox{exercice}[1][]{popbase,colframe=popink,
  before upper={\stepcounter{popexo}\popboxhead{Exercice \thepopexo}{popink}{#1}}}
% Corrigé
\newtcolorbox{corrige}[1][]{popbase,colframe=popgreen,colback=popgreenL,
  before upper={\popboxhead{Corrigé}{popgreen}{#1}}}
% Objectifs / prérequis / bilan
\newtcolorbox{objectifs}{popbase,colframe=popink,colback=poporangeL,
  before upper={\popboxhead{Objectifs de la leçon}{popink}{}}}
\newtcolorbox{prerequis}{popbase,colframe=popink,colback=popblueL,
  before upper={\popboxhead{Prérequis}{popblue}{}}}
\newtcolorbox{bilan}{popbase,colframe=popink,colback=white,boxrule=2.8pt,
  before upper={\popboxhead{Fiche bilan}{poporange}{L'essentiel à connaître pour l'examen}}}
% Figure encadrée avec légende
\newtcolorbox{popfigure}[1][]{enhanced,breakable=false,colback=white,colframe=popline,boxrule=1.4pt,arc=8pt,
  left=10pt,right=10pt,top=10pt,bottom=8pt,before skip=10pt,after skip=12pt,halign=center,
  lower separated=false,halign lower=center,
  fontlower=\popsemi\fontsize{9}{11.5}\selectfont\color{popmuted},
  #1}
\newcommand{\legende}[1]{\par\vspace{6pt}{\popsemi\fontsize{9}{11.5}\selectfont\color{popmuted}#1}}

% ============================================================
% Signature OnBuch+
% ============================================================
\newcommand{\poplogoinline}[1]{{\poplogo\fontsize{#1}{#1}\selectfont\color{popink}OnBuch\color{poporange}+}}
\newcommand{\signatureonbuch}{%
  \begin{tcolorbox}[enhanced,colback=popink,colframe=popink,boxrule=2.2pt,arc=10pt,
      drop shadow=poporange,left=14pt,right=14pt,top=10pt,bottom=10pt,before skip=0pt,after skip=0pt]
    \tikz[baseline=-0.7ex]\node[circle,fill=popgreen,draw=popcream,line width=1.3pt,minimum size=22pt,inner sep=0]{\popcheck{white}};%
    \hspace{9pt}\begin{minipage}[c]{0.62\linewidth}
      {\popxbold\fontsize{12}{13}\selectfont\color{popcream}Signé\ \textbullet\ {\poplogo OnBuch\color{poporange}+}}\par\vspace{2pt}
      {\raggedright\popsemi\fontsize{8}{10.5}\selectfont\color{popline}\MakeUppercase{Équipe pédagogique OnBuch+}\\\MakeUppercase{Conforme au programme officiel MINESEC}\par}
    \end{minipage}\hfill
    {\poplogo\fontsize{22}{22}\selectfont\color{popcream}OnBuch\color{poporange}+}
  \end{tcolorbox}%
}

% ============================================================
% Page de garde
% #1 titre · #2 matière · #3 niveau · #4 module · #5 n° leçon · #6 durée ·
% #7 description / objectifs (texte ou itemize)
% ============================================================
\newcommand{\pagegarde}[7]{%
  \thispagestyle{empty}
  \newgeometry{a4paper,margin=1.7cm,top=1.6cm,bottom=1.4cm}%
  \begin{tikzpicture}[remember picture,overlay]
    \fill[poporange!18] ([xshift=-3.2cm,yshift=-3.6cm]current page.north east) circle (4.6cm);
    \fill[poppurple!14] ([xshift=2.4cm,yshift=7.6cm]current page.south west) circle (3.8cm);
  \end{tikzpicture}%
  {\poplogo\fontsize{30}{30}\selectfont\color{popink}OnBuch\color{poporange}+}\hfill
  \raisebox{0.6ex}{\poppill{Cours signé \textbullet\ Premium}{popink}{popcream}}\par
  \vspace{1pt}\popsquiggle{poppurple}\par
  \vspace{8pt}
  \begin{tcolorbox}[enhanced,colback=poporange,colframe=popink,boxrule=2.8pt,arc=13pt,
      drop shadow=popink,left=18pt,right=18pt,top=12pt,bottom=13pt,after skip=10pt]
    \poppill{#2 \textbullet\ #3}{popink}{popcream}\hspace{5pt}\poppill{#4}{white}{popink}\par\vspace{8pt}
    {\popdisplay\fontsize{14}{14}\selectfont\color{popink}LEÇON #5}\par\vspace{4pt}
    {\raggedright\hyphenpenalty=10000\exhyphenpenalty=10000\popdisplay\fontsize{25}{27}\selectfont\color{popcream}\MakeUppercase{#1}\par}
    \vspace{8pt}
    {\popxbold\fontsize{9.5}{9.5}\selectfont\color{popink}\MakeUppercase{Durée conseillée : #6}}
  \end{tcolorbox}
  \begin{tcolorbox}[enhanced,colback=white,colframe=popink,boxrule=2.2pt,arc=10pt,
      drop shadow=popink,left=16pt,right=16pt,top=10pt,bottom=10pt,after skip=8pt]
    \poppill{Dans cette leçon}{poppurple}{white}\par\vspace{5pt}
    \popsemi\fontsize{9.3}{12.6}\selectfont\color{popink}#7
  \end{tcolorbox}
  \noindent\begin{minipage}[t]{0.487\linewidth}
    \begin{tcolorbox}[enhanced,colback=popgreen,colframe=popink,boxrule=2.2pt,arc=10pt,
        drop shadow=popink,left=11pt,right=11pt,top=10pt,bottom=10pt,height=3.1cm,valign=center]
      \tcbox[colback=white,colframe=popink,boxrule=1.3pt,arc=5pt,boxsep=0pt,nobeforeafter,
        left=3pt,right=3pt,top=3pt,bottom=3pt,valign=center]{\qrcode[height=1.9cm]{https://play.google.com/store/apps/details?id=com.luvvix.onbuch&hl=fr}}%
      \hspace{8pt}\begin{minipage}[c]{0.5\linewidth}
        {\popdisplay\fontsize{11}{12.5}\selectfont\color{white}\MakeUppercase{Télécharge OnBuch}}\par\vspace{4pt}
        {\raggedright\popsemi\fontsize{8.4}{11}\selectfont\color{white}Gratuit \textbullet\ Google Play\\Tuteur IA, annales, concours, résultats\par}
      \end{minipage}
    \end{tcolorbox}
  \end{minipage}\hfill
  \begin{minipage}[t]{0.487\linewidth}
    \begin{tcolorbox}[enhanced,colback=poppurple,colframe=popink,boxrule=2.2pt,arc=10pt,
        drop shadow=popink,left=11pt,right=11pt,top=10pt,bottom=10pt,height=3.1cm,valign=center]
      \tikz[baseline=-0.7ex]\node[circle,fill=white,draw=popink,line width=1.3pt,minimum size=1.6cm,inner sep=0]{\popdisplay\fontsize{24}{24}\selectfont\color{poppurple}+};%
      \hspace{8pt}\begin{minipage}[c]{0.6\linewidth}
        {\popdisplay\fontsize{11}{12.5}\selectfont\color{white}\MakeUppercase{Passe à OnBuch+}}\par\vspace{4pt}
        {\raggedright\popsemi\fontsize{8.4}{11}\selectfont\color{white}Tous les cours signés, Tuteur IA illimité, fiches PDF — onglet OnBuch+ de l'app\par}
      \end{minipage}
    \end{tcolorbox}
  \end{minipage}
  \vspace{8pt}
  \popcontenu
  \vfill
  \signatureonbuch
  \restoregeometry
  \newpage
}

% Rangée « Ce cours contient » (page de garde)
\newcommand{\poptile}[3]{% #1 couleur #2 gros texte #3 légende
  \begin{tcolorbox}[enhanced,colback=#1,colframe=popink,boxrule=2pt,arc=9pt,shadow={2.5pt}{-2.5pt}{0pt}{popink},
      left=6pt,right=6pt,top=9pt,bottom=9pt,halign=center,valign=center,height=1.75cm,width=\linewidth,nobeforeafter]
    {\popdisplay\fontsize{12.5}{13}\selectfont\color{white}\MakeUppercase{#2}}\par\vspace{3pt}
    {\centering\popsemi\fontsize{7.8}{9.5}\selectfont\color{white}#3\par}
  \end{tcolorbox}}
\newcommand{\popcontenu}{%
  \par\noindent{\popxboldls\fontsize{7.5}{7.5}\selectfont\color{popmuted}CE COURS SIGNÉ CONTIENT}\par\vspace{7pt}
  \noindent\begin{minipage}[t]{0.235\linewidth}\poptile{popblue}{Cours}{complet, illustré, pas à pas}\end{minipage}\hfill
  \begin{minipage}[t]{0.235\linewidth}\poptile{poporange}{Méthodes}{et exercices résolus}\end{minipage}\hfill
  \begin{minipage}[t]{0.235\linewidth}\poptile{poppink}{Exercices}{type examen + corrigés}\end{minipage}\hfill
  \begin{minipage}[t]{0.235\linewidth}\poptile{popgreen}{Fiche bilan}{l'essentiel à retenir}\end{minipage}\par}

% Sommaire
\newtcolorbox{sommaire}{enhanced,colback=white,colframe=popink,boxrule=2.2pt,arc=10pt,
  drop shadow=popink,left=18pt,right=18pt,top=13pt,bottom=13pt,before skip=6pt,after skip=20pt,
  before upper={\poppill{Sommaire}{poppurple}{white}\par\vspace{9pt}\popsemi\fontsize{10.6}{14.5}\selectfont\color{popink}}}

% Signature de fin de cours — poussée tout en bas de la dernière page
\newcommand{\findecours}{%
  \par\vfill\nopagebreak
  \begin{center}\popsquiggle{poporange}\end{center}\vspace{4pt}
  \signatureonbuch
}
\AtBeginDocument{\def\llbracket{\mathopen{[\mkern-3.5mu[}}\def\rrbracket{\mathclose{]\mkern-3.5mu]}}} % intervalles d'entiers (glyphe ⟦ absent de la police)
% Glyphes absents de Fira Math : redéfinis à partir de glyphes disponibles
\AtBeginDocument{%
  \def\setminus{\mathbin{\backslash}}\let\smallsetminus\setminus
  \def\vdots{\mathord{\vbox{\baselineskip3.2pt\lineskiplimit0pt\kern4pt\hbox{.}\hbox{.}\hbox{.}}}}%
  \def\ddots{\mathinner{\mkern1mu\raise6.5pt\hbox{.}\mkern2mu\raise3.6pt\hbox{.}\mkern2mu\raise0.7pt\hbox{.}\mkern1mu}}%
  \def\triangle{\mathord{\scalebox{0.9}{$\symup{\Delta}$}}}%
  \def\nabla{\mathord{\raisebox{\depth}{\scalebox{1}[-1]{$\symup{\Delta}$}}}}%
  \def\checkmark{\popcheck{popdark}}%
  \def\star{\mathbin{\ast}}\def\bigstar{\mathord{\ast}}%
  \def\diamond{\mathbin{\scalebox{0.6}{$\Diamond$}}}%
  \def\bigcup{\mathop{\mathchoice{\raisebox{-0.3em}{\scalebox{1.7}{$\cup$}}}{\scalebox{1.25}{$\cup$}}{\cup}{\cup}}\displaylimits}%
  \def\bigcap{\mathop{\mathchoice{\raisebox{-0.3em}{\scalebox{1.7}{$\cap$}}}{\scalebox{1.25}{$\cap$}}{\cap}{\cap}}\displaylimits}%
  \def\lesseqqgtr{\mathrel{\lessgtr}}\def\bigoplus{\mathop{\oplus}\displaylimits}%
}
\providecommand{\ssi}{si et seulement si} % abréviation fréquente
\AtBeginDocument{\providecommand{\wideparen}{\overparen}} % arc de cercle

%% FIGFIX-BEGIN v4 — correctifs globaux des figures (docs/onbuch-generation/tools/figfix)
\usepackage{adjustbox}\usepackage{etoolbox}
% toute figure TikZ/pgfplots plus large que la ligne est réduite à la largeur disponible
\BeforeBeginEnvironment{tikzpicture}{\begin{adjustbox}{max width=\linewidth}}
\AfterEndEnvironment{tikzpicture}{\end{adjustbox}}
% légendes pgfplots lisibles par-dessus les courbes
\pgfplotsset{every axis/.append style={legend style={fill=white,fill opacity=0.92,text opacity=1,draw=black!15,inner sep=3pt}}}
% étiquettes d'arêtes (syntaxe edge["..."]) avec fond blanc
\tikzset{every edge quotes/.append style={fill=white,inner sep=1.2pt}}
%% FIGFIX-END
\begin{document}

\pagegarde{Leçon 21 — Vocabulaire et compréhension de texte : démocratie}{Chinois}{Tle A}{Module 3 : Citoyenneté et ouverture au monde}{ZH21}{2 h}{Cette leçon fait découvrir le vocabulaire chinois du thème de la démocratie (élections, vote, droits et devoirs du citoyen) à travers un texte support original. Elle entraîne la lecture compréhensive, la prononciation des tons et l'emploi du lexique dans des phrases authentiques.
\vspace{4pt}
\begin{multicols}{2}\small
\begin{itemize}
\item À la fin de cette leçon, je sais lire à voix haute un court texte chinois sur la démocratie en respectant les tons.
\item À la fin de cette leçon, je sais reconnaître et traduire au moins 15 mots du thème (民主, 选举, 投票, 公民, 权利, 义务...).
\item À la fin de cette leçon, je sais identifier les radicaux et tracer correctement les caractères clés du thème (民, 主, 选, 权).
\item À la fin de cette leçon, je sais répondre par écrit à des questions de compréhension sur le texte.
\item À la fin de cette leçon, je sais employer les collocations du thème (参加选举, 行使权利, 履行义务) dans des phrases complètes.
\item À la fin de cette leçon, je sais comparer brièvement le vocabulaire de la vie civique au Cameroun et en Chine.
\end{itemize}\end{multicols}}

\begin{sommaire}
\begin{multicols}{2}
\begin{enumerate}
\item Mise en route et découverte du thème
\item Texte support : lecture guidée
\item Glossaire du thème
\item Clés et ordre des traits
\item Compréhension du texte
\item Collocations et production guidée
\item Activité d'intégration
\item Méthodes et exercices résolus
\item Exercices
\item Corrigés des exercices
\item Fiche bilan
\end{enumerate}
\end{multicols}
\end{sommaire}

\begin{objectifs}
\begin{itemize}
\item À la fin de cette leçon, je sais lire à voix haute un court texte chinois sur la démocratie en respectant les tons.
\item À la fin de cette leçon, je sais reconnaître et traduire au moins 15 mots du thème (民主, 选举, 投票, 公民, 权利, 义务...).
\item À la fin de cette leçon, je sais identifier les radicaux et tracer correctement les caractères clés du thème (民, 主, 选, 权).
\item À la fin de cette leçon, je sais répondre par écrit à des questions de compréhension sur le texte.
\item À la fin de cette leçon, je sais employer les collocations du thème (参加选举, 行使权利, 履行义务) dans des phrases complètes.
\item À la fin de cette leçon, je sais comparer brièvement le vocabulaire de la vie civique au Cameroun et en Chine.
\item À la fin de cette leçon, je sais rédiger 3 à 4 phrases guidées pour exprimer une opinion simple sur la démocratie.
\end{itemize}
\end{objectifs}

\begin{prerequis}
\begin{itemize}
\item Vocabulaire de base de la vie scolaire et sociale (学校, 学生, 国家, 人们) vu en Première
\item Structures de phrase simples : sujet + verbe + complément ; 是 / 有
\item Particules de base : 的, 了, 吗 ; négation 不 / 没
\item Notions sur les radicaux et l'ordre des traits (leçons d'écriture de Première)
\item Verbes modaux 应该, 可以, 必须 (début de Terminale)
\end{itemize}
\end{prerequis}

\newpage

\coursec{Mise en route et découverte du thème}

\courssub{Situation d'accroche et problématique}

Imagine la scène. Nous sommes en février. Dans ta classe de Terminale, le censeur annonce l'élection du délégué : deux candidats se présentent, chacun fait un petit discours, puis tous les élèves écrivent un nom sur un papier, plient le bulletin et le glissent dans une boîte. On dépouille les votes au tableau, devant tout le monde. Le candidat majoritaire est proclamé, on applaudit. Tu viens de vivre, à ton échelle, un moment de \cle{démocratie}.

Ce mot, tu le rencontres partout : dans les journaux de Yaoundé et de Douala, dans les discours officiels, dans les débats entre amis au sujet de la vie de la cité. Mais si un journaliste chinois venait interviewer ta classe sur la vie démocratique de ton lycée, saurais-tu le comprendre ? Comment dit-on \emph{voter}, \emph{élection}, \emph{citoyen}, \emph{droits}, \emph{devoirs} en chinois ?

\textbf{Problématique de la leçon :} comment lire et comprendre un texte chinois simple sur la démocratie, et comment employer le vocabulaire essentiel de la vie citoyenne ?

\begin{experience}[Brainstorming d'ouverture (5 min)]
En français, réponds à la question : \emph{que signifie pour toi le mot « démocratie » ?} Écris au tableau tous les mots qui te viennent à l'esprit. Mots attendus : \textbf{vote, élection, citoyen, droits, devoirs, liberté}. Garde cette liste : nous allons la « traduire » en chinois au cours de la leçon.
\end{experience}

\courssub{Le mot-clé de la leçon : 民主 (mínzhǔ)}

Observons d'abord le mot central avant toute explication.

\begin{textebox}[Le mot-clé du jour]
民主 \quad \emph{mínzhǔ} \quad « la démocratie »\\
民 \quad \emph{mín} \quad « le peuple »\\
主 \quad \emph{zhǔ} \quad « maître, diriger »\\
民主 = littéralement « le peuple (民) est maître (主) ».
\end{textebox}

\begin{definition}[民主 mínzhǔ]
\zh{民主} est un nom composé de deux caractères : \zh{民} \emph{mín} « le peuple » (que l'on retrouve dans \zh{人民} \emph{rénmín} « le peuple ») et \zh{主} \emph{zhǔ} « maître, principal, diriger » (que l'on retrouve dans \zh{主人} \emph{zhǔrén} « le maître, l'hôte »). Le mot signifie donc exactement ce que dit le grec \emph{dêmos-kratos} : « le pouvoir du peuple ». Exemple : \zh{我们学校很民主。} \emph{Wǒmen xuéxiào hěn mínzhǔ.} « Notre école est très démocratique. »
\end{definition}

\begin{savaistu}[Un mot moderne pour une idée ancienne]
Le mot \zh{民主} est un emprunt moderne : il a été formé au XIX\textsuperscript{e} siècle pour traduire le mot occidental \emph{democracy}, en réutilisant deux caractères très anciens. Aujourd'hui, il fait partie des « douze mots des valeurs socialistes » affichés dans les écoles chinoises, aux côtés de mots comme \zh{自由} \emph{zìyóu} « liberté » et \zh{法治} \emph{fǎzhì} « État de droit ». Au Cameroun comme en Chine, la question de la participation des citoyens à la vie publique est donc un sujet de débat vivant : c'est exactement le sujet de notre texte.
\end{savaistu}

\courssub{Première écoute globale du texte support}

\begin{methode}[Comment aborder un texte chinois en quatre étapes]
\begin{enumerate}
\item \textbf{Écoute globale} : le professeur lit le texte à voix haute, tu écoutes \emph{sans texte sous les yeux}. Objectif : saisir le sujet général, les sons, les mots déjà connus.
\item \textbf{Lecture silencieuse} : tu lis seul en repérant les mots connus et en entourant les mots inconnus.
\item \textbf{Lecture guidée} : on relit ensemble avec le glossaire, phrase par phrase.
\item \textbf{Questions de compréhension} : tu vérifies ta compréhension globale puis fine.
\end{enumerate}
\end{methode}

Nous sommes à l'étape 1. Ferme les yeux ou regarde le tableau : le professeur lit l'article ci-dessous une première fois, à vitesse naturelle. Ne cherche pas à tout comprendre : écoute la mélodie des tons et essaie d'attraper les mots que tu connais déjà (\zh{学校}, \zh{学生}, \zh{我们}…).

\begin{textebox}[Texte support — lu par l'enseignant (première écoute)]
我们学校的生活很民主。\\
\emph{Wǒmen xuéxiào de shēnghuó hěn mínzhǔ.}\\
« La vie de notre école est très démocratique. »\\
每年二月，学生们选举班长。\\
\emph{Měi nián èryuè, xuéshengmen xuǎnjǔ bānzhǎng.}\\
« Chaque année, en février, les élèves élisent le délégué de classe. »\\
每个人都可以投票，也可以当候选人。\\
\emph{Měi ge rén dōu kěyǐ tóupiào, yě kěyǐ dāng hòuxuǎnrén.}\\
« Chacun peut voter, et peut aussi être candidat. »\\
投票以后，我们一起数票，票最多的人当班长。\\
\emph{Tóupiào yǐhòu, wǒmen yìqǐ shǔ piào, piào zuì duō de rén dāng bānzhǎng.}\\
« Après le vote, nous comptons les bulletins ensemble ; celui qui a le plus de voix devient délégué. »\\
公民有权利，也有义务。\\
\emph{Gōngmín yǒu quánlì, yě yǒu yìwù.}\\
« Le citoyen a des droits, mais aussi des devoirs. »\\
我们有说话的自由，但是我们也要听别人的意见。\\
\emph{Wǒmen yǒu shuōhuà de zìyóu, dànshì wǒmen yě yào tīng biérén de yìjiàn.}\\
« Nous avons la liberté de parler, mais nous devons aussi écouter l'avis des autres. »\\
民主不只是选举，也是每天尊重别人。\\
\emph{Mínzhǔ bù zhǐshì xuǎnjǔ, yěshì měitiān zūnzhòng biérén.}\\
« La démocratie, ce n'est pas seulement les élections, c'est aussi respecter les autres chaque jour. »
\end{textebox}

\courssub{Hypothèses de sens}

Après cette première écoute, formule des hypothèses. Ne crains pas l'erreur : en compréhension orale, une bonne hypothèse vaut mieux qu'un silence.

\begin{exemplebox}[Questions-guide pour les hypothèses]
\begin{itemize}
\item \textbf{De quoi parle le texte ?} (Hypothèse attendue : de la vie d'une école, d'une élection.)
\item \textbf{Qui parle ?} (Probablement un élève ou un journaliste qui décrit son école : on a entendu \zh{我们} « nous ».)
\item \textbf{Où et quand ?} (On a peut-être entendu \zh{二月} « février » et \zh{学校} « école ».)
\item \textbf{Quels mots nouveaux reviennent souvent ?} (\zh{民主}, \zh{选举}, \zh{投票}…)
\end{itemize}
\end{exemplebox}

\courssub{Carte mentale du vocabulaire du thème}

Voici, en avant-première, la constellation de mots que nous allons étudier dans cette leçon. Chaque branche correspond à un mot du brainstorming français du début : retrouve les correspondances.

\begin{popfigure}
\begin{tikzpicture}[mindmap, grow cyclic, every node/.style={concept}, concept color=popblueL, text=popdark, root concept/.append style={concept color=popblue, text=white, font=\Large\bfseries}, level 1/.append style={level distance=3.4cm, sibling angle=60}, level 1 concept/.append style={font=\small}]
\node[root concept, align=center]{民主\\ mínzhǔ\\ démocratie}
  child { node[align=center]{选举\\ xuǎnjǔ\\ élection} }
  child { node[align=center]{投票\\ tóupiào\\ voter} }
  child { node[align=center]{公民\\ gōngmín\\ citoyen} }
  child { node[align=center]{权利\\ quánlì\\ droits} }
  child { node[align=center]{义务\\ yìwù\\ devoirs} }
  child { node[align=center]{自由\\ zìyóu\\ liberté} };
\end{tikzpicture}
\legende{Carte mentale lexicale : autour de 民主 (mínzhǔ), les six mots-clés de la vie citoyenne.}
\end{popfigure}

\begin{center}
\begin{tabularx}{\linewidth}{|X|X|X|X|}
\hline
\rowcolor{popblueL} Caractères & Pinyin & Nature & Traduction \\
\hline
民主 & mínzhǔ & nom / adjectif & démocratie ; démocratique \\
\hline
选举 & xuǎnjǔ & verbe / nom & élire ; élection \\
\hline
投票 & tóupiào & verbe & voter (déposer un bulletin) \\
\hline
公民 & gōngmín & nom & citoyen \\
\hline
权利 & quánlì & nom & droit (ce qu'on a le droit de faire) \\
\hline
义务 & yìwù & nom & devoir, obligation \\
\hline
自由 & zìyóu & nom / adjectif & liberté ; libre \\
\hline
\end{tabularx}
\end{center}

\begin{attention}[Piège fréquent : 权利 ou 权力 ?]
Les francophones confondent souvent \zh{权利} \emph{quánlì} « le droit » (du citoyen : droit de vote, droit de parole) et \zh{权力} \emph{quánlì} « le pouvoir » (du gouvernement, du chef). Même prononciation, sens différents ! Retiens : \zh{利} évoque l'\emph{intérêt, l'avantage} → le droit qui te protège ; \zh{力} évoque la \emph{force} → le pouvoir qui commande. Autre piège : \zh{投票} se dit pour « voter » au sens de déposer un bulletin ; on ne dit jamais \zh{*}.
\end{attention}

\begin{exoresolu}[Première mise en bouche]
Énoncé : traduis en chinois la phrase « Les élèves votent en février. » en utilisant les mots de la carte mentale.
\tcblower
Solution détaillée : « les élèves » = \zh{学生们} \emph{xuéshengmen} ; « voter » = \zh{投票} \emph{tóupiào} ; « en février » = \zh{在二月} \emph{zài èryuè}. En chinois, le complément de temps se place \emph{avant} le verbe, contrairement au français. On obtient : \zh{学生们在二月投票。} \emph{Xuéshengmen zài èryuè tóupiào.} Erreur typique du francophone : placer \zh{在二月} après le verbe, comme en français — c'est incorrect.
\end{exoresolu}

\courssub{Objectifs et plan de travail de la leçon}

\begin{aretenir}[À l'issue de cette leçon, tu sauras…]
\begin{itemize}
\item lire et comprendre un court texte chinois sur la démocratie et la vie citoyenne ;
\item reconnaître, prononcer (pinyin et tons) et employer les mots \zh{民主}, \zh{选举}, \zh{投票}, \zh{公民}, \zh{权利}, \zh{义务}, \zh{自由} ;
\item répondre, en chinois et en français, à des questions de compréhension sur le texte.
\end{itemize}
\textbf{Plan de travail :} 2. lecture guidée du texte support ; 3. glossaire complet du thème ; 4. clés (radicaux) et ordre des traits des caractères nouveaux ; 5. questions de compréhension ; 6. collocations et production guidée.
\end{aretenir}

\coursec{Texte support : lecture guidée}

Dans la section précédente, nous avons découvert le thème de la démocratie à travers des images et des mots-clés : \zh{民主} (mínzhǔ, démocratie), \zh{选举} (xuǎnjǔ, élection), \zh{投票} (tóupiào, voter). Il est temps maintenant de rencontrer ces mots dans un texte vivant. Nous allons lire ensemble le reportage d'un journaliste chinois venu interviewer des élèves d'un lycée de Yaoundé, au lendemain de l'élection de leur délégué de classe. Ce texte est le cœur de la leçon : c'est dans lui que nous puiserons le glossaire (section 3), les clés des caractères (section 4) et les questions de compréhension (section 5).

\courssub{Le texte support : un reportage à Yaoundé}

Lisez d'abord silencieusement le texte ci-dessous. Ne cherchez pas à tout comprendre mot par mot : repérez les mots que vous connaissez déjà (\zh{我们}, \zh{学校}, \zh{老师}, \zh{同学}…) et devinez le sens général.

\begin{textebox}[Reportage : l'élection du délégué de classe à Yaoundé]
上个星期，我们学校举行了班长选举。\\
Shàng ge xīngqī, wǒmen xuéxiào jǔxíngle bānzhǎng xuǎnjǔ.\\
« La semaine dernière, notre école a organisé l'élection du délégué de classe. »\\[4pt]
每个同学都有一票，大家认真投票。\\
Měi ge tóngxué dōu yǒu yī piào, dàjiā rènzhēn tóupiào.\\
« Chaque élève avait une voix, et tous ont voté sérieusement. »\\[4pt]
老师说，选举是公民的权利，也是义务。\\
Lǎoshī shuō, xuǎnjǔ shì gōngmín de quánlì, yě shì yìwù.\\
« Le professeur a dit que voter est un droit du citoyen, mais aussi un devoir. »\\[4pt]
被选上的李明说：「我会为大家服务，听取每个人的意见。」\\
Bèi xuǎnshàng de Lǐ Míng shuō : « Wǒ huì wèi dàjiā fúwù, tīngqǔ měi ge rén de yìjiàn. »\\
« Li Ming, qui a été élu, a déclaré : Je vous servirai et j'écouterai l'avis de chacun. »\\[4pt]
同学们都觉得，民主就是人人参与、人人负责。\\
Tóngxuémen dōu juéde, mínzhǔ jiùshì rénrén cānyù, rénrén fùzé.\\
« Les élèves pensent tous que la démocratie, c'est la participation et la responsabilité de tous. »\\[4pt]
虽然我们还年轻，但是我们从小学习做负责任的公民。\\
Suīrán wǒmen hái niánqīng, dànshì wǒmen cóngxiǎo xuéxí zuò fùzérèn de gōngmín.\\
« Bien que nous soyons encore jeunes, nous apprenons dès l'enfance à devenir des citoyens responsables. »
\end{textebox}

\courssub{Lecture guidée en trois étapes}

\begin{methode}[Comment lire un texte chinois à voix haute]
\begin{enumerate}
\item \textbf{Lecture en écho.} Le professeur lit une phrase, la classe répète en chœur, en imitant les tons et le rythme. On procède phrase par phrase, puis paragraphe par paragraphe.
\item \textbf{Grouper par sens.} Ne lisez pas syllabe par syllabe ! Découpez la phrase en groupes de sens : \zh{我们学校} / \zh{举行了} / \zh{班长选举} (wǒmen xuéxiào / jǔxíngle / bānzhǎng xuǎnjǔ) — « notre école / a organisé / l'élection du délégué ». Entre chaque groupe, une micro-pause.
\item \textbf{Lecture individuelle.} Chaque élève lit ensuite une phrase à voix haute. Attention au \cle{sandhi du 3\textsuperscript{e} ton} : deux troisièmes tons qui se suivent se prononcent 2\textsuperscript{e} ton + 3\textsuperscript{e} ton. Ainsi \zh{选举} s'écrit xuǎnjǔ mais se prononce \emph{xuánjǔ} à l'oral ; de même \zh{可以} kěyǐ se prononce \emph{kéyǐ}.
\end{enumerate}
\end{methode}

\begin{attention}[Le sandhi du troisième ton]
En chinois parlé, quand deux syllabes au 3\textsuperscript{e} ton se suivent, la première passe au 2\textsuperscript{e} ton. L'écriture pinyin garde les tons d'origine, mais la bouche change ! Dans notre texte : \zh{选举} (xuǎnjǔ → \emph{xuánjǔ}), \zh{听取} (tīngqǔ, pas de changement car tīng est au 1\textsuperscript{er} ton). Erreur typique des francophones : prononcer deux tons descendants complets, ce qui coupe le rythme de la phrase.
\end{attention}

\courssub{Repérage dans le texte : souligner et encadrer}

Après la première lecture, prenez un crayon et travaillez directement sur le texte :

\begin{itemize}
\item \textbf{Soulignez} les mots du glossaire du thème (que nous étudierons en détail à la section 3) : \zh{选举} (élection), \zh{投票} (voter), \zh{票} (voix, bulletin), \zh{公民} (citoyen), \zh{权利} (droit), \zh{义务} (devoir), \zh{民主} (démocratie), \zh{参与} (participer), \zh{负责} (être responsable), \zh{意见} (avis, opinion).
\item \textbf{Encadrez} les mots déjà connus des leçons précédentes : \zh{上个星期} (la semaine dernière), \zh{学校} (école), \zh{同学} (camarade de classe), \zh{老师} (professeur), \zh{大家} (tout le monde), \zh{觉得} (penser, trouver), \zh{学习} (apprendre).
\item \textbf{Entourez d'un cercle} les mots-outils grammaticaux déjà étudiés et réemployés ici : \zh{被}, \zh{每}, \zh{都}, \zh{虽然}, \zh{但是}.
\end{itemize}

Ce repérage montre un principe fondamental de la lecture en chinois : un texte nouveau est toujours un mélange de mots connus (environ 60 \% ici) et de mots nouveaux que le contexte permet de deviner.

\courssub{Comprendre la structure du texte}

Un bon lecteur voit le « squelette » d'un texte. Notre reportage suit une progression logique en quatre temps, typique du récit journalistique : la situation, l'action, la parole, la leçon.

\begin{popfigure}
\begin{tikzpicture}[
  node distance=0.45cm,
  etape/.style={draw=popblue, fill=popblueL, rounded corners, thick, minimum width=2.9cm, minimum height=1.5cm, align=center, font=\small},
  fleche/.style={-{Stealth[length=3mm]}, thick, poporange}
]
\node[etape, align=center] (a){\textbf{1. Situation}\\举行班长选举\\jǔxíng bānzhǎng xuǎnjǔ\\organiser l'élection};
\node[etape, right=of a, fill=poporangeL, draw=poporange, align=center] (b){\textbf{2. Action}\\大家认真投票\\dàjiā rènzhēn tóupiào\\tous votent sérieusement};
\node[etape, right=of b, fill=popgreenL, draw=popgreen, align=center] (c){\textbf{3. Déclaration}\\李明说：为大家服务\\Lǐ Míng shuō : wèi dàjiā fúwù\\Li Ming : servir tous};
\node[etape, right=of c, fill=poppurpleL, draw=poppurple, align=center] (d){\textbf{4. Leçon}\\民主 = 参与 + 负责\\mínzhǔ = cānyù + fùzé\\démocratie = participer + être responsable};
\draw[fleche] (a) -- (b);
\draw[fleche] (b) -- (c);
\draw[fleche] (c) -- (d);
\end{tikzpicture}
\legende{Structure du texte : de la situation (l'élection) à la leçon de citoyenneté (la démocratie).}
\end{popfigure}

Ce schéma vous servira de plan pour raconter le texte de mémoire à l'oral : retenez les quatre verbes-charnières \zh{举行} (organiser) → \zh{投票} (voter) → \zh{说} (déclarer) → \zh{觉得} (penser).

\courssub{Points de langue réemployés dans le texte}

Le texte n'invente pas de nouvelle grammaire : il \emph{réemploie} trois structures déjà étudiées. Les reconnaître dans un texte authentique est la compétence clé de la compréhension de l'écrit.

\begin{propriete}[Rappel 1 : le passif avec 被 (bèi)]
Structure : \cle{Receveur + 被 + Agent + Verbe (+ complément)}. Le sujet grammatical \emph{subit} l'action.\\
Dans le texte : \zh{被选上的李明} (bèi xuǎnshàng de Lǐ Míng) — « Li Ming, qui a été élu ». Ici, l'agent (\zh{大家}, les électeurs) est sous-entendu, ce qui est très fréquent en chinois : \zh{被 + Verbe} seul suffit.\\
Comparaison : le français exige « être + participe passé » (a été élu) ; le chinois garde l'ordre des mots actif et ajoute simplement la marque \zh{被} devant le verbe.
\end{propriete}

\begin{propriete}[Rappel 2 : la concession 虽然...但是... (suīrán... dànshì...)]
Structure : \cle{虽然 + proposition 1, 但是 + proposition 2} — « bien que..., (mais)... ».\\
Dans le texte : \zh{虽然我们还年轻，但是我们从小学习做负责任的公民。} — « Bien que nous soyons encore jeunes, nous apprenons dès l'enfance à devenir des citoyens responsables. »\\
Différence majeure avec le français : en français, on dit « bien que... » \emph{ou} « mais... », jamais les deux ensemble. En chinois, les deux mots se répondent et s'emploient \textbf{ensemble} : c'est une paire obligatoire, comme les deux battants d'une porte.
\end{propriete}

\begin{propriete}[Rappel 3 : la totalité 每...都... (měi... dōu...)]
Structure : \cle{每 + classificateur + Nom + 都 + Verbe} — « chaque... (tous)... ».\\
Dans le texte : \zh{每个同学都有一票} (měi ge tóngxué dōu yǒu yī piào) — « chaque élève a une voix » ; \zh{听取每个人的意见} — « écouter l'avis de chacun ».\\
Piège : \zh{每} ne peut presque jamais se passer de \zh{都} dans la phrase. Dire \zh{每个同学有一票} sans \zh{都} sonne incomplet pour une oreille chinoise.
\end{propriete}

\begin{exemplebox}[Deux phrases modèles à retenir par cœur]
\zh{每个同学都有一票。}\\
\emph{Měi ge tóngxué dōu yǒu yī piào.}\\
« Chaque élève a une voix. »\\[6pt]
\zh{民主就是人人参与。}\\
\emph{Mínzhǔ jiùshì rénrén cānyù.}\\
« La démocratie, c'est la participation de tous. »
\end{exemplebox}

\begin{attention}[票 (piào) : « voix » ou « ticket » ?]
Le caractère \zh{票} a deux sens principaux : « billet, ticket » (\zh{车票} chēpiào, ticket de bus ; \zh{电影票} diànyǐngpiào, billet de cinéma) et « voix, bulletin de vote » dans le contexte électoral. Dans \zh{每个同学都有一票}, il faut traduire « chaque élève a \textbf{une voix} », jamais « un ticket » ! Le contexte décide du sens : c'est la règle d'or de la traduction du chinois vers le français. De même, \zh{投票} (tóupiào, litt. « lancer la voix ») signifie « voter ».
\end{attention}

\courssub{Tableau récapitulatif des structures du texte}

\begin{center}
\begin{tabularx}{\linewidth}{|X|X|X|}
\hline
\rowcolor{popblueL} \textbf{Structure} & \textbf{Exemple du texte (caractères + pinyin)} & \textbf{Traduction} \\
\hline
被 + Verbe (passif) & 被选上的李明 (bèi xuǎnshàng de Lǐ Míng) & Li Ming, qui a été élu \\
\hline
每 + 个 + Nom + 都 & 每个同学都有一票 (měi ge tóngxué dōu yǒu yī piào) & chaque élève a une voix \\
\hline
虽然... 但是... & 虽然我们还年轻，但是... (suīrán wǒmen hái niánqīng, dànshì...) & bien que nous soyons jeunes, (mais)... \\
\hline
是... 也是... & 是公民的权利，也是义务 (shì gōngmín de quánlì, yě shì yìwù) & c'est un droit du citoyen, et aussi un devoir \\
\hline
为 + qqn + Verbe & 为大家服务 (wèi dàjiā fúwù) & servir tout le monde \\
\hline
就是 (insistance) & 民主就是人人参与 (mínzhǔ jiùshì rénrén cānyù) & la démocratie, c'est justement la participation de tous \\
\hline
\end{tabularx}
\end{center}

\courssub{Exercice résolu : retrouver et traduire}

\begin{exoresolu}[Repérage de structures dans le texte]
Énoncé : (1) Retrouvez dans le texte la phrase qui contient la structure passive et traduisez-la. (2) Retrouvez la phrase avec \zh{虽然...但是...} et expliquez pourquoi le français ne traduit qu'un seul des deux mots. (3) Traduisez : \zh{听取每个人的意见。}
\tcblower
Solution détaillée :\\
(1) La phrase passive est \zh{被选上的李明说...} (bèi xuǎnshàng de Lǐ Míng shuō...) : « Li Ming, qui a été élu, a déclaré... ». La marque du passif est \zh{被}, placée devant le verbe \zh{选上} (élire) ; l'agent (les camarades votants) est sous-entendu.\\
(2) \zh{虽然我们还年轻，但是我们从小学习做负责任的公民。} Le chinois emploie la paire \zh{虽然} (bien que) + \zh{但是} (mais) ensemble ; le français n'en traduit qu'un : « Bien que nous soyons jeunes, nous apprenons... » ou « Nous sommes jeunes, mais nous apprenons... ». Traduire les deux (« bien que... mais... ») serait une faute en français.\\
(3) \zh{听取每个人的意见。} (tīngqǔ měi ge rén de yìjiàn) : « écouter l'avis de chacun ». On reconnaît la structure \zh{每 + 个 + 人} (chaque personne) ; ici \zh{都} est absent car le groupe \zh{每个人的意见} est un complément de nom (l'avis de chacun), non un sujet de phrase.
\end{exoresolu}

\begin{aretenir}[L'essentiel de la lecture guidée]
\begin{itemize}
\item Lire à voix haute = grouper par sens + appliquer le sandhi du 3\textsuperscript{e} ton (\zh{选举} → \emph{xuánjǔ}).
\item Un texte nouveau se lit en s'appuyant sur les mots connus ; le contexte fait deviner les mots nouveaux.
\item Trois structures réemployées à savoir repérer : le passif avec \zh{被}, la concession \zh{虽然...但是...}, la totalité \zh{每...都...}.
\item \zh{票} = « voix » dans un contexte d'élection, « ticket » dans un contexte de transport.
\end{itemize}
\end{aretenir}

Maintenant que le texte est lu, compris dans sa structure et débarrassé de ses difficultés grammaticales, passons à la section 3 : le glossaire complet du thème « démocratie », mot par mot.

\coursec{Glossaire du thème : la démocratie}

Après avoir lu et écouté le texte support sur la démocratie, tu as rencontré plusieurs mots nouveaux. Cette section les rassemble, les explique un à un et t'apprend à les prononcer correctement, à les distinguer de leurs faux amis et à les employer dans des phrases. Prends le temps de lire chaque mot à voix haute : en chinois, un mot mal prononcé (mauvais ton) peut devenir un autre mot !

\courssub{Tableau de vocabulaire du thème}

Voici les seize mots essentiels de la leçon. Observe d'abord le tableau, puis lis chaque ligne à voix haute en suivant le pinyin.

\begin{center}
\begin{tabularx}{\linewidth}{|X|X|X|X|}
\hline
\rowcolor{popblueL} Caractères & Pinyin & Nature & Traduction \\
\hline
民主 & mínzhǔ & n. & démocratie \\
\hline
选举 & xuǎnjǔ & n./v. & élection / élire \\
\hline
投票 & tóupiào & v. & voter \\
\hline
票 & piào & n. & bulletin, voix, ticket \\
\hline
公民 & gōngmín & n. & citoyen \\
\hline
权利 & quánlì & n. & droit (légitime) \\
\hline
义务 & yìwù & n. & devoir \\
\hline
自由 & zìyóu & n. & liberté \\
\hline
意见 & yìjiàn & n. & opinion, avis \\
\hline
参与 & cānyù & v. & participer \\
\hline
负责 & fùzé & v./adj. & être responsable (de) \\
\hline
服务 & fúwù & v. & servir \\
\hline
举行 & jǔxíng & v. & organiser (une élection, une cérémonie) \\
\hline
国家 & guójiā & n. & pays, État \\
\hline
法律 & fǎlǜ & n. & loi, législation \\
\hline
人民 & rénmín & n. & peuple \\
\hline
\end{tabularx}
\end{center}

\begin{aretenir}[Trois familles de mots à mémoriser ensemble]
Le vocabulaire de ce thème se range naturellement en trois familles : la \cle{vie électorale} (\zh{民主}, \zh{选举}, \zh{投票}, \zh{票}, \zh{举行}), le \cle{statut du citoyen} (\zh{公民}, \zh{权利}, \zh{义务}, \zh{自由}, \zh{意见}, \zh{参与}), et le \cle{cadre de l'État} (\zh{国家}, \zh{法律}, \zh{人民}, \zh{负责}, \zh{服务}). Apprendre les mots par famille est plus efficace que de les apprendre dans le désordre.
\end{aretenir}

\begin{popfigure}
\begin{tikzpicture}[mindmap, grow cyclic, every node/.style={concept}, root concept/.append style={concept color=popblue, font=\large\bfseries}, level 1/.append style={level distance=3.4cm, sibling angle=120}, level 2/.append style={level distance=2.4cm, sibling angle=45}]
\node[concept color=popblue, text=white, align=center]{民主\\mínzhǔ\\démocratie}
  child[concept color=popgreen] { node {élections}
    child { node[concept color=popgreenL] {选举 xuǎnjǔ} }
    child { node[concept color=popgreenL] {投票 tóupiào} }
    child { node[concept color=popgreenL] {举行 jǔxíng} }
  }
  child[concept color=poporange] { node {citoyen}
    child { node[concept color=poporangeL] {公民 gōngmín} }
    child { node[concept color=poporangeL] {权利 quánlì} }
    child { node[concept color=poporangeL] {义务 yìwù} }
    child { node[concept color=poporangeL] {自由 zìyóu} }
  }
  child[concept color=poppurple] { node {État}
    child { node[concept color=poppurpleL] {国家 guójiā} }
    child { node[concept color=poppurpleL] {法律 fǎlǜ} }
    child { node[concept color=poppurpleL] {人民 rénmín} }
  };
\end{tikzpicture}
\legende{Carte mentale du lexique de la démocratie : trois familles de mots autour de 民主.}
\end{popfigure}

\courssub{Étude des mots clés}

\begin{definition}[公民 gōngmín : citoyen]
Le mot \zh{公民} se compose de deux caractères : \zh{公} (gōng), qui signifie « public, commun », et \zh{民} (mín), qui signifie « peuple ». Un \zh{公民} est donc littéralement un « membre public du peuple » : une personne qui appartient à un État et qui possède des droits et des devoirs. On retrouve \zh{民} dans \zh{人民} (rénmín, le peuple) et dans \zh{民主} (mínzhǔ, démocratie, littéralement « le peuple dirige »).
\end{definition}

\begin{exemplebox}[Deux phrases modèles à connaître par cœur]
\zh{投票是公民的权利。}\\
\emph{Tóupiào shì gōngmín de quánlì.}\\
« Voter est un droit du citoyen. »\\[4pt]
\zh{每个公民都有义务遵守法律。}\\
\emph{Měi ge gōngmín dōu yǒu yìwù zūnshǒu fǎlǜ.}\\
« Tout citoyen a le devoir de respecter la loi. »
\end{exemplebox}

Observe la structure de ces deux phrases : \textbf{Sujet + 是 + complément} pour définir (« voter est un droit »), et \textbf{每 + classificateur + nom + 都 + verbe} pour exprimer « tout / chaque ». Le mot \zh{都} (dōu) est obligatoire après \zh{每} : c'est une erreur classique de francophone de l'oublier.

\begin{attention}[权利 et 权力 : deux mots, une seule prononciation]
\zh{权利} (quánlì, le droit, ce qui est légitimement accordé) et \zh{权力} (quánlì, le pouvoir, l'autorité) se prononcent exactement pareil : ce sont des \cle{homophones parfaits}. Seuls les caractères et le contexte permettent de les distinguer. Compare : \zh{公民有投票的权利。} « Les citoyens ont le droit de voter » (droit) et \zh{总统有权力。} « Le président a le pouvoir » (autorité). À l'écrit comme au Bac, choisir le mauvais caractère est une faute : retiens \zh{利} (bénéfice, intérêt) pour le droit, \zh{力} (force) pour le pouvoir.
\end{attention}

\begin{attention}[意见 et 意思 : ne pas confondre « avis » et « sens »]
\zh{意见} (yìjiàn) signifie « opinion, avis » : \zh{我想听听你的意见。} \emph{Wǒ xiǎng tīngting nǐ de yìjiàn.} « Je voudrais entendre ton avis. » \zh{意思} (yìsi) signifie « sens, signification » : \zh{这个词是什么意思？} \emph{Zhège cí shì shénme yìsi?} « Que signifie ce mot ? » Un francophone qui demande « ton avis » avec \zh{意思} demande en réalité « ton sens », ce qui ne veut rien dire !
\end{attention}

\courssub{Prononciation : les tons difficiles du glossaire}

Trois mots de cette leçon posent des problèmes de tons aux francophones. Travaille-les à voix haute.

\begin{propriete}[Règle d'assimilation des 3\textsuperscript{e} tons]
Quand deux troisièmes tons se suivent, le premier se prononce comme un deuxième ton. Ainsi \zh{选举} s'écrit xuǎnjǔ (3 + 3) mais se prononce \emph{xuánjǔ} (2 + 3). C'est la même règle que dans \zh{你好} (nǐ hǎo, prononcé ní hǎo). Attention : on écrit toujours les tons d'origine dans le pinyin, on ne change que la prononciation.
\end{propriete}

\begin{itemize}
\item \zh{选举} xuǎnjǔ : écris 3 + 3, prononce 2 + 3 (\emph{xuánjǔ}).
\item \zh{权利} quánlì : 2 + 4 ; le deuxième ton monte, le quatrième tombe brusquement. Ne confonds pas avec quánlí qui ne veut rien dire.
\item \zh{义务} yìwù : 4 + 4 ; deux tons descendants secs, comme deux ordres courts. Évite de faire monter le second ton.
\end{itemize}

\begin{exemplebox}[Entraînement rapide]
\zh{公民参与选举。} \emph{Gōngmín cānyù xuǎnjǔ.} « Les citoyens participent aux élections. »\\[3pt]
\zh{国家举行选举，人民投票。} \emph{Guójiā jǔxíng xuǎnjǔ, rénmín tóupiào.} « L'État organise les élections, le peuple vote. »\\[3pt]
\zh{每个人都有发表意见的自由。} \emph{Měi ge rén dōu yǒu fābiǎo yìjiàn de zìyóu.} « Chacun a la liberté d'exprimer son opinion. »
\end{exemplebox}

\courssub{Petit texte de réinvestissement}

\begin{textebox}[Le vote dans notre quartier]
上个星期天，我们区举行了选举。早上八点，很多公民来到学校投票。票上写着候选人的名字。\\
\emph{Shàng ge xīngqītiān, wǒmen qū jǔxíngle xuǎnjǔ. Zǎoshang bā diǎn, hěn duō gōngmín láidào xuéxiào tóupiào. Piào shang xiězhe hòuxuǎnrén de míngzi.}\\
« Dimanche dernier, notre quartier a organisé des élections. À huit heures du matin, de nombreux citoyens sont venus voter à l'école. Sur les bulletins étaient écrits les noms des candidats. »\\[4pt]
投票是公民的权利，也是义务。每个人都有发表意见的自由。年轻人也参与了这次活动。他们说：「我们是国家的人民，我们要为国家服务。」\\
\emph{Tóupiào shì gōngmín de quánlì, yě shì yìwù. Měi ge rén dōu yǒu fābiǎo yìjiàn de zìyóu. Niánqīngrén yě cānyùle zhè cì huódòng. Tāmen shuō : « Wǒmen shì guójiā de rénmín, wǒmen yào wèi guójiā fúwù. »}\\
« Voter est un droit du citoyen, et aussi un devoir. Chacun a la liberté d'exprimer son avis. Les jeunes aussi ont participé à cet événement. Ils ont dit : "Nous sommes le peuple du pays, nous voulons servir le pays." »\\[4pt]
选举以后，大家都要遵守法律，负责地工作。民主需要每一个公民。\\
\emph{Xuǎnjǔ yǐhòu, dàjiā dōu yào zūnshǒu fǎlǜ, fùzé de gōngzuò. Mínzhǔ xūyào měi yī ge gōngmín.}\\
« Après les élections, tout le monde doit respecter la loi et travailler de façon responsable. La démocratie a besoin de chaque citoyen. »
\end{textebox}

\emph{Note :} \zh{候选人} (hòuxuǎnrén, « candidat ») n'est pas dans la liste de base mais est transparent morphologiquement (\zh{候} attendre + \zh{选} élire + \zh{人} personne).

\courssub{Compréhension du texte}

\begin{exercice}[Questions de compréhension]
Réponds en français (ou en chinois si tu le peux) aux questions suivantes d'après le texte ci-dessus.
\begin{enumerate}
\item Quand et où a eu lieu l'élection ?
\item Qu'est-ce qui était écrit sur les bulletins de vote ?
\item Quels sont les deux statuts du vote pour le citoyen mentionnés dans le texte ?
\item Que disent les jeunes à propos de leur identité et de leur volonté ?
\item Quelle est la condition finale pour que la démocratie fonctionne, selon la dernière phrase ?
\end{enumerate}
\end{exercice}

\begin{corrige}[Exercice Compréhension du texte]
\begin{enumerate}
\item L'élection a eu lieu dimanche dernier (\zh{上个星期天}) dans le quartier (\zh{我们区}), plus précisément à l'école (\zh{学校}).
\item Sur les bulletins (\zh{票}) étaient écrits les noms des candidats (\zh{候选人的名字}).
\item Voter est un \textbf{droit} (\zh{权利}) et un \textbf{devoir} (\zh{义务}).
\item Ils disent : « Nous sommes le peuple du pays (\zh{国家的人民}), nous voulons servir le pays (\zh{为国家服务}) ». 
\item La démocratie a besoin de \textbf{chaque citoyen} (\zh{每一个公民}) qui respecte la loi et travaille responsablement.
\end{enumerate}
\end{corrige}

\courssub{Clés (radicaux) et ordre des traits}

Connaître la clé (radical) et l'ordre des traits facilite la mémorisation, la recherche dans le dictionnaire et l'écriture correcte. Voici six caractères stratégiques du thème.

\begin{center}
\begin{tabularx}{\linewidth}{|X|X|X|X|X|}
\hline
\rowcolor{popblueL} Caractère & Pinyin & Clé (Radical) & Nb traits & Ordre des traits (description) \\
\hline
民 & mín & \zh{氏} (shì, clan) & 5 & Horizontal (一), vertical (丨), horizontal (一), horizontal (一), horizontal (一). \emph{Idée : l'œil percé (aveugle) au-dessus du peuple soumis.} \\
\hline
选 & xuǎn & \zh{辶} (chuò, marche) & 9 & Clé \zh{辶} en dernier (point, horizontal, point). Partie droite \zh{巽} : horizontal, vertical, horizontal, horizontal, horizontal, vertical, horizontal. \\
\hline
票 & piào & \zh{示} (shì, autel/rite) & 11 & Haut \zh{示} : point, horizontal, vertical, horizontal, horizontal. Bas \zh{是} : vertical, horizontal, vertical, horizontal, horizontal, horizontal. \\
\hline
权 & quán & \zh{木} (mù, bois) & 6 & Horizontal, vertical, horizontal, horizontal, vertical, horizontal. \emph{Anciennement \zh{權} (bois + poids), simplifié en bois seul.} \\
\hline
法 & fǎ & \zh{氵} (shuǐ, eau) & 8 & Points d'eau (丶 丶 丨), puis \zh{去} : horizontal, vertical, horizontal, horizontal, horizontal. \emph{L'eau qui va (équité).} \\
\hline
责 & zé & \zh{贝} (bèi, cauri/valeur) & 11 & Haut \zh{贝} : horizontal, vertical, horizontal, horizontal, point, horizontal. Bas \zh{朿} (cī, épine) : vertical, horizontal, vertical, horizontal. \\
\hline
\end{tabularx}
\end{center}

\begin{methode}[Comment réviser l'écriture]
\begin{enumerate}
\item Écris chaque caractère 5 fois en respectant scrupuleusement l'ordre des traits ci-dessus.
\item Dis le nom des traits à voix haute en écrivant (ex. : « horizontal, vertical, horizontal… »).
\item Associe le caractère à sa clé : \zh{民} (clé \zh{氏}), \zh{选} (clé \zh{辶}), \zh{票} (clé \zh{示}), \zh{权} (clé \zh{木}), \zh{法} (clé \zh{氵}), \zh{责} (clé \zh{贝}).
\item Teste-toi : cache le pinyin, regarde le caractère, donne le son, le sens et la clé.
\end{enumerate}
\end{methode}

\courssub{Collocations et production guidée}

\begin{center}
\begin{tabularx}{\linewidth}{|X|X|}
\hline
\rowcolor{popblueL} Mot pivot & Collocations fréquentes (Caractères | Pinyin | Traduction) \\
\hline
选举 & \zh{举行选举} (jǔxíng xuǎnjǔ, organiser des élections) ; \zh{参加选举} (cānjiā xuǎnjǔ, participer aux élections) \\
\hline
投票 & \zh{去投票} (qù tóupiào, aller voter) ; \zh{投票率} (tóupiào lǜ, taux de participation) \\
\hline
公民 & \zh{公民权利} (gōngmín quánlì, droits civiques) ; \zh{合格公民} (hégé gōngmín, citoyen modèle) \\
\hline
权利/义务 & \zh{享有权利} (xiǎngyǒu quánlì, jouir de droits) ; \zh{履行义务} (lǚxíng yìwù, s'acquitter de ses devoirs) \\
\hline
自由 & \zh{言论自由} (yánlùn zìyóu, liberté d'expression) ; \zh{自由表达} (zìyóu biǎodá, s'exprimer librement) \\
\hline
法律 & \zh{遵守法律} (zūnshǒu fǎlǜ, respecter la loi) ; \zh{法律面前人人平等} (fǎlǜ miànqián rénrén píngděng, égalité devant la loi) \\
\hline
负责 & \zh{负责任} (fùzé rèn, responsable) ; \zh{负责这件事} (fùzé zhè jiàn shì, être chargé de cette affaire) \\
\hline
服务 & \zh{为人民服务} (wèi rénmín fúwù, servir le peuple) ; \zh{服务态度} (fúwù tàidù, attitude de service) \\
\hline
\end{tabularx}
\end{center}

\begin{exercice}[Production écrite guidée : « Mon engagement citoyen »]
Rédige un court paragraphe (60--80 caractères chinois) sur le thème : « En tant que jeune citoyen, comment participer à la vie démocratique ? ». Utilise \textbf{au moins 6 mots} du glossaire (surligne-les ou mets-les en gras dans ta copie). Respecte l'ordre des mots chinois (Sujet -- Temps -- Lieu -- Manière -- Verbe -- Complément).
\end{exercice}

\begin{corrige}[Exercice Production écrite guidée]
\emph{Modèle de réponse :}\\
\zh{作为青年公民，我要参加选举，行使投票权利，履行义务。我也要遵守法律，自由发表意见，为国家服务。民主需要我们每个人负责地参与。}\\
\emph{Zuòwéi qīngnián gōngmín, wǒ yào cānjiā xuǎnjǔ, xíngshǐ tóupiào quánlì, lǚxíng yìwù. Wǒ yě yào zūnshǒu fǎlǜ, zìyóu fābiǎo yìjiàn, wèi guójiā fúwù. Mínzhǔ xūyào wǒmen měi ge rén fùzé de cānyù.}\\
« En tant que jeune citoyen, je dois participer aux élections, exercer mon droit de vote, m'acquitter de mes devoirs. Je dois aussi respecter la loi, exprimer librement mon opinion, servir le pays. La démocratie a besoin que chacun de nous participe de façon responsable. »
\\[4pt]
\emph{Mots du glossaire utilisés :} \zh{公民}, \zh{选举}, \zh{投票} (dans \zh{投票权利}), \zh{权利}, \zh{义务}, \zh{法律}, \zh{自由}, \zh{意见}, \zh{国家}, \zh{服务}, \zh{民主}, \zh{负责}, \zh{参与}. (13 mots utilisés).
\end{corrige}

\begin{experience}[Débat structuré : « Vote : droit ou devoir ? »]
En classe entière, divisez-vous en deux groupes.
\begin{itemize}
\item \textbf{Groupe A (Droit) :} Défendez l'idée que voter est un \zh{权利} (liberté de choisir, pas de contrainte). Utilisez : \zh{自由}, \zh{意见}, \zh{选择}.
\item \textbf{Groupe B (Devoir) :} Défendez l'idée que voter est un \zh{义务} (responsabilité envers la \zh{国家}, le \zh{人民}). Utilisez : \zh{负责}, \zh{法律}, \zh{公民}.
\end{itemize}
Chaque groupe prépare 3 arguments en chinois (phrases courtes). Le professeur note le vocabulaire du thème employé correctement. Objectif : réemployer \zh{权利/义务}, \zh{自由/法律}, \zh{参与/负责} à l'oral.
\end{experience}

\begin{exoresolu}[Complète avec le mot du glossaire qui convient]
Énoncé. Complète chaque phrase avec le mot qui convient : \zh{权利}, \zh{义务}, \zh{选举}, \zh{意见}, \zh{法律}.\\
1. 明天我们学校举行……，我们要选班长。\\
2. 遵守……是每个公民的责任。\\
3. 投票是公民的……。\\
4. 帮助父母是孩子的……。\\
5. 对这个计划，你有什么……？
\tcblower
Solution détaillée. 1. \zh{选举} : on élit (选) le chef de classe, c'est une élection. 2. \zh{法律} : on « respecte » (遵守) la loi ; le verbe 遵守 appelle presque toujours 法律 ou 规则. 3. \zh{权利} : voter est un droit, pas un pouvoir (attention à ne pas écrire 权力 !). 4. \zh{义务} : aider ses parents est présenté comme un devoir. 5. \zh{意见} : « quel avis as-tu sur ce plan ? » ; 意思 (sens) serait un faux ami ici.
\end{exoresolu}

\begin{experience}[Quiz flash en binômes : caractères → pinyin → traduction]
Par groupes de deux, utilisez seize petites cartes (une par mot du glossaire). Élève A montre le caractère (par exemple \zh{投票}) ; élève B doit donner en moins de cinq secondes le pinyin avec les tons (\emph{tóupiào}) puis la traduction (« voter »). Un point par réponse complète et correcte. Après huit cartes, inversez les rôles. Variante difficile : élève A dit la traduction française, élève B doit donner le caractère et le pinyin. Le binôme qui obtient le plus de points gagne le titre de « meilleur citoyen du lexique » !
\end{experience}

\begin{savaistu}[Vote et écriture de la loi]
En Chine comme au Cameroun, l'âge minimum pour voter aux élections locales est de dix-huit ans : à dix-huit ans, on devient pleinement \zh{公民}. Et savais-tu que le caractère \zh{法} (fǎ, la loi) contient la clé de l'eau \zh{氵} ? Dans la tradition chinoise ancienne, l'eau symbolise l'équité : la loi doit être « plane comme la surface de l'eau », c'est-à-dire la même pour tous. Une belle image pour retenir à la fois le caractère et l'idée de justice !
\end{savaistu}

\begin{aretenir}[Ce qu'il faut retenir de ce glossaire]
Seize mots, trois familles (élections, citoyen, État), deux pièges (\zh{权利}/\zh{权力} et \zh{意见}/\zh{意思}), une règle de prononciation (3 + 3 → 2 + 3 dans \zh{选举}), six caractères analysés par clé et ordre des traits. Si tu sais lire, prononcer, écrire et employer ces mots dans les phrases modèles et le texte de réinvestissement, tu es prêt pour l'expression écrite et orale du module.
\end{aretenir}

\coursec{Clés et ordre des traits}

Après avoir découvert le glossaire du thème de la démocratie, il est temps de passer à l'\emph{écriture} : au baccalauréat, certains caractères du programme sont exigibles en production écrite. Cette section vous apprend à écrire correctement cinq caractères essentiels du chapitre : \zh{民} (mín), \zh{主} (zhǔ), \zh{选} (xuǎn), \zh{权} (quán) et \zh{票} (piào). Pour chacun, nous identifierons la \cle{clé} (radical), qui donne le sens, et nous décrirons l'\cle{ordre des traits}, qui garantit une écriture rapide, équilibrée et lisible.

\courssub{Pourquoi apprendre les clés et l'ordre des traits ?}

Un caractère chinois n'est pas un dessin libre : c'est une construction réglée. Deux principes organisent tout :

\begin{itemize}
\item La \cle{clé} (部首 \zh{bùshǒu}) : c'est l'élément qui classe le caractère dans le dictionnaire et qui indique souvent la famille de sens. Exemple : la clé \zh{木} (bois) dans \zh{权}, la clé \zh{辶} (mouvement) dans \zh{选}.
\item L'\cle{ordre des traits} (笔顺 \zh{bǐshùn}) : chaque caractère s'écrit selon des règles générales : de haut en bas, de gauche à droite, l'horizontal avant le vertical, l'extérieur avant l'intérieur, et certaines clés (comme \zh{辶}) se tracent en dernier.
\end{itemize}

\begin{methode}[Comment mémoriser durablement un caractère]
\begin{enumerate}
\item \textbf{Observer} : identifier la clé (qui donne le sens) et la partie restante (souvent phonétique). Exemple : \zh{权} = \zh{木} (clé : bois) + \zh{又} (partie phonétique et graphique).
\item \textbf{Compter} : connaître le nombre exact de traits (ex. \zh{主} : 5 traits).
\item \textbf{Tracer} : écrire le caractère 5 fois dans une grille 米字格, en comptant les traits à voix haute : « un, deux, trois… » (一、二、三…).
\item \textbf{Comparer} : vérifier que l'on ne confond pas avec un caractère proche (\zh{主} / \zh{王}).
\item \textbf{Employer} : écrire une phrase entière avec le mot, par exemple \zh{人民有选举的权利。} (Rénmín yǒu xuǎnjǔ de quánlì. — « Le peuple a le droit de vote. »)
\end{enumerate}
\end{methode}

\courssub{Étude caractère par caractère}

\textbf{1. 民 (mín) — « peuple » — 5 traits}\par
Clé : \zh{氏} (shì, clan) ; certains dictionnaires classent \zh{民} sous sa propre clé. Ce caractère ancien désigne le peuple, la population. On le retrouve dans \zh{人民} (rénmín, le peuple), \zh{民主} (mínzhǔ, la démocratie), \zh{公民} (gōngmín, le citoyen).\par
Ordre des traits :
\begin{enumerate}
\item 横折 (héngzhé) : trait horizontal suivi d'un angle vertical — la partie supérieure ;
\item 横 (héng) : trait horizontal ;
\item 竖提 (shùtí) : trait vertical terminé par une petite montée vers la droite ;
\item 横 (héng) : trait horizontal ;
\item 斜钩 (xiégōu) : grand trait oblique terminé par un crochet.
\end{enumerate}

\textbf{2. 主 (zhǔ) — « maître, principal » — 5 traits}\par
Clé : \zh{丶} (diǎn, le point). Le caractère se décompose en \zh{丶} + \zh{王} : on écrit d'abord le point, puis \zh{王} (4 traits : trois horizontaux et un vertical).\par
Ordre des traits :
\begin{enumerate}
\item 点 (diǎn) : le point en haut ;
\item 横 (héng) : premier horizontal ;
\item 横 (héng) : deuxième horizontal ;
\item 竖 (shù) : le vertical central ;
\item 横 (héng) : l'horizontal du bas.
\end{enumerate}

\begin{attention}[主 a 5 traits, pas 4 !]
Le point \zh{丶} compte comme un trait à part entière et s'écrit \emph{en premier}. Erreur fréquente des francophones : oublier le point et écrire \zh{王} (wáng, « roi », 4 traits) à la place de \zh{主} (zhǔ, « maître »). Résultat : \zh{民主} (démocratie) devient \zh{民王}, un mot qui n'existe pas ! Retenez : le peuple (\zh{民}) a un point d'autorité au-dessus du roi (\zh{王}) : c'est le « maître » (\zh{主}).
\end{attention}

\textbf{3. 选 (xuǎn) — « choisir, élire » — 9 traits}\par
Clé : \zh{辶} (chuò), la clé du mouvement, que l'on retrouve dans \zh{进} (jìn, entrer), \zh{这} (zhè, ceci), \zh{道} (dào, chemin). La partie intérieure est \zh{先} (xiān, « d'abord »), qui donne le son.\par
Règle absolue : on écrit \textbf{d'abord} la partie intérieure \zh{先} (6 traits), \textbf{puis} la clé \zh{辶} (3 traits : point, angle, et long trait ondulé final) \emph{en dernier}.

\begin{propriete}[Règle d'écriture : la clé 辶 se trace toujours en dernier]
Pour tous les caractères contenant la clé du mouvement \zh{辶} (appelée « clé chuò »), on écrit d'abord la partie intérieure, puis la clé \zh{辶} en trois traits pour terminer : 点 (point) → 横折折撇 (angle double) → 捺 (long trait ondulé qui « porte » le caractère).\par
Exemples : \zh{选} (xuǎn, choisir), \zh{进} (jìn, entrer), \zh{这} (zhè, ceci), \zh{道} (dào, chemin), \zh{远} (yuǎn, loin).\par
Exception à ne pas inventer : il n'y en a pas — cette règle est sans exception ; c'est l'une des règles les plus régulières de l'écriture chinoise.
\end{propriete}

\textbf{4. 权 (quán) — « droit, pouvoir » — 6 traits}\par
Clé : \zh{木} (mù, bois), placée à gauche ; à droite, \zh{又} (yòu). Règle générale : de gauche à droite — on écrit d'abord \zh{木} (4 traits : horizontal, vertical, oblique gauche, oblique droite ramenée à un point en position de clé), puis \zh{又} (2 traits). On le retrouve dans \zh{权利} (quánlì, le droit) et \zh{权力} (quánlì, le pouvoir) — deux homophones à distinguer par le sens !

\textbf{5. 票 (piào) — « bulletin, billet » — 11 traits}\par
Clé : \zh{示} (shì, montrer, l'autel), placée en bas. Règle générale : de haut en bas — on écrit d'abord la partie supérieure \zh{覀} (xī, 6 traits), puis \zh{示} (5 traits) en dessous. On le retrouve dans \zh{选票} (xuǎnpiào, le bulletin de vote) et \zh{车票} (chēpiào, le ticket de transport).

\courssub{Tableau récapitulatif}

\begin{center}
\begin{tabularx}{\linewidth}{|X|X|X|X|}
\hline
\rowcolor{popblueL} Caractère & Pinyin et sens & Clé & Nombre de traits et règle dominante \\
\hline
民 & mín, peuple & 氏 & 5 traits ; finir par le 斜钩 oblique \\
\hline
主 & zhǔ, maître & 丶 & 5 traits ; le point d'abord, puis 王 \\
\hline
选 & xuǎn, choisir & 辶 & 9 traits ; 先 d'abord, 辶 en dernier \\
\hline
权 & quán, droit & 木 & 6 traits ; gauche puis droite \\
\hline
票 & piào, bulletin & 示 & 11 traits ; haut (覀) puis bas (示) \\
\hline
\end{tabularx}
\end{center}

\courssub{Schémas de tracé en grille 米字格}

La grille 米字格 (« grille du caractère mǐ ») est le cahier d'écriture traditionnel : les diagonales et les médianes aident à placer chaque trait au bon endroit.

\begin{popfigure}
\begin{tikzpicture}[scale=0.9]
% Grille 1 : 主 étape 1
\begin{scope}
\draw[popline] (0,0) rectangle (2,2);
\draw[popline, dashed] (0,1)--(2,1); \draw[popline, dashed] (1,0)--(1,2);
\draw[popline, dashed] (0,0)--(2,2); \draw[popline, dashed] (0,2)--(2,0);
\draw[popink, very thick] (1,1.75)--(1.15,1.6);
\node[popmuted, font=\small] at (1,-0.4) {1. 点};
\end{scope}
% Grille 2 : 主 étapes 2
\begin{scope}[xshift=2.8cm]
\draw[popline] (0,0) rectangle (2,2);
\draw[popline, dashed] (0,1)--(2,1); \draw[popline, dashed] (1,0)--(1,2);
\draw[popline, dashed] (0,0)--(2,2); \draw[popline, dashed] (0,2)--(2,0);
\draw[popink, very thick] (1,1.75)--(1.15,1.6);
\draw[popink, very thick] (0.45,1.35)--(1.55,1.35);
\node[popmuted, font=\small] at (1,-0.4) {2. 横};
\end{scope}
% Grille 3 : 主 étape 3
\begin{scope}[xshift=5.6cm]
\draw[popline] (0,0) rectangle (2,2);
\draw[popline, dashed] (0,1)--(2,1); \draw[popline, dashed] (1,0)--(1,2);
\draw[popline, dashed] (0,0)--(2,2); \draw[popline, dashed] (0,2)--(2,0);
\draw[popink, very thick] (1,1.75)--(1.15,1.6);
\draw[popink, very thick] (0.45,1.35)--(1.55,1.35);
\draw[popink, very thick] (0.55,0.95)--(1.45,0.95);
\node[popmuted, font=\small] at (1,-0.4) {3. 横};
\end{scope}
% Grille 4 : 主 étape 4
\begin{scope}[xshift=8.4cm]
\draw[popline] (0,0) rectangle (2,2);
\draw[popline, dashed] (0,1)--(2,1); \draw[popline, dashed] (1,0)--(1,2);
\draw[popline, dashed] (0,0)--(2,2); \draw[popline, dashed] (0,2)--(2,0);
\draw[popink, very thick] (1,1.75)--(1.15,1.6);
\draw[popink, very thick] (0.45,1.35)--(1.55,1.35);
\draw[popink, very thick] (0.55,0.95)--(1.45,0.95);
\draw[popink, very thick] (1,1.35)--(1,0.45);
\node[popmuted, font=\small] at (1,-0.4) {4. 竖};
\end{scope}
% Grille 5 : 主 complet
\begin{scope}[xshift=11.2cm]
\draw[popline] (0,0) rectangle (2,2);
\draw[popline, dashed] (0,1)--(2,1); \draw[popline, dashed] (1,0)--(1,2);
\draw[popline, dashed] (0,0)--(2,2); \draw[popline, dashed] (0,2)--(2,0);
\draw[popink, very thick] (1,1.75)--(1.15,1.6);
\draw[popink, very thick] (0.45,1.35)--(1.55,1.35);
\draw[popink, very thick] (0.55,0.95)--(1.45,0.95);
\draw[popink, very thick] (1,1.35)--(1,0.45);
\draw[popink, very thick] (0.35,0.45)--(1.65,0.45);
\node[popmuted, font=\small] at (1,-0.4) {5. 横 → 主};
\end{scope}
\end{tikzpicture}
\legende{Tracé de 主 (zhǔ) en 5 étapes : le point 丶 d'abord, puis 王.}
\end{popfigure}

\begin{popfigure}
\begin{tikzpicture}[scale=0.9]
% Grille : 选 avec 先 d'abord
\begin{scope}
\draw[popline] (0,0) rectangle (2.4,2.4);
\draw[popline, dashed] (0,1.2)--(2.4,1.2); \draw[popline, dashed] (1.2,0)--(1.2,2.4);
\draw[popline, dashed] (0,0)--(2.4,2.4); \draw[popline, dashed] (0,2.4)--(2.4,0);
\node[popink] at (1.35,1.45) {\zh{先}};
\node[popmuted, font=\small] at (1.2,-0.4) {Étape 1 : écrire 先 (6 traits)};
\end{scope}
% Flèche
\draw[-{Stealth[length=3mm]}, poporange, very thick] (2.7,1.2)--(3.6,1.2);
% Grille : 选 complet avec 辶 en orange
\begin{scope}[xshift=3.9cm]
\draw[popline] (0,0) rectangle (2.4,2.4);
\draw[popline, dashed] (0,1.2)--(2.4,1.2); \draw[popline, dashed] (1.2,0)--(1.2,2.4);
\draw[popline, dashed] (0,0)--(2.4,2.4); \draw[popline, dashed] (0,2.4)--(2.4,0);
\node[popink] at (1.35,1.45) {\zh{先}};
\node[poporange] at (0.85,0.95) {\zh{辶}};
\node[popmuted, font=\small] at (1.2,-0.4) {Étape 2 : 辶 en dernier (3 traits)};
\end{scope}
% Résultat
\draw[-{Stealth[length=3mm]}, poporange, very thick] (6.6,1.2)--(7.5,1.2);
\begin{scope}[xshift=7.8cm]
\draw[popline] (0,0) rectangle (2.4,2.4);
\draw[popline, dashed] (0,1.2)--(2.4,1.2); \draw[popline, dashed] (1.2,0)--(1.2,2.4);
\draw[popline, dashed] (0,0)--(2.4,2.4); \draw[popline, dashed] (0,2.4)--(2.4,0);
\node[popink, font=\Large] at (1.2,1.2) {\zh{选}};
\node[popmuted, font=\small] at (1.2,-0.4) {Résultat : 选 (9 traits)};
\end{scope}
\end{tikzpicture}
\legende{Tracé de 选 (xuǎn) : la partie intérieure 先 d'abord, la clé 辶 en dernier.}
\end{popfigure}

\courssub{Caractères proches à ne pas confondre}

\begin{center}
\begin{tabularx}{\linewidth}{|X|X|X|}
\hline
\rowcolor{popblueL} Paire & Différence graphique & Moyen mnémotechnique \\
\hline
主 zhǔ (maître) / 王 wáng (roi) & 主 a un point 丶 en plus (5 traits contre 4) & Le maître a « un point » d'avance sur le roi \\
\hline
民 mín (peuple) / 氏 shì (clan) & 民 a un trait horizontal de plus en haut (横) & Le peuple est « plus grand » que le clan : un trait de plus \\
\hline
\end{tabularx}
\end{center}

\begin{exemplebox}[Trois mots clés exigibles en écriture au Bac]
\zh{民主} mínzhǔ — « démocratie » : \zh{民主很重要。} (Mínzhǔ hěn zhòngyào. — « La démocratie est très importante. »)\\[2pt]
\zh{选举} xuǎnjǔ — « élection » : \zh{喀麦隆有选举。} (Kāmàilóng yǒu xuǎnjǔ. — « Le Cameroun organise des élections. »)\\[2pt]
\zh{权利} quánlì — « droit » : \zh{公民有投票的权利。} (Gōngmín yǒu tóupiào de quánlì. — « Les citoyens ont le droit de voter. »)
\end{exemplebox}

\courssub{Exercice de tracé et exercice résolu}

\begin{exercice}[Tracé guidé en classe]
Chaque élève trace dans sa grille 米字格 les cinq caractères \zh{民}, \zh{主}, \zh{选}, \zh{权}, \zh{票}, cinq fois chacun, en comptant les traits à voix haute en chinois : 一、二、三、四、五… Le professeur vérifie que la clé \zh{辶} de \zh{选} est bien tracée en dernier et que le point de \zh{主} est écrit en premier.
\end{exercice}

\begin{exoresolu}[Identifier la clé et compter les traits]
Énoncé : pour chaque caractère, donnez la clé et le nombre de traits : a) \zh{权} ; b) \zh{选} ; c) \zh{主} ; d) \zh{票}.
\tcblower
Solution détaillée :
\begin{enumerate}
\item \zh{权} : clé \zh{木} (bois) à gauche ; 6 traits (4 pour \zh{木} + 2 pour \zh{又}), tracé de gauche à droite.
\item \zh{选} : clé \zh{辶} (mouvement) ; 9 traits — \zh{先} (6 traits) d'abord, \zh{辶} (3 traits) en dernier.
\item \zh{主} : clé \zh{丶} (point) ; 5 traits — le point en premier, puis \zh{王}. Attention : ce n'est pas 4 traits !
\item \zh{票} : clé \zh{示} ; 11 traits — \zh{覀} (6 traits) en haut d'abord, puis \zh{示} (5 traits) en bas, de haut en bas.
\end{enumerate}
\end{exoresolu}

\begin{attention}[Pièges fréquents des francophones]
\begin{itemize}
\item Écrire \zh{选} en traçant la clé \zh{辶} en premier : c'est la faute la plus fréquente ; la clé du mouvement se trace toujours en dernier.
\item Oublier le point de \zh{主} ou l'écrire après \zh{王} : le point s'écrit en premier.
\item Confondre \zh{民} et \zh{氏} : vérifiez le trait horizontal supplémentaire de \zh{民}.
\item Tracer \zh{票} de bas en haut : en chinois, on écrit toujours de haut en bas.
\end{itemize}
\end{attention}

\begin{aretenir}[L'essentiel de la section]
Cinq caractères exigibles : \zh{民} (5 traits, clé 氏), \zh{主} (5 traits, clé 丶, point en premier), \zh{选} (9 traits, clé 辶 tracée en dernier), \zh{权} (6 traits, clé 木, gauche-droite), \zh{票} (11 traits, clé 示, haut-bas). Règles d'or : gauche avant droite, haut avant bas, clé 辶 en dernier. Ces cinq caractères permettent d'écrire les trois mots clés du chapitre : \zh{民主} (démocratie), \zh{选举} (élection), \zh{权利} (droit).
\end{aretenir}

\coursec{Compréhension du texte}

\courssub{De l'écriture à la lecture : transition}

Dans la section précédente, nous avons appris à écrire les caractères du thème — \zh{民} (mín, le peuple), \zh{主} (zhǔ, maître), \zh{选} (xuǎn, choisir), \zh{举} (jǔ, lever), \zh{票} (piào, bulletin de vote) — en respectant les clés et l'ordre des traits. Savoir tracer un caractère, c'est bien ; mais au baccalauréat, c'est la \cle{compréhension de texte} qui rapporte le plus de points. Nous allons maintenant relire le texte support de la leçon — l'élection du délégué de classe — et apprendre à répondre aux questions de compréhension \emph{comme un examinateur l'attend} : par des phrases complètes, en chinois, en reprenant les mots de la question.

\begin{textebox}[Rappel du texte support : 我们班的选举 (L'élection de notre classe)]
上个星期五，我们学校举行了班长选举。每个同学都有一票。\\
Shàng gè xīngqīwǔ, wǒmen xuéxiào jǔxíng le bānzhǎng xuǎnjǔ. Měi gè tóngxué dōu yǒu yí piào.\\
Vendredi dernier, notre école a organisé l'élection du délégué de classe. Chaque élève avait une voix.\\[4pt]
同学们认真地投票。最后，李明被选上了，他很高兴。\\
Tóngxuémen rènzhēn de tóupiào. Zuìhòu, Lǐ Míng bèi xuǎnshàng le, tā hěn gāoxìng.\\
Les élèves ont voté sérieusement. Finalement, Li Ming a été élu, il était très content.\\[4pt]
老师说：「选举是每个人的权利，也是责任。」李明要为大家服务，帮助同学，做一个好班长。\\
Lǎoshī shuō : « Xuǎnjǔ shì měi gè rén de quánlì, yě shì zérèn. » Lǐ Míng yào wèi dàjiā fúwù, bāngzhù tóngxué, zuò yí gè hǎo bānzhǎng.\\
Le professeur a dit : « Voter est un droit de chacun, et aussi une responsabilité. » Li Ming va servir tout le monde, aider ses camarades et être un bon délégué.\\[4pt]
民主就是人民当家作主：大家都有说话的权利，大家一起决定重要的事情。\\
Mínzhǔ jiùshì rénmín dāngjiā zuòzhǔ : dàjiā dōu yǒu shuōhuà de quánlì, dàjiā yìqǐ juédìng zhòngyào de shìqing.\\
La démocratie, c'est le peuple qui est maître chez lui : tout le monde a le droit de parler, et tout le monde décide ensemble des choses importantes.
\end{textebox}

\courssub{Méthode : répondre par une phrase complète}

\begin{methode}[Répondre à une question de compréhension en quatre étapes]
\begin{enumerate}
\item \textbf{Lire la question} et identifier le mot interrogatif : \zh{什么} (shénme, quoi), \zh{谁} (shéi, qui), \zh{几} (jǐ, combien), \zh{为什么} (wèishénme, pourquoi).
\item \textbf{Repérer le mot-clé} de la question (le nom ou le verbe principal) : dans \zh{学校举行了什么？}, le mot-clé est \zh{举行} (jǔxíng, organiser).
\item \textbf{Localiser la phrase du texte} qui contient ce mot-clé : ici, la première phrase \zh{我们学校举行了班长选举}.
\item \textbf{Rédiger une phrase complète} en reprenant les mots de la question et en remplaçant le mot interrogatif par l'information trouvée : \zh{学校举行了班长选举。}
\end{enumerate}
\end{methode}

\begin{popfigure}
\begin{tikzpicture}[node distance=0.55cm,
  etape/.style={rectangle, rounded corners=3pt, draw=popblue, fill=popblueL, align=center, minimum width=4.6cm, minimum height=0.85cm, font=\small},
  etapeLast/.style={rectangle, rounded corners=3pt, draw=popgreen, fill=popgreenL, align=center, minimum width=4.6cm, minimum height=0.85cm, font=\small},
  fleche/.style={-{Stealth[length=2.5mm]}, thick, poporange}]
\node[etape, align=center] (e1){1. Lire la question\\\zh{学校举行了什么？}};
\node[etape, below=of e1, align=center] (e2){2. Repérer le mot-clé\\\zh{举行} (organiser)};
\node[etape, below=of e2, align=center] (e3){3. Localiser la phrase dans le texte\\\zh{我们学校举行了班长选举}};
\node[etapeLast, below=of e3, align=center] (e4){4. Rédiger une phrase complète\\\zh{学校举行了班长选举。}};
\draw[fleche] (e1) -- (e2);
\draw[fleche] (e2) -- (e3);
\draw[fleche] (e3) -- (e4);
\end{tikzpicture}
\legende{Organigramme de la méthode de réponse : de la question à la phrase complète.}
\end{popfigure}

\begin{propriete}[La règle d'or : la question est le squelette de la réponse]
En chinois comme en français, une bonne réponse de compréhension \textbf{reprend les mots de la question} et remplace le mot interrogatif par la réponse. Le mot interrogatif occupe en chinois la \emph{même place} que le mot qu'il remplace — contrairement au français, il n'y a pas d'inversion ni de déplacement :
\begin{itemize}
\item \zh{学校举行了\cle{什么}？} → \zh{学校举行了\cle{班长选举}。} (le mot interrogatif et la réponse sont tous les deux après le verbe)
\item \zh{\cle{谁}被选上了？} → \zh{\cle{李明}被选上了。} (tous les deux en position de sujet)
\end{itemize}
C'est un grand avantage du chinois pour le francophone : l'ordre des mots ne change jamais entre la question et la réponse.
\end{propriete}

\begin{exemplebox}[Observer avant de répondre]
\zh{谁被选上了？} \emph{Shéi bèi xuǎnshàng le?} — Qui a été élu ?\\
Réponse complète : \zh{李明被选上了。} \emph{Lǐ Míng bèi xuǎnshàng le.} — Li Ming a été élu.\\[6pt]
\zh{每个同学有几票？} \emph{Měi gè tóngxué yǒu jǐ piào?} — Combien de voix chaque élève avait-il ?\\
Réponse complète : \zh{每个同学都有一票。} \emph{Měi gè tóngxué dōu yǒu yí piào.} — Chaque élève avait une voix.\\[6pt]
\zh{老师说选举是什么？} \emph{Lǎoshī shuō xuǎnjǔ shì shénme?} — Selon le professeur, qu'est-ce que voter ?\\
Réponse complète : \zh{老师说选举是每个人的权利，也是责任。} \emph{Lǎoshī shuō xuǎnjǔ shì měi gè rén de quánlì, yě shì zérèn.} — Le professeur a dit que voter est un droit de chacun et aussi une responsabilité.
\end{exemplebox}

\begin{attention}[Piège du Bac : la réponse en un seul mot]
À la question \zh{什么}, beaucoup d'élèves francophones répondent seulement \zh{班长选举} (« l'élection du délégué »), comme on le ferait à l'oral en français. \textbf{C'est insuffisant à l'écrit du baccalauréat} : une phrase complète avec \cle{sujet + verbe} est exigée.\\
✗ \zh{班长选举。} (réponse télégraphique, points perdus)\\
✓ \zh{学校举行了班长选举。} (sujet \zh{学校} + verbe \zh{举行} + complément)\\
Même vigilance avec \zh{谁} : ne répondez pas \zh{李明。} mais \zh{李明被选上了。}
\end{attention}

\courssub{Questions de compréhension : travail individuel}

\begin{textebox}[À traiter individuellement en 10 minutes]
\textbf{Questions de compréhension globale}\\[2pt]
1. 学校举行了什么？\emph{Xuéxiào jǔxíng le shénme?} — Qu'est-ce que l'école a organisé ?\\
2. 谁被选上了？\emph{Shéi bèi xuǎnshàng le?} — Qui a été élu ?\\[6pt]
\textbf{Questions de détail}\\[2pt]
3. 每个同学有几票？\emph{Měi gè tóngxué yǒu jǐ piào?} — Combien de voix chaque élève avait-il ?\\
4. 老师说选举是什么？\emph{Lǎoshī shuō xuǎnjǔ shì shénme?} — Selon le professeur, qu'est-ce que voter ?\\
5. 李明要做什么？\emph{Lǐ Míng yào zuò shénme?} — Que va faire Li Ming ?\\[6pt]
\textbf{Question de réflexion}\\[2pt]
6. 根据课文，民主是什么？\emph{Gēnjù kèwén, mínzhǔ shì shénme?} — Selon le texte, qu'est-ce que la démocratie ?\\[6pt]
\textbf{Vrai ou Faux (对 / 错)}\\[2pt]
a. 李明是老师。\emph{Lǐ Míng shì lǎoshī.} — Li Ming est professeur.\\
b. 同学们认真投票。\emph{Tóngxuémen rènzhēn tóupiào.} — Les élèves votent sérieusement.\\
c. 民主就是一个人决定。\emph{Mínzhǔ jiùshì yí gè rén juédìng.} — La démocratie, c'est une seule personne qui décide.
\end{textebox}

\begin{exoresolu}[Questions 1 et 2 : correction détaillée]
Énoncé : répondez par des phrases complètes. 1) \zh{学校举行了什么？} 2) \zh{谁被选上了？}
\tcblower
\textbf{Question 1.} Mot interrogatif : \zh{什么} (shénme, quoi). Mot-clé : \zh{举行} (jǔxíng, organiser). On le retrouve dans la première phrase du texte : \zh{我们学校举行了班长选举}. On remplace \zh{什么} par \zh{班长选举} :\\
\zh{学校举行了班长选举。} \emph{Xuéxiào jǔxíng le bānzhǎng xuǎnjǔ.} — L'école a organisé l'élection du délégué de classe.\\[6pt]
\textbf{Question 2.} Mot interrogatif : \zh{谁} (shéi, qui), en position de sujet. Le texte dit : \zh{李明被选上了}. On remplace \zh{谁} par \zh{李明} :\\
\zh{李明被选上了。} \emph{Lǐ Míng bèi xuǎnshàng le.} — Li Ming a été élu.\\
Remarque : la structure passive avec \zh{被} (bèi) doit être conservée dans la réponse ; ne dites pas \zh{李明选上了}, qui signifierait « Li Ming a choisi ».
\end{exoresolu}

\begin{corrige}[Questions 3 à 6 et Vrai/Faux]
\textbf{3.} \zh{每个同学都有一票。} \emph{Měi gè tóngxué dōu yǒu yí piào.} — Chaque élève avait une voix. (On reprend \zh{每个同学都有…票} et on remplace \zh{几} par \zh{一} — attention au changement de ton : \zh{yí} devant \zh{piào} 4\up{ème} ton).\\[4pt]
\textbf{4.} \zh{老师说选举是每个人的权利，也是责任。} \emph{Lǎoshī shuō xuǎnjǔ shì měi gè rén de quánlì, yě shì zérèn.} — Le professeur a dit que voter est un droit de chacun et aussi une responsabilité. (On reprend la citation du texte, introduite par \zh{老师说}.)\\[4pt]
\textbf{5.} \zh{李明要为大家服务，帮助同学，做一个好班长。} \emph{Lǐ Míng yào wèi dàjiā fúwù, bāngzhù tóngxué, zuò yí gè hǎo bānzhǎng.} — Li Ming va servir tout le monde, aider ses camarades et être un bon délégué.\\[4pt]
\textbf{6.} \zh{根据课文，民主就是人民当家作主：大家都有说话的权利，大家一起决定重要的事情。} \emph{Gēnjù kèwén, mínzhǔ jiùshì rénmín dāngjiā zuòzhǔ : dàjiā dōu yǒu shuōhuà de quánlì, dàjiā yìqǐ juédìng zhòngyào de shìqing.} — Selon le texte, la démocratie, c'est le peuple qui est maître : chacun a le droit de parler et tous décident ensemble des choses importantes. (Pour une question de réflexion, on reprend la phrase de conclusion du texte en la commençant par \zh{根据课文}, « selon le texte ».)\\[4pt]
\textbf{Vrai/Faux.} a) \textbf{Faux} (错) : Li Ming est un élève, élu délégué ; le professeur est \zh{老师}. b) \textbf{Vrai} (对) : le texte dit \zh{同学们认真地投票} (l'adverbe \zh{地} est omis à l'oral mais la phrase est correcte). c) \textbf{Faux} (错) : le texte dit le contraire — \zh{大家一起决定}, « tout le monde décide ensemble » ; une seule personne qui décide, ce n'est pas la démocratie.
\end{corrige}

\courssub{Correction collective et reformulation}

\begin{methode}[Conduire la correction au tableau]
\begin{enumerate}
\item Un élève écrit sa réponse au tableau ; la classe vérifie d'abord la \textbf{forme} : y a-t-il un sujet, un verbe, une phrase complète ?
\item On vérifie ensuite le \textbf{contenu} : l'information vient-elle bien du texte ? Le mot interrogatif a-t-il été correctement remplacé ?
\item En cas de réponse fausse ou incomplète, on \textbf{reformule} sans se moquer : on montre la phrase du texte, on souligne le mot-clé, et l'élève reconstruit sa phrase.
\end{enumerate}
\end{methode}

\begin{exemplebox}[Reformulation de réponses fausses]
Réponse d'élève : ✗ \zh{举行了班长选举。} (sujet manquant)\\
Reformulation : ✓ \zh{学校举行了班长选举。} — Qui a organisé ? L'école : il faut le sujet \zh{学校}.\\[6pt]
Réponse d'élève : ✗ \zh{李明选上了。} (sens changé : « Li Ming a choisi »)\\
Reformulation : ✓ \zh{李明被选上了。} — Le texte utilise le passif avec \zh{被} : Li Ming \emph{a été élu}, il ne choisit pas.\\[6pt]
Réponse d'élève : ✗ \zh{民主是一个人决定。} (contre-sens)\\
Reformulation : ✓ \zh{民主是大家一起决定。} — Relire la dernière phrase du texte : la démocratie, c'est décider \emph{ensemble}.
\end{exemplebox}

\begin{center}
\begin{tabularx}{\linewidth}{|X|X|X|}
\hline
\rowcolor{popblueL} Mot interrogatif & Question → Réponse (modèle) & Traduction de la réponse \\
\hline
\zh{什么} shénme (quoi) & \zh{学校举行了什么？→ 学校举行了班长选举。} & L'école a organisé l'élection du délégué. \\
\hline
\zh{谁} shéi (qui) & \zh{谁被选上了？→ 李明被选上了。} & Li Ming a été élu. \\
\hline
\zh{几} jǐ (combien) & \zh{每个同学有几票？→ 每个同学都有一票。} & Chaque élève avait une voix. \\
\hline
\zh{什么} + \zh{是} (définition) & \zh{民主是什么？→ 民主就是人民当家作主。} & La démocratie, c'est le peuple maître. \\
\hline
\zh{要} + \zh{什么} (intention) & \zh{李明要做什么？→ 李明要为大家服务。} & Li Ming va servir tout le monde. \\
\hline
\end{tabularx}
\end{center}

\begin{attention}[Comparer le français et le chinois]
En français, la question déplace les mots : « \emph{Qui} a été élu ? » met l'interrogatif en tête, alors que la réponse met « Li Ming » en sujet. En chinois, \textbf{rien ne bouge} : \zh{谁被选上了？} et \zh{李明被选上了。} ont exactement le même ordre. Profitez-en : recopiez la question, effacez le mot interrogatif, écrivez la réponse à sa place. C'est la stratégie la plus sûre — à condition de ne pas oublier le sujet et de garder les particules comme \zh{被} et \zh{了}.
\end{attention}

\begin{exercice}[À vous de jouer]
Sans regarder les corrigés, répondez par des phrases complètes en chinois : 1) \zh{谁有一票？} (Qui avait une voix ?) 2) \zh{老师说了什么？} (Qu'a dit le professeur ?) 3) \zh{李明高兴吗？} (Li Ming était-il content ? — répondez avec \zh{很高兴}). Puis échangez votre copie avec un voisin et corrigez-vous mutuellement selon la méthode en quatre étapes.
\end{exercice}

\begin{aretenir}[L'essentiel de la section]
\begin{itemize}
\item Une réponse de compréhension = \textbf{phrase complète} (sujet + verbe) qui reprend les mots de la question.
\item Le mot interrogatif chinois (\zh{什么}, \zh{谁}, \zh{几}) occupe la place exacte du mot-réponse : on le remplace, on ne déplace rien.
\item Méthode en quatre étapes : lire la question → repérer le mot-clé → localiser la phrase → rédiger la phrase complète.
\item Garder les structures du texte (\zh{被} pour le passif, \zh{了} pour l'accompli) et commencer une réponse de réflexion par \zh{根据课文} (selon le texte).
\end{itemize}
\end{aretenir}

\coursec{Collocations et production guidée}

Dans la section précédente, nous avons vérifié notre compréhension du texte sur la démocratie en répondant aux questions. Nous allons maintenant passer de la \emph{réception} à la \emph{production} : il ne suffit pas de connaître les mots \zh{选举} (élection), \zh{权利} (droit) ou \zh{义务} (devoir) ; il faut savoir avec quel \cle{verbe} chacun s'emploie. C'est ce qu'on appelle les \cle{collocations} : des couples verbe + nom qui vont naturellement ensemble. C'est la dernière étape avant de rédiger toi-même quelques phrases sur la vie démocratique de ta classe.

\courssub{Observer d'abord : quels verbes avec quels noms ?}

Lis attentivement ces phrases tirées de situations de la vie scolaire et civique. Souligne mentalement le verbe qui accompagne chaque nom en gras.

\begin{exemplebox}[Des couples qui vont ensemble]
\zh{我们学校每年举行班长选举。} \emph{Wǒmen xuéxiào měi nián jǔxíng bānzhǎng xuǎnjǔ.} — Notre école organise chaque année l'élection du délégué.\\
\zh{全班同学都参加了投票。} \emph{Quán bān tóngxué dōu cānjiā le tóupiào.} — Tous les élèves de la classe ont participé au vote.\\
\zh{公民应该行使自己的权利。} \emph{Gōngmín yīnggāi xíngshǐ zìjǐ de quánlì.} — Les citoyens doivent exercer leurs droits.\\
\zh{每个人也要履行自己的义务。} \emph{Měi gè rén yě yào lǚxíng zìjǐ de yìwù.} — Chacun doit aussi remplir ses devoirs.\\
\zh{班长常常听取同学们的意见。} \emph{Bānzhǎng chángcháng tīngqǔ tóngxuémen de yìjiàn.} — Le délégué écoute souvent les avis de ses camarades.\\
\zh{他愿意为大家服务。} \emph{Tā yuànyì wèi dàjiā fúwù.} — Il accepte de servir tout le monde.\\
\zh{好公民遵守法律。} \emph{Hǎo gōngmín zūnshǒu fǎlǜ.} — Un bon citoyen respecte la loi.
\end{exemplebox}

Tu remarques que chaque nom a « son » verbe : \zh{举行} avec \zh{选举}, \zh{参加} avec \zh{投票}, \zh{行使} avec \zh{权利}, \zh{履行} avec \zh{义务}, \zh{听取} avec \zh{意见}, \zh{服务} avec \zh{为大家}, \zh{遵守} avec \zh{法律}. Ces associations ne se devinent pas à partir du français : il faut les apprendre comme des blocs.

\begin{propriete}[Les collocations : des paires fixes verbe + complément]
En chinois, le verbe et son complément d'objet forment souvent des \cle{paires fixes} qu'on ne peut pas modifier librement. On dit \zh{行使权利} « exercer un droit », jamais \zh{做权利} ; on dit \zh{履行义务} « remplir un devoir », jamais \zh{做义务}.\\
\textbf{Règle pratique :} n'apprends jamais un nom abstrait seul ; apprends-le \textbf{avec son verbe}, comme un seul bloc : \zh{举行选举}, \zh{参加投票}, \zh{行使权利}, \zh{履行义务}, \zh{听取意见}, \zh{遵守法律}.\\
\textbf{Comparaison avec le français :} le français fonctionne pareil ! On dit « tenir une promesse » et non « faire une promesse », « prendre une décision » et non « faire une décision ». Le chinois pousse simplement ce principe plus loin, et le verbe « tout faire » \zh{做} ne peut pas remplacer les verbes spécialisés.
\end{propriete}

\begin{center}
\begin{tabularx}{\linewidth}{|X|X|X|X|}
\hline
\rowcolor{popblueL} Collocation & Pinyin & Nature & Traduction \\
\hline
\zh{举行选举} & jǔxíng xuǎnjǔ & VO & organiser une élection \\
\hline
\zh{参加投票} & cānjiā tóupiào & VO & participer au vote \\
\hline
\zh{行使权利} & xíngshǐ quánlì & VO & exercer un droit \\
\hline
\zh{履行义务} & lǚxíng yìwù & VO & remplir un devoir \\
\hline
\zh{听取意见} & tīngqǔ yìjiàn & VO & écouter les avis \\
\hline
\zh{为大家服务} & wèi dàjiā fúwù & Prep + VO & servir tout le monde \\
\hline
\zh{遵守法律} & zūnshǒu fǎlǜ & VO & respecter la loi \\
\hline
\end{tabularx}
\end{center}

\begin{popfigure}
\begin{tikzpicture}[font=\small]
% colonne verbes
\node[fill=popblueL, rounded corners, align=center, minimum width=2.2cm] (v1) at (0,4) {\zh{举行} jǔxíng};
\node[fill=popblueL, rounded corners, align=center, minimum width=2.2cm] (v2) at (0,3) {\zh{参加} cānjiā};
\node[fill=popblueL, rounded corners, align=center, minimum width=2.2cm] (v3) at (0,2) {\zh{行使} xíngshǐ};
\node[fill=popblueL, rounded corners, align=center, minimum width=2.2cm] (v4) at (0,1) {\zh{履行} lǚxíng};
\node[fill=popblueL, rounded corners, align=center, minimum width=2.2cm] (v5) at (0,0) {\zh{听取} tīngqǔ};
% colonne noms (ordre mélangé)
\node[fill=poporangeL, rounded corners, align=center, minimum width=2.2cm] (n1) at (6,4) {\zh{义务} yìwù};
\node[fill=poporangeL, rounded corners, align=center, minimum width=2.2cm] (n2) at (6,3) {\zh{选举} xuǎnjǔ};
\node[fill=poporangeL, rounded corners, align=center, minimum width=2.2cm] (n3) at (6,2) {\zh{意见} yìjiàn};
\node[fill=poporangeL, rounded corners, align=center, minimum width=2.2cm] (n4) at (6,1) {\zh{投票} tóupiào};
\node[fill=poporangeL, rounded corners, align=center, minimum width=2.2cm] (n5) at (6,0) {\zh{权利} quánlì};
% flèches
\draw[-{Stealth}, very thick, popgreen] (v1) -- (n2);
\draw[-{Stealth}, very thick, popgreen] (v2) -- (n4);
\draw[-{Stealth}, very thick, popgreen] (v3) -- (n5);
\draw[-{Stealth}, very thick, popgreen] (v4) -- (n1);
\draw[-{Stealth}, very thick, popgreen] (v5) -- (n3);
\node[align=center, popmuted] at (0,4.9) {Verbes};
\node[align=center, popmuted] at (6,4.9) {Noms};
\end{tikzpicture}
\legende{Les paires fixes : relie chaque verbe à son nom. Attention, les noms sont volontairement mélangés !}
\end{popfigure}

\courssub{Chaque collocation dans une phrase}

\begin{exemplebox}[Sept phrases modèles à réemployer]
\zh{我们班上星期举行了选举。} \emph{Wǒmen bān shàng xīngqī jǔxíng le xuǎnjǔ.} — Notre classe a organisé une élection la semaine dernière.\\
\zh{三十五个人参加了投票。} \emph{Sānshíwǔ gè rén cānjiā le tóupiào.} — Trente-cinq personnes ont participé au vote.\\
\zh{我们应该行使自己的权利。} \emph{Wǒmen yīnggāi xíngshǐ zìjǐ de quánlì.} — Nous devons exercer nos droits.\\
\zh{学生也要履行自己的义务。} \emph{Xuéshēng yě yào lǚxíng zìjǐ de yìwù.} — Les élèves doivent aussi remplir leurs devoirs.\\
\zh{班长要听取大家的意见。} \emph{Bānzhǎng yào tīngqǔ dàjiā de yìjiàn.} — Le délégué doit écouter les avis de tous.\\
\zh{老师为大家服务，我们都很感谢她。} \emph{Lǎoshī wèi dàjiā fúwù, wǒmen dōu hěn gǎnxiè tā.} — La professeure sert tout le monde, nous lui sommes tous reconnaissants.\\
\zh{每个人都应该遵守法律。} \emph{Měi gè rén dōu yīnggāi zūnshǒu fǎlǜ.} — Chacun doit respecter la loi.
\end{exemplebox}

\begin{attention}[La place du complément de lieu et de temps]
En français, on dit « organiser une élection \emph{dans notre école} » : le complément de lieu vient après. En chinois, c'est l'inverse : le complément de lieu ou de temps se place \textbf{avant le verbe}.\\
Correct : \zh{在我们学校举行选举。} \emph{Zài wǒmen xuéxiào jǔxíng xuǎnjǔ.}\\
Incorrect : \zh{举行选举在我们学校。}\\
Schéma : \textbf{[Sujet] + 在/给/跟 + Lieu/Temps + Verbe + Objet}. De même : \zh{班长在教室里听取我们的意见。} — Le délégué écoute nos avis dans la salle de classe.
\end{attention}

\begin{savaistu}[为人民服务 : une devise célèbre]
\zh{为人民服务} \emph{wèi rénmín fúwù}, « servir le peuple », est une devise très célèbre en Chine. Elle est inscrite en grands caractères à l'entrée de nombreux bâtiments publics, écoles et administrations. La structure \zh{为 + personne + 服务} est exactement celle de notre collocation \zh{为大家服务} : \zh{为} (wèi) signifie « pour, au bénéfice de ». Au Cameroun, on retrouve une idée voisine dans l'idéal du service public et du travail communautaire au village.
\end{savaistu}

\courssub{Production guidée, étape 1 : phrases à trous}

\begin{exoresolu}[Complète avec la bonne collocation]
Énoncé — Complète chaque phrase avec la collocation qui convient : \zh{举行选举}, \zh{参加投票}, \zh{行使权利}, \zh{履行义务}, \zh{听取意见}.\\
1. 明天我们班要 \_\_\_\_\_\_\_\_，选新班长。\\
2. 每个同学都应该 \_\_\_\_\_\_\_\_，不能放弃。\\
3. 好班长常常 \_\_\_\_\_\_\_\_ 同学们的想法。\\
4. 公民有权利，也有义务：我们要 \_\_\_\_\_\_\_\_，也要 \_\_\_\_\_\_\_\_。
\tcblower
Solution détaillée :\\
1. \zh{举行选举} — on « organise » une élection pour choisir (选) le nouveau délégué. \emph{Míngtiān wǒmen bān yào jǔxíng xuǎnjǔ, xuǎn xīn bānzhǎng.}\\
2. \zh{参加投票} — chaque élève doit « participer au vote » ; l'indice est \zh{不能放弃} « ne pas renoncer ». \emph{Měi gè tóngxué dōu yīnggāi cānjiā tóupiào, bù néng fàngqì.}\\
3. \zh{听取意见} — le bon délégué « écoute les avis » ; attention, \zh{意见} se place après \zh{听取}, et \zh{同学们的想法} précise de qui viennent les avis. \emph{Hǎo bānzhǎng chángcháng tīngqǔ tóngxuémen de xiǎngfǎ.}\\
4. \zh{行使权利} puis \zh{履行义务} — la phrase oppose droits et devoirs : on « exerce » un droit, on « remplit » un devoir. Piège classique : ne pas inverser les deux verbes ! \emph{Gōngmín yǒu quánlì, yě yǒu yìwù: wǒmen yào xíngshǐ quánlì, yě yào lǚxíng yìwù.}
\end{exoresolu}

\begin{exercice}[À toi de jouer]
Traduis en chinois avec les collocations : 1. Notre école organise une élection chaque année. 2. Tous les élèves ont participé au vote. 3. Le délégué doit écouter les avis de tout le monde et servir la classe. 4. Un bon citoyen respecte la loi.
\end{exercice}

\begin{corrige}[Exercice « À toi de jouer »]
1. \zh{我们学校每年举行选举。} \emph{Wǒmen xuéxiào měi nián jǔxíng xuǎnjǔ.} — le complément de temps \zh{每年} se place avant le verbe.\\
2. \zh{所有同学都参加了投票。} \emph{Suǒyǒu tóngxué dōu cānjiā le tóupiào.} — \zh{了} marque l'action accomplie.\\
3. \zh{班长要听取大家的意见，为大家服务。} \emph{Bānzhǎng yào tīngqǔ dàjiā de yìjiàn, wèi dàjiā fúwù.}\\
4. \zh{好公民遵守法律。} \emph{Hǎo gōngmín zūnshǒu fǎlǜ.}
\end{corrige}

\courssub{Production guidée, étape 2 : rédiger sur l'élection du délégué}

\begin{methode}[Rédiger 3 à 4 phrases avec la trame]
\begin{enumerate}
\item \textbf{Phrase 1 — les faits :} \zh{我们班举行了……} (notre classe a organisé…). Ajoute quand : \zh{上星期}, \zh{昨天}, \zh{开学的时候}.
\item \textbf{Phrase 2 — la participation :} \zh{每个同学……} (chaque élève…). Réemploie \zh{参加了投票} ou \zh{行使了自己的权利}.
\item \textbf{Phrase 3 — le rôle du délégué :} \zh{新班长要……} avec \zh{听取大家的意见} et \zh{为大家服务}.
\item \textbf{Phrase 4 — ton opinion :} \zh{我觉得民主就是……} (je pense que la démocratie, c'est…). Termine par une idée simple : \zh{听取每个人的意见}, \zh{大家一起做决定}, \zh{遵守规则}.
\end{enumerate}
Relis-toi : verbe avant l'objet, complément de lieu/temps avant le verbe, tons corrects sur \zh{xuǎnjǔ} (3-3) et \zh{yìjiàn} (4-4).
\end{methode}

\begin{textebox}[Modèle de production : l'élection dans ma classe]
\zh{上星期五，我们班举行了班长选举。每个同学都参加了投票，行使了自己的权利。三十八个同学里面，玛丽得了二十票，所以她当了新班长。}\\\zh{Shàng xīngqīwǔ, wǒmen bān jǔxíng le bānzhǎng xuǎnjǔ. Měi gè tóngxué dōu cānjiā le tóupiào, xíngshǐ le zìjǐ de quánlì. Sānshíbā gè tóngxué lǐmiàn, Mǎlì dé le èrshí piào, suǒyǐ tā dāng le xīn bānzhǎng.}\\\shortstack[l]{\zh{Vendredi dernier, notre classe a organisé l'élection du délégué. Chaque élève a participé au vote et a exercé son droit. Sur trente-huit élèves, Marie a obtenu vingt voix, elle est donc devenue la nouvelle déléguée.}\\ \zh{新班长说，她要常常听取大家的意见，为大家服务。我觉得民主就是：每个人都有权利说话，每个人也要履行自己的义务，遵守班级的规则。}\\ \zh{Xīn bānzhǎng shuō, tā yào chángcháng tīngqǔ dàjiā de yìjiàn, wèi dàjiā fúwù. Wǒ juéde mínzhǔ jiùshì: měi gè rén dōu yǒu quánlì shuōhuà, měi gè rén yě yào lǚxíng zìjǐ de yìwù, zūnshǒu bānjí de guīzé.}\\ \zh{La nouvelle déléguée a dit qu'elle écouterait souvent les avis de tous et servirait tout le monde. Je pense que la démocratie, c'est : chacun a le droit de parler, chacun doit aussi remplir ses devoirs et respecter les règles de la classe.}}
\end{textebox}

\courssub{Pour aller plus loin : vie civique au Cameroun et en Chine}

Ce vocabulaire te servira aussi à l'examen pour comparer les deux pays. Au Cameroun, les citoyens élisent le président de la République, les députés et les conseillers municipaux ; dans les lycées, chaque classe élit son délégué, exactement comme dans notre texte. En Chine, il existe des élections locales et des assemblées populaires, dont l'Assemblée nationale populaire, \zh{人民代表大会} \emph{rénmín dàibiǎo dàhuì}, où siègent des représentants (\zh{代表}, « délégué, représentant » — le même mot que dans \zh{课代表}, le délégué de matière !). Dans les deux pays, l'idée de base se dit avec les mêmes collocations : \zh{参加投票}, \zh{行使权利}, \zh{听取意见}, \zh{为大家服务}.

\begin{experience}[Mise en commun en classe]
Deux ou trois volontaires lisent leur production à voix haute. La classe écoute et vérifie trois points : 1) les tons, surtout \zh{xuǎnjǔ}, \zh{quánlì}, \zh{yìwù} ; 2) l'ordre des mots (complément de lieu/temps avant le verbe) ; 3) les bonnes paires verbe + nom. Le professeur corrige avec bienveillance : on reformule la phrase fautive, puis l'élève la répète correctement. Ensuite, chacun améliore sa rédaction avant de la recopier au propre.
\end{experience}

\begin{aretenir}[L'essentiel de la section]
\begin{itemize}
\item Les collocations sont des \cle{blocs fixes} : \zh{举行选举} (\emph{jǔxíng xuǎnjǔ}), \zh{参加投票} (\emph{cānjiā tóupiào}), \zh{行使权利} (\emph{xíngshǐ quánlì}), \zh{履行义务} (\emph{lǚxíng yìwù}), \zh{听取意见} (\emph{tīngqǔ yìjiàn}), \zh{为大家服务} (\emph{wèi dàjiā fúwù}), \zh{遵守法律} (\emph{zūnshǒu fǎlǜ}). Jamais \zh{做权利} !
\item Le complément de lieu ou de temps se place \cle{avant le verbe} : \zh{在我们学校举行选举}.
\item Pour rédiger : faits (\zh{举行了……}) → participation (\zh{每个同学……}) → opinion (\zh{我觉得民主就是……}).
\item \zh{为人民服务} (\emph{wèi rénmín fúwù}) : « servir le peuple », devise célèbre en Chine.
\end{itemize}
\end{aretenir}

\newpage

\coursec{Activité d'intégration}

\coursec{Leçon 21 — Vocabulaire et compréhension de texte : démocratie}

\courssub{Mise en route et découverte du thème}
\begin{experience}[Brainstorming : « Qu'est-ce que la démocratie pour toi ? »]
Le professeur affiche au tableau deux photos : (1) des élèves levant la main pour voter au Lycée Bilingue de Yaoundé, (2) une séance de l'Assemblée populaire nationale à Beijing (人民代表大会). En groupes de 4, les élèves notent 3 mots-clés en français, puis tentent de les dire en chinois avec le dictionnaire. Mise en commun : construction d'une carte mentale au tableau (mots attendus : \zh{投票}, \zh{选举}, \zh{权利}, \zh{义务}, \zh{代表}, \zh{公平}).
\begin{center}
\begin{popfigure}
\begin{tikzpicture}[mindmap, grow cyclic, every node/.style=concept, concept color=popblue, level 1/.append style={level distance=3.5cm, sibling angle=90}, level 2/.append style={level distance=2.5cm}]
\node {民主 \textit{mínzhǔ}}
    child { node {权利 \textit{quánlì}} }
    child { node {义务 \textit{yìwù}} }
    child { node {选举 \textit{xuǎnjǔ}} }
    child { node {代表 \textit{dàibiǎo}} };
\end{tikzpicture}
\legende{Carte mentale introductive : le champ lexical de la démocratie}
\end{popfigure}
\end{center}
\end{experience}

\courssub{Texte support : lecture guidée}
\begin{textebox}[Texte original : La démocratie dans nos lycées — Regards croisés Cameroun/Chine]
\zh{在喀麦隆的雅温得双语高中，也在中国的北京四中，高三的学生都要选举班长。} \\
\zh{十月十五日，雅温得双语高中高三A班在教室305举行了班级大使选举。} \\
\zh{三位候选人作了竞选演讲，介绍各自的竞选纲领：组织中非文化周，改善食堂饭菜，开设汉语角。} \\
\zh{全体同学行使民主权利，把选票投进票箱。监票人公开点票，宣布结果。} \\
\zh{马莉同学高票当选，被同学们选上了。她表示：``我为全班同学服务，这是我的承诺。''} \\
\zh{这次选举让我们体会到：民主不仅是一项权利，更是一份义务。} \\

\textit{Pinyin :} \\
Zài Kāmàilóng de Yǎwēndé Shuāngyǔ Gāozhōng, yě zài Zhōngguó de Běijīng Sìzhōng, gāosān de xuéshēng dōu yào xuǎnjǔ bānzhǎng. \\
Shíyuè shíwǔ rì, Yǎwēndé Shuāngyǔ Gāozhōng Gāosān A bān zài jiàoshì sān líng wǔ jǔxíng le bānjí dàshǐ xuǎnjǔ. \\
Sān wèi hòuxuǎnrén zuò le jìngxuǎn yǎnjiǎng, jièshào gèzì de jìngxuǎn gānglǐng: zǔzhī Zhōng-Fēi wénhuà zhōu, gǎishàn shítáng fàncài, kāishè Hànyǔ jiǎo. \\
Quántǐ tóngxué xíngshǐ le mínzhǔ quánlì, bǎ xuǎnpiào tóu jìn piàoxiāng. Jiānpiàorén gōngkāi diǎnpiào, xuānbù jiéguǒ. \\
Mǎ Lì tóngxué gāopiào dāngxuǎn, bèi tóngxuémen xuǎn shàng le. Tā biǎoshì: ``Wǒ wèi quánbān tóngxué fúwù, zhè shì wǒ de chéngnuò.'' \\
Zhè cì xuǎnjǔ ràng wǒmen tǐhuì dào: mínzhǔ bùjǐn shì yī xiàng quánlì, gèng shì yī fèn yìwù. \\

\textit{Traduction :} Au Lycée Bilingue de Yaoundé au Cameroun, comme au Lycée n°4 de Pékin en Chine, les élèves de Terminale doivent élire leur délégué de classe. Le 15 octobre, la classe de Terminale A du Lycée Bilingue de Yaoundé a organisé l'élection de l'ambassadeur de classe en salle 305. Trois candidats ont fait leur discours de campagne, présentant leurs programmes : organiser la semaine culturelle sino-camerounaise, améliorer les repas de la cantine, ouvrir un coin chinois. Tous les élèves ont exercé leur droit démocratique, glissant leur bulletin dans l'urne. Les scrutateurs ont dépouillé publiquement et annoncé les résultats. L'élève Marie a été élue à une large majorité, élue par ses camarades. Elle a déclaré : « Je servirai tous les camarades de la classe, c'est ma promesse. » Cette élection nous a fait comprendre : la démocratie n'est pas seulement un droit, c'est aussi un devoir.
\end{textebox}

\begin{methode}[Comment aborder un texte de compréhension en chinois (niveau HSK 3-4)]
\begin{enumerate}
    \item \textbf{Repérage global :} Identifier le thème (ici : élection scolaire), les protagonistes (élèves, candidats, professeurs), le lieu/date (雅温得双语高中, 十月十五日).
    \item \textbf{Lexique connu vs inconnu :} Surligner les mots du vocabulaire cible (\zh{选举}, \zh{候选人}, \zh{纲领}...). Deviner les nouveaux par le contexte (\zh{高票} = grand nombre de voix, \zh{体会到} = réaliser/profondément ressentir).
    \item \textbf{Structure des phrases :} Repérer les mots-outils (\zh{在} lieu, \zh{把} disposition, \zh{被} passif, \zh{不但/不仅...更/而且} progression).
    \item \textbf{Relecture à voix haute :} Travailler les tons (3e ton \zh{选} xuǎn, 4e ton \zh{票} piào, 3e ton \zh{主} zhǔ) et les groupements rythmiques.
    \item \textbf{Réponse aux questions :} Citer le texte (copier les caractères) ou reformuler avec ses mots.
\end{enumerate}
\end{methode}

\courssub{Glossaire du thème : Vocabulaire de la démocratie et des élections}
\begin{center}
\begin{tabularx}{\linewidth}{|X|X|X|X|}
\hline
\rowcolor{popblueL}
\textbf{Caractères} & \textbf{Pinyin} & \textbf{Nature} & \textbf{Traduction} \\
\hline
民主 & mínzhǔ & Adj/N & démocratique, démocratie \\
选举 & xuǎnjǔ & V/N & élire, élection \\
投票 & tóupiào & VO/N & voter, vote \\
候选人 & hòuxuǎnrén & N & candidat \\
选票 & xuǎnpiào & N & bulletin de vote \\
票箱 & piàoxiāng & N & urne, boîte à bulletins \\
点票 & diǎnpiào & VO/N & dépouiller, dépouillement \\
当选 & dāngxuǎn & V & être élu (intransitif) \\
宣布 & xuānbù & V & annoncer, proclamer \\
代表 & dàibiǎo & V/N & représenter, représentant/délégué \\
服务 & fúwù & V/N & servir, service \\
权利 & quánlì & N & droit \\
义务 & yìwù & N & devoir, obligation \\
竞选 & jìngxuǎn & V & faire campagne, briguer un mandat \\
纲领 & gānglǐng & N & programme politique, plateforme \\
承诺 & chéngnuò & V/N & promettre, promesse \\
监票人 & jiānpiàorén & N & scrutateur \\
记录员 & jìlùyuán & N & secrétaire de séance \\
行使 & xíngshǐ & V & exercer (un droit) \\
高票 & gāopiào & Adj/N & grand nombre de voix (souvent \zh{高票当选}) \\
体会 & tǐhuì & V & ressentir, réaliser, prendre conscience \\
\hline
\end{tabularx}
\end{center}

\courssub{Clés (radicaux) et ordre des traits : Focus sur 6 caractères clés}
\begin{definition}[Analyse graphique et mémorisation]
\begin{itemize}
    \item \textbf{\zh{选} (xuǎn, élire) :} Clé \zh{辶} (chuò, marche) — idée de chemin/choix. Phonétique \zh{巽} (xùn). Ordre : haut \zh{⺈} + \zh{人} + \zh{彐} → bas \zh{辶} (point, horizontal, pli, point). \textit{Mnémo :} Choisir le bon chemin.
    \item \textbf{\zh{票} (piào, billet/ticket) :} Clé \zh{示} (shì, montrer/rite) — document officiel. Phonétique \zh{是} (shì). Ordre : \zh{示} (haut) puis \zh{是} (bas : \zh{日} + \zh{勿}). \textit{Attention :} \zh{勿} s'écrit horizontal, vertical, horizontal, point.
    \item \textbf{\zh{举} (jǔ, lever/organiser) :} Clé \zh{廾} (gǒng, deux mains). Haut : \zh{兴} (xīng) sans le trait final. Ordre : \zh{廾} (bas, deux mains qui lèvent) puis le haut. \textit{Verbe clé :} \zh{举行} (organiser/avoir lieu), \zh{举手} (lever la main).
    \item \textbf{\zh{权} (quán, droit/pouvoir) :} Clé \zh{木} (mù, bois). Phonétique \zh{泉} (quán, source). Ordre : \zh{木} (haut) + \zh{泉} (bas : \zh{白} + \zh{水}). \textit{Lien :} Le bois (tablette de loi) + la source (origine) = autorité.
    \item \textbf{\zh{民} (mín, peuple) :} Clé \zh{氏} (shì, clan) ou \zh{乙}. Pictogramme : un œil \zh{目} percé (aveugle/obéissant) ou aiguillonné. Ordre : \zh{乙} + \zh{目} (vertical, horizontal, vertical, horizontaux).
    \item \textbf{\zh{主} (zhǔ, maître/principal) :} Clé \zh{丶} (zhǔ, point) ou \zh{王}. Pictogramme : une lampe \zh{主} sur un support. Ordre : point (haut), horizontal, vertical, horizontal (bas). \textit{Composés :} \zh{民主} (peuple maître), \zh{主席} (président), \zh{主张} (proposition).
\end{itemize}
\end{definition}

\begin{experience}[Atelier écriture : « Course aux radicaux »]
En binômes : l'un dicte le pinyin (ex: \textit{xuǎnjǔ}), l'autre écrit le caractère sur ardoise en annonçant l'ordre des traits et la clé. Inversion des rôles. Le professeur vérifie la proportion des caractères (\zh{选} : partie gauche fine, droite large).
\end{experience}

\courssub{Compréhension du texte}
\begin{exercice}[Questions sur le texte support (niveau A2/B1)]
Répondez en français ou en chinois (caractères + pinyin) en vous basant \textbf{uniquement} sur le texte.
\begin{enumerate}
    \item Où et quand a lieu l'élection au Cameroun ? (Citer la phrase exacte).
    \item Combien y a-t-il de candidats ? Quel classifieur est utilisé ?
    \item Citez deux promesses (\zh{承诺}) figurant dans les programmes (\zh{纲领}).
    \item Quelle structure grammaticale exprime l'action de voter ? (Indiquez le mot-outil et la phrase du texte).
    \item Comment s'appelle la personne qui compte les voix ? Quel est son rôle ?
    \item Quelle structure exprime le résultat pour Marie ? (Citez les deux formes utilisées dans le texte).
    \item Traduisez la dernière phrase : \zh{民主不仅是一项权利，更是一份义务。}
    \item \textbf{Expression personnelle (oral) :} Selon toi, pourquoi le texte dit-il que la démocratie est aussi une « obligation » (\zh{义务}) ? Donne un exemple concret au lycée.
\end{enumerate}
\end{exercice}

\begin{corrige}[Exercice Compréhension]
\begin{enumerate}
    \item \zh{十月十五日，雅温得双语高中高三A班在教室305举行了班级大使选举。} (Le 15 octobre, la classe de Terminale A du Lycée Bilingue de Yaoundé a organisé l'élection de l'ambassadeur de classe en salle 305.)
    \item Trois candidats (\zh{三位候选人}). Classifieur \zh{位} (wèi), poli pour les personnes.
    \item Exemples : \zh{组织中非文化周} (organiser la semaine culturelle sino-camerounaise), \zh{改善食堂饭菜} (améliorer les repas de la cantine), \zh{开设汉语角} (ouvrir un coin chinois).
    \item Structure \zh{把} (disposition) : \zh{把选票投进票箱} (mettre le bulletin dans l'urne). Focus sur la manipulation de l'objet \zh{选票}.
    \item \zh{监票人} (jiānpiàorén, scrutateur). Rôle : \zh{公开点票} (dépouiller publiquement), \zh{宣布结果} (annoncer les résultats).
    \item 1. Passif \zh{被} : \zh{被同学们选上了} (a été élue par les camarades). 2. Verbe intransitif \zh{当选} : \zh{高票当选} (élue à une large majorité).
    \item « La démocratie n'est pas seulement un droit, c'est encore plus un devoir. » Structure \zh{不但/不仅 A，更/而且 B} (non seulement A, mais encore plus B).
    \item \textit{Piste :} Voter n'est pas facultatif ; il faut s'informer, participer, accepter le résultat, contrôler l'élu. Ex: nettoyer la classe si l'ambassadeur l'organise (\zh{义务}).
\end{enumerate}
\end{corrige}

\courssub{Collocations et structures grammaticales clés}
\begin{propriete}[Collocations électorales (Blocs lexicaux figés à mémoriser par cœur)]
Ces associations mots+mots sont la clé de la fluidité (HSK 4+).
\begin{center}
\begin{tabularx}{\linewidth}{|X|X|X|}
\hline
\rowcolor{popblueL}
\textbf{Collocation (Caractères + Pinyin)} & \textbf{Structure grammaticale} & \textbf{Traduction} \\
\hline
\zh{行使权利} (xíngshǐ quánlì) & V + O & Exercer son droit \\
\zh{履行义务} (lǚxíng yìwù) & V + O & S'acquitter de son devoir \\
\zh{竞选演讲} (jìngxuǎn yǎnjiǎng) & N + N (modif.) & Discours de campagne \\
\zh{竞选纲领} (jìngxuǎn gānglǐng) & N + N & Programme de campagne \\
\zh{竞选承诺} (jìngxuǎn chéngnuò) & N + N & Promesse de campagne \\
\zh{投出选票} (tóu chū xuǎnpiào) & V + Comp. dir. \zh{出} + O & Déposer son bulletin (vers l'extérieur) \\
\zh{把选票投进票箱} (bǎ xuǎnpiào tóu jìn piàoxiāng) & \zh{把} O V + Comp. dir. \zh{进} & Mettre le bulletin dans l'urne \\
\zh{公开点票} (gōngkāi diǎnpiào) & Adv + VO & Dépouiller publiquement \\
\zh{宣布结果} (xuānbù jiéguǒ) & V + O & Annoncer les résultats \\
\zh{高票当选} (gāopiào dāngxuǎn) & Adj + V (Chengyu) & Être élu(e) à une large majorité \\
\zh{被选上} (bèi xuǎn shàng) & Passif \zh{被} + V + Rés. & Être élu (focus sur le résultat subi) \\
\hline
\end{tabularx}
\end{center}
\end{propriete}

\begin{propriete}[Structures grammaticales réinvesties (Niveau Terminale A)]
\begin{itemize}
    \item \textbf{1. Disposition \cle{把} (Objet connu, action aboutie) :} Sujet + \textbf{把} + Objet + Verbe + Complément de résultat/lieu.
    \zh{请把选票放进票箱。} (Qǐng bǎ xuǎnpiào fàng jìn piàoxiāng.) — Mettez le bulletin dans l'urne.
    \zh{他把承诺写在纲领里。} (Tā bǎ chéngnuò xiě zài gānglǐng lǐ.) — Il a écrit sa promesse dans le programme.
    \item \textbf{2. Passif \cle{被} (Subir, résultat, souvent sans agent) :} Sujet (patient) + \textbf{被} + (Agent) + Verbe + \zh{了} / Complément.
    \zh{马莉被同学们选上了。} (Mǎ Lì bèi tóngxuémen xuǎn shàng le.) — Marie a été élue par les camarades.
    \zh{选票被点完了。} (Xuǎnpiào bèi diǎn wán le.) — Les bulletins ont été tous comptés.
    \item \textbf{3. Modaux d'engagement et de capacité (Discours) :}
    \begin{tabular}{lll}
    \zh{会} (huì) & capacité acquise, savoir-faire & \zh{我会组织活动。} (Je sais organiser / Je saurai organiser.) \\
    \zh{能} (néng) & capacité physique, permission, possibilité & \zh{我能代表大家。} (Je peux représenter tout le monde.) \\
    \zh{要} (yào) & volonté forte, intention, obligation morale & \zh{我要为同学服务。} (Je veux / Je dois servir les camarades.) \\
    \zh{想} (xiǎng) & souhait, désir & \zh{我想改变食堂菜单。} (Je veux changer le menu de la cantine.)
    \end{tabular}
    \item \textbf{4. Préposition \cle{为} (wèi, pour/bénéficiaire) :} \zh{为} + Bénéficiaire + Verbe.
    \zh{我为班级做贡献。} (Wǒ wèi bānjí zuò gòngxiàn.) — Je contribue pour la classe.
    \zh{我们为人民服务。} (Wǒmen wèi rénmín fúwù.) — Nous servons le peuple (devise classique).
    \item \textbf{5. Structure de progression \cle{不但/不仅...而且/更...} :}
    \zh{民主不仅是权利，更是义务。} (Mínzhǔ bùjǐn shì quánlì, gèng shì yìwù.)
\end{itemize}
\end{propriete}

\begin{attention}[Pièges fréquents pour francophones — Check-list de révision]
\begin{itemize}
    \item \textbf{Ordre Temps / Lieu / Verbe :} \textbf{Jamais} \zh{我们举行了选举在教室}. \textbf{Toujours} \zh{我们在教室举行了选举}. (Lieu \zh{在}... précède le verbe).
    \item \textbf{\cle{被} vs \cle{把} :} \cle{把} = Je fais quelque chose \textbf{à} l'objet (actif, manipulation). \cle{被} = L'objet \textbf{subit} l'action (passif, focus résultat).
    \begin{itemize}
        \item \zh{我把选票投了。} (J'ai voté / J'ai mis le bulletin — focus action).
        \item \zh{选票被投了。} (Le bulletin a été mis — focus objet/résultat).
    \end{itemize}
    \item \textbf{\cle{当选} (intransitif) vs \cle{选上} (verbe+résultat) :}
    \begin{itemize}
        \item Sujet = L'élu : \zh{他当选了} (Il a été élu) / \zh{他被选上了} (Il a été élu par qqn).
        \item Sujet = L'électeur : \zh{我们选上了他} (Nous l'avons élu). \textbf{Pas} \zh{我们当选了他}.
    \end{itemize}
    \item \textbf{Classificateurs (obligatoires après nombre/démonstratif) :}
    \zh{一张选票} (plat), \zh{一位候选人} (personne polie), \zh{一份纲领} (document), \zh{三票} (\zh{票} comme classifieur propre aux voix).
    \item \textbf{Tons critiques :}
    \zh{选} (xuǎn, 3e ton \textit{bas-remontant}) $\neq$ \zh{宣} (xuān, 1er ton \textit{haut} dans \zh{宣布}).
    \zh{票} (piào, 4e ton \textit{chutant}) $\neq$ \zh{漂} (piāo, 1er ton).
    \zh{仿} (fǎng, 3e ton) vs \zh{放} (fàng, 4e ton — dans \zh{放进}).
    \item \textbf{Ponctuation :} Pas d'espace après la virgule chinoise \zh{，} ni le point \zh{。}. Guillemets \zh{「 」} ou \zh{《 》}.
\end{itemize}
\end{attention}

\courssub{Collocations et production guidée : Tâche finale — Simulation électorale}
\begin{methode}[Organisation de la simulation « Élection de l'Ambassadeur de classe » (2h)]
Cette tâche finale mobilise \textbf{toutes les compétences} : compréhension (consignes), expression écrite (affiche, bulletin, compte rendu), expression orale (discours, annonce), interaction (campagne, vote), culture (procédure démocratique).
\begin{enumerate}
    \item \textbf{Préparation (10 min) :} Distribution du matériel (cartes A5 pour affiches, bulletins modèles, boîte « urne », liste électorale). Désignation de 2 \zh{监票人} (scrutateurs) et 1 \zh{记录员} (secrétaire).
    \item \textbf{Étape 1 — Candidature \& Affiche (Écrit, 15 min) :} Voir consignes détaillées ci-dessous. Affichage au mur.
    \item \textbf{Étape 2 — Campagne \& Discours (Oral, 20 min prep + 1 min/élève) :} Lecture des affiches (lecture silencieuse), puis discours debout.
    \item \textbf{Étape 3 — Vote secret (Pratique, 10 min) :} Remplissage bulletin (\zh{候选人姓名：\_\_\_\_\_\_} \zh{同意/不同意}), pliage, dépôt dans \zh{票箱} en disant \zh{我投票} ou \zh{我投给[X]一票}.
    \item \textbf{Étape 4 — Dépouillement \& Annonce (Oral collectif, 10 min) :} Scrutateurs lisent : \zh{这张票投给[X]} ; Secrétaire note au tableau ; Professeur annonce : \zh{[X]被选上了！} / \zh{[X]当选了！} \zh{恭喜[X]同学！}
    \item \textbf{Étape 5 — Compte rendu écrit (Écrit individuel, 15 min, noté /10) :} Rédaction pour le journal de l'établissement (voir modèle et grille).
\end{enumerate}
\end{methode}

\begin{textebox}[Document authentique : Affiche officielle de l'élection (Modèle professeur)]
\begin{center}
\textbf{雅温得双语高中 高三A班 班级大使选举公告} \\
\vspace{0.2cm}
\textit{Yǎwēndé Shuāngyǔ Gāozhōng Gāosān A bān Bānjí Dàshǐ Xuǎnjǔ Gōnggào} \\
\vspace{0.3cm}
\begin{tabular}{@{}l@{}}
选举日期：2025年10月15日 \\
投票时间：上午 10:00 -- 11:00 \\
地点：教室 305 \\
主持人：王老师 \\
尊敬的选民同学们：请行使你们的民主权利，投出宝贵的一票！
\end{tabular}
\end{center}
\textit{Traduction : Avis d'élection de l'Ambassadeur de la classe de Terminale A, Lycée Bilingue de Yaoundé. Date : 15 octobre 2025. Heure : 10h-11h. Lieu : Salle 305. Président : M. Wang. Chers électeurs, exercez votre droit démocratique, déposez votre précieux bulletin !}
\end{textebox}

\courssub{Consignes détaillées du Portfolio Élève (Noté /20)}
\begin{enumerate}
    \item \textbf{Affiche de campagne (A5, manuscrit, 3 phrases exactes) — /5 pts}
    \begin{itemize}
        \item Phrase 1 : Identité + candidature. \zh{我叫[Nom], 我是候选人。}
        \item Phrase 2 : Capacité (\zh{会}/\zh{能}) + Action pour la classe. \zh{我会组织[活动]。}
        \item Phrase 3 : Engagement (\zh{为} + Bénéficiaire + Verbe). \zh{我为全班同学服务。}
    \end{itemize}
    \textit{Critères :} Caractères corrects, pinyin au-dessus si hésitation, ponctuation chinoise, 3 phrases seulement.

    \item \textbf{Discours oral (4 phrases, notes bristol autorisées) — /5 pts}
    \begin{enumerate}
        \item Salutations formelles + Identité + Poste. \zh{尊敬的王老师，亲爱的同学们，我是[X]，我竞选班级大使。}
        \item Qualité / Expérience. \zh{我有[经验]，很负责任。} / \zh{我当过[小组长]。}
        \item Programme (\zh{要}/\zh{会} + V + O). \zh{我要组织[足球赛和汉语角]。}
        \item Appel au vote + Remerciements. \zh{请大家投我一票，谢谢大家！}
    \end{enumerate}
    \textit{Critères :} Prononciation (tons), fluidité, vocabulaire cible, contact visuel.

    \item \textbf{Bulletin de vote \& Participation orale au vote — /3 pts}
    Remplissage correct en caractères (\zh{候选人姓名：马莉}, cocher \zh{同意}). Phrase orale : \zh{我投给马莉一票。}

    \item \textbf{Compte rendu écrit pour le journal « Vie lycéenne » (4-5 phrases, 6 mots cible surlignés) — /7 pts}
    \textbf{Trame imposée :}
    \begin{itemize}
        \item P1 (Contexte) : \zh{时间 + 地点 + 主语 + 举行了 + 选举。}
        \item P2 (Déroulement) : \zh{几位候选人 + 作了 + 竞选演讲。}
        \item P3 (Vote) : \zh{同学们 + 投出了 + 选票 / 行使了权利。}
        \item P4 (Résultat) : \zh{…… + 被选上了 / 高票当选了。}
        \item P5 (Réflexion civique - optionnelle) : \zh{我们 + 体会到 + 民主的意义。}
    \end{itemize}
\end{enumerate}

\begin{center}
\begin{popfigure}
\begin{tikzpicture}[
    node distance=1.2cm and 1.5cm,
    box/.style={rectangle, draw=popblue, thick, fill=popblueL, rounded corners, minimum width=2.8cm, minimum height=1cm, align=center, font=\small},
    arrow/.style={-Stealth, thick, popdark}
]
\node[box, align=center] (e1){Étape 1\\Candidature \& Affiche};
\node[box, right=of e1, align=center] (e2){Étape 2\\Discours de campagne};
\node[box, right=of e2, align=center] (e3){Étape 3\\Vote secret};
\node[box, below=of e2, align=center] (e4){Étape 4\\Dépouillement \& Résultats};
\node[box, below=of e3, align=center] (e5){Étape 5\\Compte rendu écrit};

\draw[arrow] (e1) -- (e2);
\draw[arrow] (e2) -- (e3);
\draw[arrow] (e3) -- (e4);
\draw[arrow] (e4) -- (e5);
\draw[arrow] (e2) -- (e4); % lien visuel campagne -> résultat
\end{tikzpicture}
\legende{Flux de la simulation électorale (Tâche finale)}
\end{popfigure}
\end{center}

\courssub{Proposition de production et corrigé commenté (Modèle « Excellence »)}
\begin{exoresolu}[Portfolio de l'élève « Marie » (\zh{马莉} Mǎ Lì), candidate élue]
\textbf{1. Affiche de campagne (Manuscrite A5)} \\
\zh{我叫马莉，我是候选人。} (Wǒ jiào Mǎ Lì, wǒ shì hòuxuǎnrén.) \\
\zh{我会组织汉语角和足球赛。} (Wǒ huì zǔzhī Hànyǔ jiǎo hé zúqiú sài.) \\
\zh{我为全班同学服务！} (Wǒ wèi quánbān tóngxué fúwù!) \\
\textit{Commentaire :} Respect strict du cahier des charges (3 phrases). Phrase 2 : \zh{会} (capacité) + \zh{组织} + objets coordonnés par \zh{和}. Phrase 3 : Structure \cle{为} + \zh{全班同学} (bénéficiaire) + \zh{服务}. Ponctuation chinoise \zh{。} \zh{！}.

\textbf{2. Discours oral (Lu sur bristol, ~45 sec)} \\
\zh{尊敬的王老师，亲爱的同学们，我是马莉，我竞选班级大使。} (Zūnjìng de Wáng lǎoshī, qīn'ài de tóngxuémen, wǒ shì Mǎ Lì, wǒ jìngxuǎn bānjí dàshǐ.) \\
\zh{我当过小组长，有经验，很负责任。} (Wǒ dāngguò xiǎozǔzhǎng, yǒu jīngyàn, hěn fùzérèn.) \\
\zh{我要组织中非文化周，也要改善食堂饭菜。} (Wǒ yào zǔzhī Zhōng-Fēi wénhuà zhōu, yě yào gǎishàn shítáng fàncài.) \\
\zh{请大家投我一票，谢谢大家！} (Qǐng dàjiā tóu wǒ yī piào, xièxie dàjiā!) \\
\textit{Commentaire :} Registre formel (\zh{尊敬的/亲爱的}). Argumentaire : \zh{当过} (expérience passée) + \zh{有经验} + \zh{负责任}. Programme concret ancré au Cameroun (\zh{中非文化周}, \zh{食堂饭菜}). \zh{也要} (aussi). Appel au vote direct \zh{投我一票}. Tons soignés : \zh{竞选} (jìngxuǎn), \zh{负责} (fùzé), \zh{改善} (gǎishàn).

\textbf{3. Annonce des résultats (Par le professeur / scrutateur)} \\
\zh{同学们，现在宣布结果。马莉得了十五票，张华得了八票。马莉被选上了！恭喜马莉同学当选班级大使！} (Tóngxuémen, xiànzài xuānbù jiéguǒ. Mǎ Lì dé le shíwǔ piào, Zhāng Huá dé le bā piào. Mǎ Lì bèi xuǎn shàng le! Gōngxǐ Mǎ Lì tóngxué dāngxuǎn bānjí dàshǐ!) \\
\textit{Commentaire :} \zh{得了} (dé le, obtenir des voix). Passif \zh{被选上了} (focus sur l'élue). Verbe \zh{当选} en actif final (nouveau statut). Vocabulaire précis : \zh{宣布结果}.

\textbf{4. Compte rendu écrit pour le journal (Version finale corrigée, mots cible \textbf{surlignés})} \\
\zh{十月十五日，我们班在教室三零五举行了\textbf{班级大使民主选举}。} (Shíyuè shíwǔ rì, wǒmen bān zài jiàoshì sān líng wǔ jǔxíng le \textbf{bānjí dàshǐ mínzhǔ xuǎnjǔ}.) \\
\zh{三位\textbf{候选人}作了精彩的\textbf{竞选演讲}，介绍了各自的\textbf{竞选纲领}。} (Sān wèi \textbf{hòuxuǎnrén} zuò le jīngcǎi de \textbf{jìngxuǎn yǎnjiǎng}, jièshào le gèzì de \textbf{jìngxuǎn gānglǐng}.) \\
\zh{全体同学\textbf{行使了民主权利}，投出了宝贵的\textbf{选票}。} (Quántǐ tóngxué \textbf{xíngshǐ le mínzhǔ quánlì}, tóu chū le bǎoguì de \textbf{xuǎnpiào}.) \\
\zh{经过公开\textbf{点票}，马莉同学\textbf{高票当选}了。} (Jīngguò gōngkāi \textbf{diǎnpiào}, Mǎ Lì tóngxué \textbf{gāopiào dāngxuǎn} le.) \\
\zh{我们深刻\textbf{体会到}：民主不仅是\textbf{权利}，更是\textbf{义务}。} (Wǒmen shēnkè \textbf{tǐhuì dào}: mínzhǔ bùjǐn shì \textbf{quánlì}, gèng shì \textbf{yìwù}.) \\

\textit{Commentaire détaillé (Grille d'évaluation /7 pts) :} \\
\begin{itemize}
    \item \textbf{Phrase 1 (1,5 pt) :} Ordre canonique \textbf{Temps} (\zh{十月十五日}) — \textbf{Lieu} (\zh{在教室三零五}) — Sujet — Verbe \zh{举行了} — Objet \zh{民主选举} (collocation). \zh{三零五} (lecture chiffre par chiffre pour n° salle). \textit{Mots cible : 选举.}
    \item \textbf{Phrase 2 (1,5 pt) :} \zh{三位} (classifieur \cle{位}). \zh{各自的} (chacun son). \zh{竞选演讲}, \zh{竞选纲领} (collocations cibles). \textit{Mots cible : 候选人, 竞选演讲, 竞选纲领.}
    \item \textbf{Phrase 3 (1,5 pt) :} \zh{行使权利} (collocation). \zh{投出} (verbe directionnel \cle{出}). \zh{宝贵的} (adj + \cle{的}). \textit{Mots cible : 权利, 选票.}
    \item \textbf{Phrase 4 (1,5 pt) :} \zh{经过} (préposition « après »). \zh{公开点票}. \zh{高票当选} (chengyu / collocation forte). \textit{Mots cible : 点票, 当选 (via 高票当选).}
    \item \textbf{Phrase 5 (1 pt bonus) :} Réflexion civique. Structure \zh{不但/不仅……更/而且……}. \zh{深刻体会到}. \textit{Mots cible : 体会, 权利, 义务.}
    \item \textbf{Total mots cible utilisés : 10/6 requis. Qualité de l'écriture : Caractères lisibles, proportions respectées, pinyin absent (niveau Terminale).}
\end{itemize}
\end{exoresolu}

\courssub{Bilan de la leçon}
\begin{aretenir}[Bilan linguistique, citoyen et évaluatif]
Cette leçon 21 a permis d'atteindre les objectifs du Module 3 « Citoyenneté et ouverture au monde » pour la Terminale A :
\begin{itemize}
    \item \textbf{Lexique (HSK 3-4) :} Ancrage durable des 18 mots-clés + 12 collocations via l'usage répétitif multisensoriel (visuel : affiche/tableau ; auditif : discours/dépouillement ; kinesthésique : écriture bulletin/compte rendu ; graphique : caractères). Les blocs \zh{行使权利}, \zh{投出选票}, \zh{竞选纲领}, \zh{高票当选} sont désormais des unités de récupération automatique.
    \item \textbf{Grammaire (Niveau Terminale) :} Application contextualisée et contrastée :
    \begin{itemize}
        \item \cle{把} (disposition active sur objet précis : \zh{把选票投进票箱}).
        \item \cle{被} (passif résultat/subie : \zh{被选上了}).
        \item Modaux \cle{要}/\cle{会}/\cle{能}/\cle{想} (nuances engagement/capacité).
        \item \cle{为} (bénéficiaire).
        \item Ordre T-L-V respecté spontanément.
    \end{itemize}
    \item \textbf{Écriture :} Renforcement motricité fine caractères à haute fréquence (\zh{选}, \zh{票}, \zh{举}, \zh{行}, \zh{权}, \zh{义}, \zh{民}, \zh{主}) via rédaction authentique (affiche, bulletin, compte rendu).
    \item \textbf{Oral (Épreuve E3L) :} Prise de parole en public (discours 1 min), interaction ritualisée (vote), prosodie (tons, liaison \zh{投给} \textit{tóu gěi}), registre formel (\zh{尊敬的}).
    \item \textbf{Culture / Citoyenneté (Socio-culturel) :} Vécu d'une procédure démocratique complète en L2. Comparaison terminologique : \zh{班级大使} / \zh{班长} (Chine) vs « Délégué » / « Représentant » (Cameroun). Réflexion sur le couple \zh{权利}/\zh{义务} (Droits/Dévoirs) dans les deux constitutions. Ouverture sur \zh{全过程人民民主} (Démocratie populaire à processus intégral, concept chinois actuel).
\end{itemize}
\textbf{Prolongements possibles (Différenciation / Approfondissement) :}
\begin{enumerate}
    \item \textbf{Écriture formelle :} Rédiger la \zh{感谢信} (lettre de remerciement) de l'élu aux électeurs (structure \zh{感谢……对……的支持}).
    \item \textbf{Débat oral :} « \zh{大使应该做什么？} » (Que doit faire l'ambassadeur ?) — Réinvestir \zh{应该} + V, \zh{必须} + V, \zh{计划} + V.
    \item \textbf{Projet interdisciplinaire (EMC/Histoire-Géo) :} Comparer le système électoral camerounais (Élections à un tour, majorité relative) et le système chinois (Élections indirectes au niveau national, suffrage direct à la base) via un tableau bilingue.
\end{enumerate}
\end{aretenir}

\begin{savaistu}[Regard croisé : Vocabulaire institutionnel Cameroun / Chine]
\begin{center}
\begin{tabularx}{\linewidth}{|X|X|X|}
\hline
\rowcolor{popblueL}
\textbf{Concept} & \textbf{Cameroun (Français / Chinois)} & \textbf{Chine (Chinois / Français)} \\
\hline
Délégué de classe & Délégué / \zh{班长} (bānzhǎng) ou \zh{班级大使} (bānjí dàshǐ) & \zh{班长} (bānzhǎng) — Élu par la classe, responsable discipline/études \\
Représentant élu & Député / \zh{议员} (yìyuán) & \zh{人大代表} (réndà dàibiǎo) — Député à l'Assemblée Populaire Nationale \\
Scrutin & Vote / \zh{投票} (tóupiào) & \zh{选举} (xuǎnjǔ) — Terme générique ; \zh{差额选举} (élection compétitive) \\
Programme & Programme / \zh{纲领} (gānglǐng) & \zh{竞选纲领} (jìngxuǎn gānglǐng) / \zh{施政纲领} (programme de gouvernance) \\
\hline
\end{tabularx}
\end{center}
\textit{Note :} Au Cameroun, le terme \zh{大使} (ambassadeur) est valorisant pour un rôle de liaison culturelle. En Chine, \zh{班长} a une autorité administrative forte (notation, gestion). Le terme \zh{民主} (mínzhǔ) existe dans les deux constitutions mais les réalités procédurales diffèrent (suffrage universel direct vs démocratie consultative/délibérative).
\end{savaistu}

\newpage

\coursec{Méthodes et exercices résolus}

\coursec{Leçon 21 — Vocabulaire et compréhension de texte : démocratie}

\courssub{Mise en route et découverte du thème}

\begin{experience}[Brainstorming et activation du schéma notionnel]
\textbf{Consigne :} En binôme, associez chaque mot-clé français à son équivalent chinois (caractères + pinyin) affiché au tableau. Puis, classez-les en deux colonnes : \textbf{Droits} / \textbf{Devoirs \& Institutions}.
\begin{center}
\begin{tabularx}{\linewidth}{|X|X|X|}
\hline
\rowcolor{popblueL} \textbf{Français} & \textbf{Chinois (Caractères + Pinyin)} & \textbf{Catégorie} \\
\hline
Citoyen & \zh{公民 gōngmín} & \textit{Acteur} \\
\hline
Droit & \zh{权利 quánlì} & Droits \\
\hline
Devoir & \zh{义务 yìwù} & Devoirs \\
\hline
Élection & \zh{选举 xuǎnjǔ} & Institutions \\
\hline
Vote / Voter & \zh{投票 tóupiào} & Droits / Action \\
\hline
Loi & \zh{法律 fǎlǜ} & Institutions \\
\hline
Représentant & \zh{代表 dàibiǎo} & Institutions \\
\hline
Candidat & \zh{候选人 hòuxuǎnrén} & Institutions \\
\hline
Équitable / Justice & \zh{公平 gōngpíng / 公正 gōngzhèng} & Valeurs \\
\hline
Harmonie & \zh{和谐 xiéméi} & Valeurs \\
\hline
\end{tabularx}
\end{center}
\textbf{Prolongement oral :} \zh{你认为“民主”最重要的特征是什么？} (À votre avis, quelle est la caractéristique la plus importante de la « démocratie » ?) — Réponse attendue : \zh{权力属于人民} / \zh{选举公平} / \zh{法律面前人人平等}.
\end{experience}

\courssub{Texte support : lecture guidée}

\begin{textebox}[Texte support : 什么是民主？]
\zh{民主是一种政治制度。} (Mínzhǔ shì yī zhǒng zhèngzhì zhìdù.) \\
\zh{首先，它意味着权力属于人民。} (Shǒuxiān, tā yìwèizhe quánlì shǔyú rénmín.) \\
\zh{公民通过选举选择代表。} (Gōngmín tōngguò xuǎnjǔ xuǎnzé dàibiǎo.) \\
\zh{选举是依法进行的。} (Xuǎnjǔ shì yīfǎ jìnxíng de.) \\
\zh{其次，选举必须公平、公正、公开。} (Qícì, xuǎnjǔ bìxū gōngpíng, gōngzhèng, gōngkāi.) \\
\zh{每一位公民都有投票的权利和被选举的权利。} (Měi yī wèi gōngmín dōu yǒu tóupiào de quánlì hé bèi xuǎnjǔ de quánlì.) \\
\zh{此外，法律面前人人平等。} (Cǐwài, fǎlǜ miànqián rénrén píngděng.) \\
\zh{政府必须依法办事，保障公民的自由。} (Zhèngfǔ bìxū yīfǎ bànshì, bǎozhàng gōngmín de zìyóu.) \\
\zh{总之，民主不仅是一种投票的方式，更是一种尊重权利、促进社会和谐发展的理念。} (Zǒngzhī, mínzhǔ bùjǐn shì yī zhǒng tóupiào de fāngshì, gèng shì yī zhǒng zūnzhòng quánlì, cùjìn shèhuì xiéméi fāzhǎn de lǐniàn.) \\
\zh{在喀麦隆，如同在中国，人民都在为建设更美好的民主社会而努力。} (Zài Kāmàilóng, rútóng zài Zhōngguó, rénmín dōu zài wèi jiànshè gèng měihǎo de mínzhǔ shèhuì ér nǔlì.)

\textbf{Traduction :} La démocratie est un système politique. Premièrement, cela signifie que le pouvoir appartient au peuple. Les citoyens choisissent leurs représentants par le biais d'élections. Les élections se déroulent conformément à la loi. Deuxièmement, les élections doivent être équitables, justes et transparentes. Chaque citoyen a le droit de voter et le droit d'être élu. De plus, tous sont égaux devant la loi. Le gouvernement doit agir selon la loi et garantir la liberté des citoyens. En résumé, la démocratie n'est pas seulement une manière de voter, c'est davantage un idéal qui respecte les droits et favorise le développement harmonieux de la société. Au Cameroun, comme en Chine, les peuples s'efforcent tous de construire une société démocratique plus belle.
\end{textebox}

\begin{methode}[Stratégie de lecture guidée : du global au détaillé]
\textbf{Objectif :} Lire et comprendre un court texte chinois sur la démocratie (niveau HSK 3-4).
\begin{enumerate}
    \item \textbf{Observation paratextuelle :} Lire le titre (\zh{什么是民主？}), repérer la structure (paragraphes courts, énumération), identifier le registre (exposé informatif, langue standard).
    \item \textbf{Lecture globale (skimming) :} Lire une première fois sans s'arrêter. Repérer les \cle{mots-clés} récurrents (\zh{民主}、\zh{选举}、\zh{公民}、\zh{权利}、\zh{法律}). Saisir l'idée générale : définition $\rightarrow$ mécanismes (élection, loi) $\rightarrow$ droits/devoirs $\rightarrow$ synthèse valeur + exemples Cameroun/Chine.
    \item \textbf{Lecture détaillée (scanning) :} Relire phrase par phrase. Souligner les \cle{connecteurs logiques} (\zh{首先}、\zh{其次}、\zh{此外}、\zh{总之}). Encadrer les structures grammaticales cibles : \zh{是……的} (phrase 4, insistance sur la manière \zh{依法}), \zh{不仅……而且……} (phrase 9, progression), \zh{被} (phrase 6, passif \zh{被选举}), \zh{为……而……} (phrase 10, finalité).
    \item \textbf{Gestion de l'inconnu :} Pour un caractère inconnu, chercher le \cle{radical} (indice de sens) et le \cle{phonétique} (indice de son). Ex: \zh{候选人} $\rightarrow$ \zh{候} (attendre, radical \zh{亻}) + \zh{选} (choisir, radical \zh{讠}) + \zh{人} (personne) = candidat.
    \item \textbf{Synthèse orale :} Résumer le texte en 3 phrases en chinois à l'oral (Pinyin + tons corrects + sandhi appliqué).
\end{enumerate}
\end{methode}

\begin{exoresolu}[Application : Lecture et analyse du texte support]
\textbf{Énoncé :} Lisez le texte ci-dessus. 1) Donnez le thème principal en une phrase française. 2) Relevez 4 connecteurs logiques, donnez leur pinyin et traduisez-les. 3) Expliquez la structure \zh{是……的} dans la phrase 4. 4) Traduisez les phrases 9 et 10 en français.

\tcblower
\textbf{Solution détaillée :}
\begin{enumerate}
    \item \textbf{Thème :} Le texte définit la démocratie comme un système politique fondé sur la souveraineté populaire, des élections libres/équitables/transparentes organisées selon la loi, l'égalité devant la loi, la protection des libertés, et présente cela comme un idéal de développement harmonieux partagé par le Cameroun et la Chine.
    \item \textbf{Connecteurs logiques :}
    \begin{itemize}
        \item \zh{首先} (shǒuxiān) -- « Premièrement / Tout d'abord » (Marque le premier argument majeur).
        \item \zh{其次} (qícì) -- « Deuxièmement / Ensuite » (Marque l'argument suivant).
        \item \zh{此外} (cǐwài) -- « De plus / En outre » (Ajoute une dimension juridique/institutionnelle).
        \item \zh{总之} (zǒngzhī) -- « En résumé / En conclusion / Bref » (Introduit une synthèse valorisante).
    \end{itemize}
    \item \textbf{Structure \zh{是……的} (Phrase 4 : \zh{选举是依法进行的}) :} Cette structure met l'accent sur un \cle{détail de circonstance} (manière, lieu, temps, auteur) d'une action \textbf{accomplie / passée / factuelle}.
    \textbf{Schéma :} Sujet (\zh{选举}) + \zh{是} + Complément de circonstance (\zh{依法} « selon la loi / légalement ») + Verbe (\zh{进行} « se dérouler ») + \zh{的}.
    \textbf{Traduction ciblée :} « Les élections \textbf{sont menées} conformément à la loi » (insistance sur la légalité de la procédure).
    \item \textbf{Traduction phrases 9-10 :} « En résumé, la démocratie n'est pas seulement une manière de voter, c'est davantage un idéal qui respecte les droits et favorise le développement harmonieux de la société. Au Cameroun, comme en Chine, les peuples s'efforcent tous de construire une société démocratique plus belle. »
\end{enumerate}
\textbf{Conclusion :} La lecture active (repérage connecteurs + structures cibles) permet de dépasser le déchiffrage pour accéder à la logique argumentative du texte.
\end{exoresolu}

\courssub{Glossaire du thème}

\begin{center}
\begin{tabularx}{\linewidth}{|X|X|X|X|}
\hline
\rowcolor{popblueL} \textbf{Caractères} & \textbf{Pinyin} & \textbf{Nature} & \textbf{Traduction} \\
\hline
民主 & mínzhǔ & n. / adj. & démocratie / démocratique \\
\hline
政治制度 & zhèngzhì zhìdù & n. & système politique \\
\hline
权力 & quánlì & n. & pouvoir (politique, autorité de l'État) \\
\hline
人民 & rénmín & n. & le peuple \\
\hline
公民 & gōngmín & n. & citoyen \\
\hline
选举 & xuǎnjǔ & v. / n. & élire / élection \\
\hline
代表 & dàibiǎo & n. / v. & représentant / représenter \\
\hline
候选人 & hòuxuǎnrén & n. & candidat (à une élection) \\
\hline
投票 & tóupiào & v. / n. & voter / vote \\
\hline
权利 & quánlì & n. & droit (juridique, avantage citoyen) \\
\hline
义务 & yìwù & n. & devoir, obligation \\
\hline
公平 & gōngpíng & adj. & équitable, impartial \\
\hline
公正 & gōngzhèng & adj. & juste, impartial (procédure, jugement) \\
\hline
公开 & gōngkāi & adj. / v. & public, transparent / rendre public \\
\hline
法律 & fǎlǜ & n. & loi, législation \\
\hline
依法 & yīfǎ & adv. & conformément à la loi, légalement \\
\hline
平等 & píngděng & adj. / n. & égal / égalité \\
\hline
政府 & zhèngfǔ & n. & gouvernement \\
\hline
办事 & bànshì & v. & gérer des affaires, agir \\
\hline
保障 & bǎozhàng & v. & garantir, assurer \\
\hline
自由 & zìyóu & n. / adj. & liberté / libre \\
\hline
方式 & fāngshì & n. & manière, mode, façon \\
\hline
尊重 & zūnzhòng & v. & respecter \\
\hline
促进 & cùjìn & v. & promouvoir, favoriser \\
\hline
社会 & shèhuì & n. & société \\
\hline
和谐 & xiéméi & adj. & harmonieux \\
\hline
发展 & fāzhǎn & v. / n. & développer / développement \\
\hline
理念 & lǐniàn & n. & idéal, concept, philosophie \\
\hline
建设 & jiànshè & v. / n. & construire / construction \\
\hline
努力 & nǔlì & v. / adv. & faire des efforts / avec effort \\
\hline
喀麦隆 & Kāmàilóng & n.p. & Cameroun \\
\hline
如同 & rútóng & conj. & tout comme, de même que \\
\hline
\end{tabularx}
\end{center}

\courssub{Clés (radicaux) et ordre des traits}

\begin{propriete}[Analyse morphologique : Radicaux sémantiques et ordre des traits]
\textbf{Principaux radicaux du thème :}
\begin{itemize}
    \item \zhbf{亻} (rén, « personne ») : \zh{公民} (\zh{公} + \zh{民}), \zh{候选人} (\zh{候}, \zh{人}), \zh{代表} (\zh{代} \rightarrow \zh{亻} dans \zh{伟}? Non, \zh{代} radical \zh{亻}, \zh{表} radical \zh{衣}). \textit{Note :} \zh{民} radical \zh{氏} (clan), \zh{主} radical \zh{丶} / \zh{王}.
    \item \zhbf{讠} (yán, « parole, langage ») : \zh{选} (choisir), \zh{举} (radical \zh{廾} gǒng « mains jointes» mais \zh{选举} souvent groupé par \zh{选}), \zh{表决} (bǐaojué, voter/décider), \zh{竞选} (jìngxuǎn, briguer).
    \item \zhbf{力} (lì, « force, pouvoir ») : \zh{权} (quán, pouvoir/droit -- haut \zh{木} bois + bas \zh{force}), \zh{义} (yì, devoir/justice -- radical \zh{礻} shì « autel/rite » + \zh{我}), \zh{务} (wù, affaire/devoir -- radical \zh{力}).
    \item \zhbf{水} (shuǐ, « eau ») : \zh{法} (fǎ, loi -- \zh{氵} + \zh{去}), \zh{治} (zhì, gouverner -- \zh{氵} + \zh{台}).
    \item \zhbf{心} (xīn, « cœur ») : \zh{义} (yì, contient \zh{心} en bas dans l'écriture traditionnelle \zh{義}, simplifié \zh{礻}+\zh{我}), \zh{忠} (zhōng, loyauté).
\end{itemize}

\textbf{Ordre des traits (règles clés) :}
\begin{enumerate}
    \item \zh{民} (5 traits) : Horizontal (一) $\rightarrow$ Vertical (丨) $\rightarrow$ Horizontal (一) $\rightarrow$ Horizontal (一) $\rightarrow$ Crochet vertical (乚). \textit{Règle :} Haut $\rightarrow$ Bas ; Horizontal avant vertical qui traverse.
    \item \zh{主} (5 traits) : Point (丶) $\rightarrow$ Horizontal (一) $\rightarrow$ Vertical (丨) $\rightarrow$ Horizontal (一) $\rightarrow$ Horizontal (一). \textit{Règle :} Le trait principal central (vertical) avant les horizontales qui le croisent ? Non, pour \zh{主} : Point, puis horizontale haute, puis verticale, puis deux horizontales basses.
    \item \zh{权} (6 traits) : \zh{木} (4 traits :横竖撇捺) + \zh{力} (2 traits :横折钩). \textit{Règle :} Composant du haut (\zh{木}) avant composant du bas (\zh{力}).
    \item \zh{选} (9 traits) : \zh{讠} (2 traits :点、提) + \zh{巽} (7 traits). \textit{Règle :} Gauche (\zh{讠}) avant droite.
    \item \zh{举} (9 traits) : \zh{兴} (6 traits) + \zh{廾} (3 traits). \textit{Règle :} Haut avant bas.
    \item \zh{法} (8 traits) : \zh{氵} (3 traits :点、点、提) + \zh{去} (5 traits). \textit{Règle :} Gauche avant droite.
\end{enumerate}
\textbf{Activité :} Sur cahier, écrire 5 fois chaque caractère ci-dessus en respectant l'ordre des traits et en prononçant la syllabe (ton compris).
\end{propriete}

\courssub{Compréhension du texte}

\begin{methode}[Répondre aux questions de compréhension : localiser, reformuler, justifier]
\textbf{Objectif :} Répondre par écrit (en chinois, caractères + pinyin) à des questions fermées et ouvertes.
\begin{enumerate}
    \item \textbf{Analyse de la question :} Identifier le \cle{mot-interrogatif} (\zh{谁} qui, \zh{什么} quoi, \zh{为什么} pourquoi, \zh{怎么} comment, \zh{哪里} où, \zh{哪些} lesquels). Repérer le \cle{mot-clé} du texte visé.
    \item \textbf{Localisation dans le texte :} Scanner le texte pour trouver la phrase contenant le mot-clé ou son synonyme. Surligner la zone de réponse.
    \item \textbf{Construction de la réponse (Chinois) :}
        \begin{itemize}
            \item \textbf{Question fermée (Oui/Non / Choix) :} Réponse courte + phrase complète. Ex: \zh{是的，民主意味着权力属于人民。}
            \item \textbf{Question ouverte (\zh{为什么} / \zh{怎么}) :} Reprendre le sujet de la question + structure \zh{因为……所以……} ou \zh{由于……因此……}. \cle{Reformuler} avec ses mots (synonymes, changement de voix active/passive) plutôt que copier-coller.
        \end{itemize}
    \item \textbf{Vérification grammaticale (Check-list) :} Ordre Sujet -- Temps -- Lieu -- Manière -- Verbe -- Complément. Classificateurs corrects (\zh{个/位/项/部}). \zh{的/得/地} distincts. Ponctuation chinoise (。 ， 、).
    \item \textbf{Auto-correction :} Relire : « Est-ce que cela répond \emph{exactement} à la question ? » « Les tons (et sandhi) sont-ils notés sur le pinyin ? »
\end{enumerate}
\end{methode}

\begin{exoresolu}[Application : Questions de compréhension écrites]
\textbf{Énoncé :} Répondez en chinois (caractères + pinyin) aux questions suivantes sur le texte \zh{什么是民主？}.

\begin{enumerate}
    \item \zh{根据课文，民主制度的核心含义是什么？} (D'après le texte, quel est le sens central du système démocratique ?)
    \item \zh{选举必须具备哪三个原则？} (Quels sont les trois principes que l'élection doit posséder ?)
    \item \zh{选举是怎么进行的？} (Comment les élections se déroulent-elles ? \textit{Cible : structure 是……的})
    \item \zh{为什么说“民主不仅是一种投票的方式”？请用课文信息解释。} (Pourquoi dit-on que la démocratie n'est pas seulement une façon de voter ? Expliquez avec le texte.)
    \item \zh{请用“不仅……而且……”造句，谈谈你对喀麦隆民主建设的看法。} (Utilisez « non seulement... mais aussi... » pour parler de la construction démocratique au Cameroun.)
\end{enumerate}

\tcblower
\textbf{Solution détaillée :}
\begin{enumerate}
    \item \textbf{Localisation :} Phrases 1-2 (\zh{民主是一种政治制度……权力属于人民}).
    \textbf{Réponse :} \zh{民主制度的核心含义是权力属于人民。} (Mínzhǔ zhìdù de héxīn hányì shì quánlì shǔyú rénmín.)
    \item \textbf{Localisation :} Phrase 5 (\zh{选举必须公平、公正、公开}).
    \textbf{Réponse :} \zh{选举必须具备公平、公正、公开三个原则。} (Xuǎnjǔ bìxū jùbèi gōngpíng, gōngzhèng, gōngkāi sān gè yuánzé.) \textbf{Note :} Classificateur \zh{个} pour \zh{原则}.
    \item \textbf{Localisation :} Phrase 4 (\zh{选举是依法进行的}).
    \textbf{Réponse :} \zh{选举是依法进行的。} (Xuǎnjǔ shì yīfǎ jìnxíng de.) \textbf{Analyse :} Structure \zh{是……的} insistant sur la manière (\zh{依法}).
    \item \textbf{Localisation :} Phrase 9 (\zh{民主不仅是一种投票的方式，更是一种尊重权利、促进社会和谐发展的理念}).
    \textbf{Réponse :} \zh{因为民主更是一种尊重权利、促进社会和谐发展的理念。} (Yīnwèi mínzhǔ gèng shì yī zhǒng zūnzhòng quánlì, cùjìn shèhuì xiéméi fāzhǎn de lǐniàn.) \textbf{Structure :} \zh{因为……} répond à \zh{为什么}.
    \item \textbf{Production personnelle (Modèle) :}
    \zh{喀麦隆的民主建设不仅需要完善的法律，而且需要每个公民的积极参与。} (Kāmàilóng de mínzhǔ jiànshè bùjǐn xūyào wánshàn de fǎlǜ, érqiě xūyào měi gè gōngmín de jījí cānyù.)
    \textbf{Analyse :} \zh{不仅} placé avant le verbe \zh{需要} (sujet \zh{民主建设} omis par topicalisation) + \zh{而且} + \zh{需要}... Parallélisme parfait (V + O).
\end{enumerate}
\textbf{Conclusion :} La réponse en chinois doit être une phrase complète, grammaticalement correcte, qui réutilise le vocabulaire du texte tout en montrant une capacité de reformulation syntaxique.
\end{exoresolu}

\courssub{Collocations et production guidée}

\begin{methode}[Mémorisation active du vocabulaire : radicaux, tons, collocations]
\textbf{Objectif :} Reconnaître, prononcer (pinyin, tons, sandhi) et employer le lexique de la démocratie.
\begin{enumerate}
    \item \textbf{Analyse morphologique (Radicaux) :} Grouper par radical pour créer des réseaux sémantiques (voir boîte « Clés » ci-dessus).
    \item \textbf{Entraînement prosodique (Tons \& Sandhi) :} Isoler les paires tonales difficiles. Répétition chorale puis individuelle.
        \begin{itemize}
            \item \zh{权利} (quánlì, 2-4) vs \zh{权力} (quánlì, 2-4) \emph{mêmes tons, sens diff., radicaux diff.}
            \item \zh{民主} (mínzhǔ, 2-3).
            \item \zh{选举} (xuǎn\textbf{3} jǔ\textbf{3} $\rightarrow$ sandhi xuán\textbf{2} jǔ\textbf{3}).
            \item \zh{一位} (yī\textbf{1} $\rightarrow$ yí\textbf{2} wèi).
            \item \zh{不参加} (bù\textbf{4} $\rightarrow$ bú\textbf{2} cānjiā).
        \end{itemize}
    \item \textbf{Collocations fixes (搭配) -- À mémoriser par cœur :}
        \begin{center}
        \begin{tabularx}{\linewidth}{|X|X|X|}
        \hline
        \rowcolor{popblueL} \textbf{Verbe} & \textbf{Objet (Nom)} & \textbf{Traduction} \\
        \hline
        \zh{行使} (xíngshǐ) & \zh{权利} (quánlì) & exercer un droit \\
        \hline
        \zh{履行} (lǚxíng) & \zh{义务} (yìwù) & s'acquitter d'un devoir \\
        \hline
        \zh{举行} (jǔxíng) / \zh{进行} (jìnxíng) & \zh{选举} (xuǎnjǔ) & organiser / mener une élection \\
        \hline
        \zh{投票给} (tóupiào gěi) / \zh{选} (xuǎn) & \zh{候选人} (hòuxuǎnrén) / \zh{代表} (dàibiǎo) & voter pour / élire \\
        \hline
        \zh{遵守} (zūnshǒu) / \zh{依法} (yīfǎ) & \zh{法律} (fǎlǜ) & respecter la loi / agir selon la loi \\
        \hline
        \zh{保障} (bǎozhàng) & \zh{自由} (zìyóu) / \zh{权利} (quánlì) & garantir la liberté / les droits \\
        \hline
        \zh{促进} (cùjìn) & \zh{发展} (fāzhǎn) / \zh{和谐} (xiéméi) & favoriser le développement / l'harmonie \\
        \hline
        \end{tabularx}
        \end{center}
    \item \textbf{Production guidée :} Écrire 3 phrases personnelles en utilisant 3 collocations différentes. Vérifier l'ordre : Sujet + Temps + Lieu + Manière + Verbe + Complément.
\end{enumerate}
\end{methode}

\begin{exoresolu}[Application : Vocabulaire, radicaux, collocations et phrases]
\textbf{Énoncé :} 1) Complétez le tableau pour 5 mots. 2) Choisissez le bon mot entre parenthèses. 3) Écrivez une phrase avec \zh{行使权利} et une avec \zh{履行义务}.

\begin{center}
\begin{tabularx}{\linewidth}{|X|X|X|X|}
\hline
\rowcolor{popblueL} \textbf{Caractères} & \textbf{Pinyin} & \textbf{Radical (Clé)} & \textbf{Collocation type (V + O)} \\
\hline
公民 & gōngmín & 亻 (rén) & \zh{享有权利} (xiǎngyǒu quánlì) \\
\hline
选举 & xuǎnjǔ & 讠 (yán) & \zh{举行选举} (jǔxíng xuǎnjǔ) \\
\hline
权利 & quánlì & 木 (mù) / 力 (lì) & \zh{行使权利} (xíngshǐ quánlì) \\
\hline
义务 & yìwù & 礻 (shì) / 力 (lì) & \zh{履行义务} (lǚxíng yìwù) \\
\hline
候选人 & hòuxuǎnrén & 亻 (rén) & \zh{成为候选人} (chéngwéi hòuxuǎnrén) \\
\hline
\end{tabularx}
\end{center}

\textbf{Choix du mot juste :}
a) 每个 \zh{(公民 / 权利)} 都有投票的权利。 \\
b) 我们要 \zh{(遵守 / 遵守)} 法律。 (Vérifier le verbe qui colle avec \zh{法律}) $\rightarrow$ \zh{遵守法律}. \\
c) 他是一位优秀的 \zh{(候选人 / 选举)}。 \\

\tcblower
\textbf{Solution détaillée :}
\begin{enumerate}
    \item \textbf{Tableau complété :} (Voir ci-dessus). Notez le radical de \zh{权} (\zh{木} au-dessus de \zh{力}) et \zh{义} (\zh{礻/示} « autel/rite » + \zh{我}).
    \item \textbf{Choix justifiés :}
        \begin{itemize}
            \item a) \zh{公民} (nom humain). \zh{权利} est un nom abstrait, ne peut pas être sujet de \zh{每个...都有...} (chaque... a...).
            \item b) \zh{遵守} (verbe « respecter/observer » pour loi/règles). Collocation fixe \zh{遵守法律}.
            \item c) \zh{候选人} (nom « candidat »). \zh{选举} est nom « élection » ou verbe « élire », incompatible avec \zh{一位} (classificateur humain).
        \end{itemize}
    \item \textbf{Phrases produites :}
        \begin{itemize}
            \item \zh{作为公民，我们应该依法行使权利。} (Zuòwéi gōngmín, wǒmen yīnggāi yīfǎ xíngshǐ quánlì.) « En tant que citoyens, nous devons exercer nos droits selon la loi. »
            \item \zh{纳税是每个公民必须履行的义务。} (Nàshuì shì měi gè gōngmín bìxū lǚxíng de yìwù.) « Payer des impôts est un devoir que tout citoyen doit accomplir. »
        \end{itemize}
\end{enumerate}
\textbf{Conclusion :} La maîtrise des radicaux aide à deviner le sens ; les collocations évitent les erreurs de collocation verbale (ex: *\zh{做义务} $\rightarrow$ \zh{履行义务} ; *\zh{用权利} $\rightarrow$ \zh{行使权利}).
\end{exoresolu}

\begin{methode}[Production guidée : Mini-argumentaire écrit (Expression écrite)]
\textbf{Objectif :} Produire un paragraphe court (70-100 caractères chinois) sur la citoyenneté en utilisant le vocabulaire et les structures complexes du thème (niveau Bac).
\begin{enumerate}
    \item \textbf{Planification (Mind-map) :} Idée centrale : \zh{公民责任}. Branches : \zh{权利} (\zh{投票}、\zh{监督}), \zh{义务} (\zh{纳税}、\zh{遵法}、\zh{服兵役}), \zh{参与} (\zh{社区}、\zh{选举}).
    \item \textbf{Sélection des structures cibles (Bac HSK 4) :}
        \begin{itemize}
            \item \zh{不仅……而且……} (Progressivité, addition).
            \item \zh{只有……才……} (Condition nécessaire forte).
            \item \zh{随着……，……} (Corrélation temporelle).
            \item \zh{被} (Passif) : \zh{权利被保障}.
            \item \zh{是……的} (Insistance sur circonstance accomplie).
        \end{itemize}
    \item \textbf{Rédaction (Brouillon caractères) :} Écrire sans pinyin d'abord. Vérifier l'ordre des mots (STLVO).
    \item \textbf{Relecture \cle{Check-list Bac} :}
        \begin{itemize}
            \item [ ] Ponctuation chinoise (。 ， 、) ?
            \item [ ] Classificateurs corrects (\zh{一位}公民, \zh{一次}选举, \zh{项}义务, \zh{部}法律) ?
            \item [ ] \zh{的/得/地} distincts ? (\zh{积极地参与}, \zh{重要的权利}, \zh{做得好}) ?
            \item [ ] Tons notés sur le pinyin final (avec sandhi) ?
            \item [ ] Pas de calque français (ex: \zh{在雅温得投票} pas \zh{投票在雅温得}) ?
        \end{itemize}
    \item \textbf{Version finale :} Recopier proprement (Caractères + Pinyin + Traduction).
\end{enumerate}
\end{methode}

\begin{exoresolu}[Application : Mini-production écrite « Mon engagement citoyen »]
\textbf{Énoncé :} Rédigez un court paragraphe (environ 80-100 caractères chinois) intitulé \zh{我的公民责任}. Utilisez \textbf{au moins 6 mots du vocabulaire} de la leçon et \textbf{deux structures complexes} (\zh{不仅……而且……} / \zh{只有……才……}).

\tcblower
\textbf{Solution détaillée (Modèle de rédaction) :}

\begin{textebox}[Production élève type : 我的公民责任]
\zh{作为公民，我深知权利与义务统一。} (Zuòwéi gōngmín, wǒ shēnzhī quánlì yǔ yìwù tǒngyī.) \\
\zh{不仅要依法行使投票权，而且要自觉履行纳税、遵守法律的义务。} (Bùjǐn yào yīfǎ xíngshǐ tóupiàoquán, érqiě yào zìjué lǚxíng nàshuì, zūnshǒu fǎlǜ de yìwù.) \\
\zh{只有积极参与社区建设和选举，民主才能真正促进社会和谐发展。} (Zhǐyǒu jījí cānyù shèqū jiànshè hé xuǎnjǔ, mínzhǔ cáinéng zhēnzhèng cùjìn shèhuì xiéméi fāzhǎn.)
\end{textebox}

\textbf{Analyse linguistique du modèle :}
\begin{itemize}
    \item \textbf{Vocabulaire utilisé (10/6 requis) :} \zh{公民}、\zh{权利}、\zh{义务}、\zh{统一} (tǒngyī, unifié), \zh{依法}、\zh{行使}、\zh{投票权} (tóupiàoquán), \zh{自觉} (zìjué, consciemment), \zh{履行}、\zh{纳税} (nàshuì), \zh{遵守法律}、\zh{积极}、\zh{参与}、\zh{社区} (shèqū, communauté/quartier), \zh{建设}、\zh{选举}、\zh{民主}、\zh{真正}、\zh{促进}、\zh{社会}、\zh{和谐发展}.
    \item \textbf{Structure 1 :} \zh{不仅……而且……} (Phrase 2) : Lie deux obligations complémentaires (droit de vote + devoirs fiscaux/légaux). Sujet \zh{我} omis par ellipse (style écrit formel). \zh{要} (devoir) répété pour la parallélisme.
    \item \textbf{Structure 2 :} \zh{只有……才……} (Phrase 3) : Condition nécessaire. \zh{只有} + Condition (\zh{积极参与...}) + \zh{才} + Résultat (\zh{真正促进...}).
    \item \textbf{Grammaire :} \zh{依法} (adv. pré-verbal), \zh{自觉} (adv. pré-verbal), \zh{真正} (adv. de degré). \zh{的} correctement placé avant \zh{义务}.
    \item \textbf{Ponctuation :} Virgules chinoises pour énumération interne, points finaux.
\end{itemize}
\textbf{Traduction :} « En tant que citoyen, je sais bien que droits et devoirs sont unis. Il faut non seulement exercer son droit de vote selon la loi, mais aussi s'acquitter consciemment du devoir de payer des impôts et de respecter la loi. Ce n'est qu'en participant activement à la construction communautaire et aux élections que la démocratie pourra véritablement promouvoir le développement harmonieux de la société. »
\textbf{Conclusion :} Ce modèle obtient la note maximale car il respecte la contrainte de longueur, intègre le lexique ciblé naturellement, utilise les structures complexes exigées sans erreur d'ordre des mots, et présente une argumentation logique et cohérente.
\end{exoresolu}

\begin{exercice}[Entraînement personnel : Expression écrite]
\textbf{Sujet :} \zh{谈谈我对“公平、公正、公开”选举的理解} (Parlez de votre compréhension des élections « équitables, justes, transparentes »).
\textbf{Consignes :}
\begin{enumerate}
    \item Rédigez un paragraphe de 60-80 caractères.
    \item Utilisez \textbf{au moins 4 mots} du glossaire (ex: \zh{选举}、\zh{公民}、\zh{权利}、\zh{法律}、\zh{监督} jiāndū).
    \item Utilisez \textbf{une structure} au choix : \zh{是……的} (insistance) OU \zh{如果……就……} (condition).
    \item Écrivez : Caractères \rightarrow Pinyin (tons + sandhi) \rightarrow Traduction française.
\end{enumerate}
\end{exercice}

\begin{corrige}[Exercice : Expression écrite]
\textbf{Modèle de réponse :}
\begin{textebox}[Corrigé type]
\zh{选举是依法进行的，必须公平、公正、公开。} (Xuǎnjǔ shì yīfǎ jìnxíng de, bìxū gōngpíng, gōngzhèng, gōngkāi.) \\
\zh{如果选举不公正，公民的权利就得不到保障。} (Rúguǒ xuǎnjǔ bù gōngzhèng, gōngmín de quánlì jiù dé bù dào bǎozhàng.) \\
\zh{所以我们要依法监督每一次选举。} (Suǒyǐ wǒmen yào yīfǔ jiāndū měi yī cì xuǎnjǔ.)
\end{textebox}
\textbf{Analyse :} Phrase 1 : \zh{是……的} (insistance sur \zh{依法}). Phrase 2 : \zh{如果……就……} (condition/conséquence). Phrase 3 : \zh{所以} (conséquence), \zh{监督} (superviser - nouveau mot HSK 4), \zh{每一次} (classificateur \zh{次} pour événement).
\end{corrige}

\courssub{Points de vigilance pour le Bac}

\begin{attention}[Les 5 erreurs qui coûtent des points au Bac — Spécial « Démocratie \& Citoyenneté »]
\begin{enumerate}
    \item \textbf{Confusion \zh{权利} vs \zh{权力} (Mêmes tons \textit{quánlì}, sens opposés) :}
    \zh{权利} = Droit (juridique, \textbf{avantage} pour le citoyen) $\rightarrow$ Radical \zh{木} (bois/structure) + \zh{力}.
    \zh{权力} = Pouvoir (politique, \textbf{autorité} de l'État) $\rightarrow$ Radical \zh{力} (force).
    \textit{Piège Bac :} « Le gouvernement exerce ses \zh{权利} » $\rightarrow$ \textbf{FAUX}. Écrire : \zh{政府行使权力} / \zh{公民享有权利}.
    
    \item \textbf{Ordre des mots : Complément de lieu / temps AVANT le verbe (Calque français « Je vote \textbf{à Yaoundé} ») :}
    \textbf{FAUX :} \zh{我投票在雅温得。} \quad \textbf{JUSTE :} \zh{我在雅温得投票。} (Sujet + \zh{在} Lieu + Verbe).
    \textit{Règle d'or :} \cle{Topique $\rightarrow$ Temps $\rightarrow$ Lieu $\rightarrow$ Manière $\rightarrow$ Verbe}.
    
    \item \textbf{Mots-outils \zh{的 / 得 / 地} (Les 3 \textit{de} homophones) :}
    \begin{itemize}
        \item \zh{公平\textbf{的}选举} (Adj + \zh{的} + Nom) : « élection équitable ».
        \item \zh{选举进行得\textbf{很公平}} (Verbe + \zh{得} + Degré/Résultat) : « l'élection s'est déroulée très équitablement ».
        \item \zh{公平\textbf{地}投票} (Adv + \zh{地} + Verbe) : « voter équitablement ».
    \end{itemize}
    \textit{Erreur fréquente :} \zh{公平的投票} (Acceptable comme NP « vote équitable », mais \zh{公平地投票} requis pour l'action verbale).
    
    \item \textbf{Classificateurs omis ou erronés après nombres/démonstratifs :}
    \textbf{FAUX :} \zh{一候选人、这法律、三权利。}
    \textbf{JUSTE :} \zh{一\textbf{位}候选人} (humain respect), \zh{这\textbf{部}法律} (texte de loi/code officiel), \zh{三\textbf{项}权利} (articles/clauses listés), \zh{一次选举} (événement).
    \textit{Astuce :} \zh{位} (humain), \zh{部} (livres/documents officiels), \zh{项} (projets/articles), \zh{次} (fois/événement).
    
    \item \textbf{Tons « Sandhi » (变调) non appliqués à l'oral / Pinyin mal noté à l'écrit :}
    \begin{itemize}
        \item \zh{选举} (xuǎn\textbf{3} jǔ\textbf{3} $\rightarrow$ lu \textit{xuán\textbf{2} jǔ\textbf{3}}).
        \item \zh{民主} (mín\textbf{2} zhǔ\textbf{3}) -- pas de sandhi.
        \item \zh{一\textbf{位}} (yī\textbf{1} $\rightarrow$ yí\textbf{2} wèi / yì\textbf{4} wèi selon ton suivant).
        \item \zh{不\textbf{参加}} (bù\textbf{4} $\rightarrow$ bú\textbf{2} cānjiā).
        \item \zh{七\textbf{个}} (qī\textbf{1} $\rightarrow$ qí\textbf{2} gè / qì\textbf{4} gè).
    \end{itemize}
    \textit{Sanction :} Perte de points à l'oral (prononciation) et à l'écrit (pinyin incorrect = mot mal orthographié).
\end{enumerate}
\textbf{Conseil final :} Relisez votre copie en vous demandant : « Ai-je écrit \zh{权力} au lieu de \zh{权利} ? Mon verbe est-il bien \textbf{après} le lieu ? Mes \zh{的/得/地} sont-ils justifiés ? Mes classificateurs sont-ils présents ? Mes tons sandhi sont-ils notés correctement (ex: \zh{xuánjǔ}, \zh{yí wèi}, \zh{bú cānjiā}) ? »
\end{attention}

\begin{savaistu}[Regards croisés : Démocratie en Chine et au Cameroun]
\textbf{Chine (RPC) :} La Constitution définit la Chine comme un « État de dictature démocratique du peuple » (\zh{人民民主专政}). Le système politique repose sur le système des congrès du peuple (\zh{人民代表大会制度}) : les députés sont élus au suffrage universel direct au niveau local (base), et au suffrage indirect aux niveaux supérieurs. Le Parti communiste chinois (PCC) exerce le leadership (\zh{领导核心}). Le vocabulaire officiel insiste sur la « démocratie populaire à processus entier » (\zh{全过程人民民主}) : élection, consultation, décision, gestion, supervision.
\textbf{Cameroun :} République décentralisée unitaire. Démocratie multipartite depuis 1990. Président élu au suffrage universel direct (7 ans, renouvelable). Parlement bicaméral (Assemblée nationale + Sénat). Décentralisation avec des collectivités territoriales décentralisées (CTD) élues.
\textbf{Comparaison :} Les deux pays affirment la souveraineté populaire (\zh{权力属于人民} / « Le peuple est souverain »). Différence majeure : système de parti unique dirigeant (Chine) vs multipartisme compétitif (Cameroun). Point commun actuel : lutte contre la corruption, décentralisation/déconcentration, stabilité comme condition du développement.
\textbf{Vocabulaire bonus :} \zh{全过程人民民主} (démocratie populaire à processus entier), \zh{人大代表} (député à l'APN), \zh{政协} (Conférence consultative politique), \zh{多党制} (multipartisme), \zh{分权} (séparation des pouvoirs / décentralisation).
\end{savaistu}

\newpage

\coursec{Exercices}

\courssub{Je vérifie mes connaissances}

\begin{exercice}[QCM : Vocabulaire de la démocratie]
Chaque item ne comporte qu'une seule réponse correcte. Choisis la bonne proposition (A, B, C ou D).
\begin{enumerate}
    \item Quel terme désigne le « droit » (au sens légal, ex: droit de vote) ?
    \begin{itemize}
        \item[A.] 权力 (quánlì)
        \item[B.] 权利 (quánlì)
        \item[C.] 特权 (tèquán)
        \item[D.] 权威 (quánwēi)
    \end{itemize}
    \item Le verbe qui s'associe à \zh{选举} (élection) pour dire « organiser/tenir une élection » est :
    \begin{itemize}
        \item[A.] 开 (kāi)
        \item[B.] 举行 (jǔxíng)
        \item[C.] 做 (zuò)
        \item[D.] 参加 (cānjiā)
    \end{itemize}
    \item \zh{公民} (gōngmín) signifie :
    \begin{itemize}
        \item[A.] Gouvernement
        \item[B.] Citoyen
        \item[C.] Peuple
        \item[D.] Électeur
    \end{itemize}
    \item L'antonyme de \zh{自由} (zìyóu, liberté) dans le contexte des devoirs civiques est :
    \begin{itemize}
        \item[A.] 义务 (yìwù)
        \item[B.] 法律 (fǎlǜ)
        \item[C.] 限制 (xiànzhì)
        \item[D.] 民主 (mínzhǔ)
    \end{itemize}
    \item Pour exprimer « écouter les opinions » (\zh{听取意见}), on utilise le verbe :
    \begin{itemize}
        \item[A.] 听 (tīng)
        \item[B.] 听取 (tīngqǔ)
        \item[C.] 听说 (tīngshuō)
        \item[D.] 听见 (tīngjiàn)
    \end{itemize}
\end{enumerate}
\end{exercice}

\begin{exercice}[Vrai ou Faux : Justifiez par une phrase du texte]
Lis les affirmations suivantes. Coche \textbf{Vrai} ou \textbf{Faux} et recopie la phrase exacte du texte support qui justifie ta réponse.
\begin{enumerate}
    \item La démocratie signifie que le peuple exerce le pouvoir directement sans représentants.
    \item Les citoyens ont des droits mais n'ont pas d'obligations légales.
    \item L'élection du délégué de classe est un exemple concret de pratique démocratique à l'école.
    \item La loi protège la liberté d'expression des citoyens.
    \item En Chine, le système politique est une démocratie occidentale multipartite.
\end{enumerate}
\end{exercice}

\begin{exercice}[Vocabulaire : Associer caractères, pinyin et traduction]
Relie chaque caractère (colonne 1) à son pinyin correct (colonne 2) et à sa traduction (colonne 3) en traçant des flèches ou en recopiant dans l'ordre.
\begin{center}
\begin{tabularx}{\linewidth}{|X|X|X|X|}
\hline
\rowcolor{popblueL} \textbf{Caractères} & \textbf{Pinyin} & \textbf{Nature} & \textbf{Traduction} \\ \hline
民主 & mínzhǔ & n. & droit (légal) \\ \hline
选举 & gōngmín & v./n. & élection / élire \\ \hline
投票 & yìwù & v./n. & vote / voter \\ \hline
公民 & zìyóu & n. & citoyen \\ \hline
权利 & fǎlǜ & n. & liberté \\ \hline
义务 & quánlì & n. & devoir / obligation \\ \hline
自由 & tóupiào & n. & loi \\ \hline
意见 & mínzhǔ & n. & opinion / avis \\ \hline
负责 & jiǔxíng & v. & être responsable de \\ \hline
法律 & yìjiàn & n. & démocratie \\ \hline
\end{tabularx}
\end{center}
\emph{Note : Le tableau ci-dessus est mélangé intentionnellement. Classe correctement les 10 lignes.}
\end{exercice}

\begin{exercice}[Pinyin et tons : Dictée de mots du thème]
Écris les caractères chinois correspondants au pinyin suivant. Attention aux tons !
\begin{enumerate}
    \item mínzhǔ \hfill (\_\_\_\_\_\_\_\_\_\_)
    \item xuǎnjǔ \hfill (\_\_\_\_\_\_\_\_\_\_)
    \item tóupiào \hfill (\_\_\_\_\_\_\_\_\_\_)
    \item gōngmín \hfill (\_\_\_\_\_\_\_\_\_\_)
    \item quánlì \hfill (\_\_\_\_\_\_\_\_\_\_)
    \item yìwù \hfill (\_\_\_\_\_\_\_\_\_\_)
    \item zìyóu \hfill (\_\_\_\_\_\_\_\_\_\_)
    \item yìjiàn \hfill (\_\_\_\_\_\_\_\_\_\_)
    \item fùzé \hfill (\_\_\_\_\_\_\_\_\_\_)
    \item fǎlǜ \hfill (\_\_\_\_\_\_\_\_\_\_)
\end{enumerate}
\end{exercice}

\courssub{Je m'entraîne}

\begin{exercice}[Traduction : Français $\rightarrow$ Chinois]
Traduis les phrases suivantes en chinois (caractères + pinyin). Utilise le vocabulaire de la leçon.
\begin{enumerate}
    \item Voter est un droit et un devoir du citoyen.
    \item Notre école a organisé une élection la semaine dernière.
    \item Le délégué de classe doit écouter les opinions de tout le monde.
    \item La loi protège la liberté des citoyens.
    \item Nous devons assumer nos responsabilités.
\end{enumerate}
\end{exercice}

\begin{exercice}[Transformation de phrases : Structure \zh{是……的} (insistance sur le passé) et \zh{把}]
Transforme les phrases selon le modèle indiqué.
\begin{enumerate}
    \item \textbf{Modèle \zh{是……的} (action passée accomplie) :} \zh{我们上周举行了选举。} $\rightarrow$ \zh{选举是我们上周举行的。}
    \begin{itemize}
        \item[a.] \zh{他们昨天投了票。} $\rightarrow$ \zh{票是\_\_\_\_\_\_\_\_\_\_。}
        \item[b.] \zh{班长听取了同学们的意见。} $\rightarrow$ \zh{意见是\_\_\_\_\_\_\_\_\_\_。}
    \end{itemize}
    \item \textbf{Modèle \zh{把} (disposition de l'objet) :} \zh{请遵守法律。} $\rightarrow$ \zh{请把法律遵守。} (Note: \zh{遵守} ne prend pas toujours \zh{把} facilement, préférer verbes de résultat/manipulation). \textbf{Correction du modèle :} \zh{请把票投进箱子里。} (\zh{投票} $\rightarrow$ \zh{把票投进箱子里}).
    \begin{itemize}
        \item[a.] \zh{请填写选票。} $\rightarrow$ \zh{请把\_\_\_\_\_\_\_\_\_\_。}
        \item[b.] \zh{大家选举了新班长。} $\rightarrow$ \zh{大家把\_\_\_\_\_\_\_\_\_\_。}
    \end{itemize}
\end{enumerate}
\end{exercice}

\begin{exercice}[Textes à trous : Collocations verbales]
Complète les phrases avec le verbe approprié parmi la liste : \zh{举行 (jǔxíng)}, \zh{履行 (lǚxíng)}, \zh{听取 (tīngqǔ)}, \zh{行使 (xíngshǐ)}, \zh{遵守 (zūnshǒu)}. Mets le verbe à la forme convenable (éventuellement + \zh{了} / \zh{过}).
\begin{enumerate}
    \item 学校下周将\_\_\_\_\_\_\_\_\_\_新一届学生会选举。
    \item 公民必须\_\_\_\_\_\_\_\_\_\_法律规定的义务。
    \item 代表们认真\_\_\_\_\_\_\_\_\_\_了群众的意见和建议。
    \item 每个公民都有权\_\_\_\_\_\_\_\_\_\_自己的选举权和被选举权。
    \item 我们在日常生活中要自觉\_\_\_\_\_\_\_\_\_\_公共秩序。
\end{enumerate}
\end{exercice}

\begin{exercice}[Remise en ordre \& Appariements]
\begin{enumerate}
    \item \textbf{Remets les mots dans l'ordre pour former une phrase correcte.}
    \begin{itemize}
        \item[a.] 选举 / 我们 / 班长 / 刚才 / 举行 / 了 / 的
        \item[b.] 权利 / 和 / 义务 / 公民 / 是 / 不可分割 / 的
        \item[c.] 投票 / 箱子 / 请 / 选票 / 把 / 投进 / 里 / 面
    \end{itemize}
    \item \textbf{Apparie le mot-clé à sa définition contextuelle.}
    \begin{center}
    \begin{tabularx}{\linewidth}{|X|X|}
    \hline
    \rowcolor{popblueL} \textbf{Mot} & \textbf{Définition} \\ \hline
    1. 民主 & A. Document écrit qui fixe les règles de la société \\ \hline
    2. 法律 & B. Action de choisir ses représentants par vote \\ \hline
    3. 选举 & C. Système où le pouvoir appartient au peuple \\ \hline
    4. 义务 & D. Ce que l'on doit faire légalement ou moralement \\ \hline
    5. 意见 & E. Idée ou point de vue sur une question \\ \hline
    \end{tabularx}
    \end{center}
\end{enumerate}
\end{exercice}

\courssub{Je me prépare au Bac}

\begin{exercice}[Compréhension écrite type Baccalauréat — Texte original]
Lis le texte suivant (environ 6 lignes) et réponds aux questions en phrases complètes en chinois (caractères + pinyin).

\begin{textebox}[Élection du délégué de classe à Yaoundé]
我们班上周一举行了班长选举。这是一次真正的民主活动。全班四十五位同学都是公民，大家都有权利投票，也有义务参加。选举前，三位候选人分别介绍了自己的想法，并回答了同学们的提问。投票时，我们把选票投进箱子里，过程公开、公平、公正。最后，李明以三十票当选。他当选后表示：``我一定负责听取大家的意见，为同学服务。'' 这次选举让我深刻理解了民主和法治的意义。
\end{textebox}

\begin{enumerate}
    \item 这个选举是在什么时候举行的？(Quand a eu lieu l'élection ?)
    \item 参加选举的同学有多少人？Quelle est leur qualité (statut) ?(Combien d'élèves ont participé ? Quel est leur statut ?)
    \item 候选人在选举前做了什么准备？(Que ont fait les candidats avant le vote ?)
    \item 投票的过程体现了哪三个原则？(Le processus de vote a reflété quels trois principes ?)
    \item 李明当选后做出了什么承诺？(Quelle promesse Li Ming a-t-il faite après son élection ?)
\end{enumerate}
\end{exercice}

\begin{exercice}[Thème et Version : Traduction Bac]
\begin{enumerate}
    \item \textbf{Thème (Français $\rightarrow$ Chinois) :} Traduis les deux phrases suivantes en chinois. Soigne les structures (\zh{是……的}, \zh{把}, compléments de degré) et le vocabulaire précis.
    \begin{itemize}
        \item[a.] Voter est un droit sacré et un devoir légal pour chaque citoyen camerounais.
        \item[b.] La semaine dernière, notre lycée a organisé l'élection des délégués de manière transparente et équitable.
    \end{itemize}
    \item \textbf{Version (Chinois $\rightarrow$ Français) :} Traduis le paragraphe suivant en français courant et correct.
    \begin{textebox}[Extrait à traduire]
    民主不是装饰品，而是用来解决人民要解决的问题的。一个国家的民主不民主，关键在于是不是真正做到了人民当家作主。中国的全过程人民民主，保证了人民在选举、协商、决策、管理、监督等各个环节、各个领域享有广泛权利和自由。
    \end{textebox}
\end{enumerate}
\end{exercice}

\begin{exercice}[Expression écrite guidée : \zh{我们班的选举} (80--100 caractères)]
Rédige un court récit (80 à 100 caractères chinois, ponctuation comprise) pour raconter l'élection du délégué de ta classe.
\begin{itemize}
    \item \textbf{Contraintes obligatoires :} Utilise \textbf{au moins 5 mots} du glossaire de la leçon (ex: \zh{民主, 选举, 投票, 公民, 权利, 义务, 自由, 意见, 负责, 法律, 举行, 履行, 听取, 候选人, 当选, 公平, 公正}).
    \item \textbf{Structure suggérée :} Contexte (date, lieu) $\rightarrow$ Déroulement (candidats, vote, principes) $\rightarrow$ Résultat $\rightarrow$ Réaction / Sentiment personnel.
    \item \textbf{Indication de longueur :} Environ 6--8 phrases courtes.
\end{itemize}
\begin{textebox}[Cadre de rédaction]
\begin{center}
\begin{tabularx}{\linewidth}{|X|}
\hline
\zh{我们班的选举} \\ \hline
\\ \hline
\\ \hline
\\ \hline
\\ \hline
\\ \hline
\end{tabularx}
\end{center}
\end{textebox}
\end{exercice}

\newpage

\coursec{Corrigés des exercices}

\begin{corrige}[Exercice 1 — QCM : Vocabulaire de la démocratie]
\begin{enumerate}
    \item \textbf{Réponse B :} \zh{权利} (quánlì) = droit (légal, subjectif). \zh{权力} (quánlì) désigne le \emph{pouvoir} (autorité, pouvoir d'État) — homophone parfait, piège classique. \zh{特权} = privilège ; \zh{权威} = autorité (prestige).
    \item \textbf{Réponse B :} \zh{举行} (jǔxíng) = organiser, tenir (une cérémonie, une élection, une réunion). Collocation figée : \zh{举行选举}. \zh{参加} signifie « participer à », pas « organiser ».
    \item \textbf{Réponse B :} \zh{公民} (gōngmín) = citoyen (statut juridique). \zh{人民} = peuple ; \zh{政府} = gouvernement ; \zh{选民} = électeur.
    \item \textbf{Réponse A :} \zh{义务} (yìwù) = devoir, obligation. Dans le couple civique \zh{权利与义务} (droits et devoirs), c'est le terme qui s'oppose à \zh{自由} / \zh{权利}.
    \item \textbf{Réponse B :} \zh{听取} (tīngqǔ) = écouter attentivement, recueillir (registre soutenu) : \zh{听取意见} (écouter les avis). \zh{听说} = entendre dire ; \zh{听见} = entendre (perception) ; \zh{听} seul est trop faible pour cette collocation.
\end{enumerate}
\end{corrige}

\begin{corrige}[Exercice 2 — Vrai ou Faux]
\begin{enumerate}
    \item \textbf{Faux.} La démocratie moderne fonctionne par représentants élus. Justification attendue (type) : \zh{人民通过选举选出代表来行使权力。} (Rénmín tōngguò xuǎnjǔ xuǎnchū dàibiǎo lái xíngshǐ quánlì.) « Le peuple élit des représentants qui exercent le pouvoir. »
    \item \textbf{Faux.} Justification : \zh{公民既有权利，也有义务。} (Gōngmín jì yǒu quánlì, yě yǒu yìwù.) « Le citoyen a à la fois des droits et des devoirs. »
    \item \textbf{Vrai.} Justification : \zh{我们班上周一举行了班长选举。这是一次真正的民主活动。} (Wǒmen bān shàng zhōu yī jǔxíng le bānzhǎng xuǎnjǔ. Zhè shì yī cì zhēnzhèng de mínzhǔ huódòng.) « Notre classe a organisé l'élection du délégué lundi dernier. C'était une véritable activité démocratique. »
    \item \textbf{Vrai.} Justification : \zh{法律保护公民的言论自由。} (Fǎlǜ bǎohù gōngmín de yánlùn zìyóu.) « La loi protège la liberté d'expression des citoyens. »
    \item \textbf{Faux.} La Chine pratique \zh{全过程人民民主} (la démocratie populaire de processus complet) sous la direction du Parti communiste chinois, et non un système multipartite de type occidental. Justification : \zh{中国实行的是全过程人民民主。} (Zhōngguó shíxíng de shì quánguòchéng rénmín mínzhǔ.) « La Chine applique la démocratie populaire de processus complet. »
\end{enumerate}
\end{corrige}

\begin{corrige}[Exercice 3 — Associer caractères, pinyin et traduction]
\begin{center}
\begin{tabularx}{\linewidth}{|X|X|X|X|}
\hline
\rowcolor{popblueL} \textbf{Caractères} & \textbf{Pinyin} & \textbf{Nature} & \textbf{Traduction} \\ \hline
民主 & mínzhǔ & n. & démocratie \\ \hline
选举 & xuǎnjǔ & v./n. & élection / élire \\ \hline
投票 & tóupiào & v./n. & vote / voter \\ \hline
公民 & gōngmín & n. & citoyen \\ \hline
权利 & quánlì & n. & droit (légal) \\ \hline
义务 & yìwù & n. & devoir / obligation \\ \hline
自由 & zìyóu & n. & liberté \\ \hline
意见 & yìjiàn & n. & opinion / avis \\ \hline
负责 & fùzé & v. & être responsable de \\ \hline
法律 & fǎlǜ & n. & loi \\ \hline
\end{tabularx}
\end{center}
\textbf{Points de vigilance :} ne pas confondre \zh{选举} (xuǎnjǔ, 3\textsuperscript{e} + 3\textsuperscript{e} ton, avec changement de ton à l'oral : le premier 3\textsuperscript{e} ton devient un 2\textsuperscript{e} ton) avec un ton montant écrit ; \zh{负责} se prononce fùzé (4\textsuperscript{e} + 2\textsuperscript{e} ton). Attention à ne pas confondre avec \zh{举行} (jǔxíng, 3\textsuperscript{e} + 2\textsuperscript{e} ton) : les deux mots n'ont ni le même sens ni la même prononciation.
\end{corrige}

\begin{corrige}[Exercice 4 — Dictée de mots du thème]
\begin{enumerate}
    \item mínzhǔ $\rightarrow$ \zh{民主}
    \item xuǎnjǔ $\rightarrow$ \zh{选举}
    \item tóupiào $\rightarrow$ \zh{投票}
    \item gōngmín $\rightarrow$ \zh{公民}
    \item quánlì $\rightarrow$ \zh{权利} (attention : pas \zh{权力} !)
    \item yìwù $\rightarrow$ \zh{义务}
    \item zìyóu $\rightarrow$ \zh{自由}
    \item yìjiàn $\rightarrow$ \zh{意见}
    \item fùzé $\rightarrow$ \zh{负责}
    \item fǎlǜ $\rightarrow$ \zh{法律}
\end{enumerate}
\textbf{Conseil méthode :} en dictée, repère d'abord les tons : 4\textsuperscript{e} ton + 4\textsuperscript{e} ton (\zh{义务}, \zh{意见}) ; deux 3\textsuperscript{e} tons (\zh{选举}) ; 1\textsuperscript{er} + 1\textsuperscript{er} (\zh{公民}). Vérifie ensuite les caractères proches : \zh{权} (clé 木, le bois) dans \zh{权利}, \zh{义} (3 traits) dans \zh{义务}, et ne confonds pas \zh{律} (clé 彳) avec \zh{津}.
\end{corrige}

\begin{corrige}[Exercice 5 — Traduction : Français $\rightarrow$ Chinois]
\begin{enumerate}
    \item \zh{投票是公民的权利和义务。} (Tóupiào shì gōngmín de quánlì hé yìwù.) Structure : Sujet (\zh{投票}) + \zh{是} + complément du nom avec \zh{的}.
    \item \zh{我们学校上周举行了选举。} (Wǒmen xuéxiào shàng zhōu jǔxíng le xuǎnjǔ.) Le complément de temps \zh{上周} se place avant le verbe ; \zh{了} marque l'accomplissement.
    \item \zh{班长必须听取大家的意见。} (Bānzhǎng bìxū tīngqǔ dàjiā de yìjiàn.) « doit » = \zh{必须} ou \zh{应该} ; \zh{大家} = tout le monde.
    \item \zh{法律保护公民的自由。} (Fǎlǜ bǎohù gōngmín de zìyóu.)
    \item \zh{我们必须承担自己的责任。} (Wǒmen bìxū chéngdān zìjǐ de zérèn.) On peut aussi dire : \zh{我们要履行自己的义务。} (Wǒmen yào lǚxíng zìjǐ de yìwù.) « Nous devons remplir nos devoirs. »
\end{enumerate}
\textbf{Points de vigilance :} pas de \zh{了} après \zh{必须} + verbe (action non accomplie) ; l'ordre des mots reste Sujet + Temps + Verbe + Objet.
\end{corrige}

\begin{corrige}[Exercice 6 — Transformation : \zh{是……的} et \zh{把}]
\begin{enumerate}
    \item \textbf{Modèle \zh{是……的} :}
    \begin{itemize}
        \item[a.] \zh{票是他们昨天投的。} (Piào shì tāmen zuótiān tóu de.) « C'est hier qu'ils ont voté. »
        \item[b.] \zh{意见是班长听取的。} (Yìjiàn shì bānzhǎng tīngqǔ de.) « C'est le délégué qui a écouté les avis. »
    \end{itemize}
    Règle : Objet + \zh{是} + Sujet + Temps + Verbe + \zh{的}. On insiste sur les circonstances (qui, quand, comment) d'une action passée ; pas de \zh{了} dans cette structure.
    \item \textbf{Modèle \zh{把} :}
    \begin{itemize}
        \item[a.] \zh{请把选票填写好。} (Qǐng bǎ xuǎnpiào tiánxiě hǎo.) « Remplissez bien le bulletin. » Le verbe après \zh{把} doit être suivi d'un complément de résultat (\zh{好}) — jamais un verbe nu.
        \item[b.] \zh{大家把新班长选出来了。} (Dàjiā bǎ xīn bānzhǎng xuǎn chūlai le.) « Tout le monde a élu le nouveau délégué. » Complément directionnel/résultatif \zh{出来} obligatoire ; on préfère ici le verbe simple \zh{选} à \zh{选举} devant ce complément.
    \end{itemize}
\end{enumerate}
\textbf{Attention :} avec \zh{把}, l'objet doit être \emph{défini} (connu) et le verbe ne peut jamais rester seul : il faut \zh{了}, un complément de résultat ou de direction.
\end{corrige}

\begin{corrige}[Exercice 7 — Textes à trous : Collocations verbales]
\begin{enumerate}
    \item \zh{学校下周将举行新一届学生会选举。} (Xuéxiào xià zhōu jiāng jǔxíng xīn yī jiè xuéshēnghuì xuǎnjǔ.) Collocation figée : \zh{举行选举}.
    \item \zh{公民必须履行法律规定的义务。} (Gōngmín bìxū lǚxíng fǎlǜ guīdìng de yìwù.) \zh{履行义务} = remplir ses devoirs (registre soutenu).
    \item \zh{代表们认真听取了群众的意见和建议。} (Dàibiǎomen rènzhēn tīngqǔ le qúnzhòng de yìjiàn hé jiànyì.) \zh{了} car l'action est accomplie.
    \item \zh{每个公民都有权行使自己的选举权和被选举权。} (Měi ge gōngmín dōu yǒu quán xíngshǐ zìjǐ de xuǎnjǔquán hé bèi xuǎnjǔquán.) \zh{行使权利} = exercer un droit ; \zh{有权} + verbe = avoir le droit de.
    \item \zh{我们在日常生活中要自觉遵守公共秩序。} (Wǒmen zài rìcháng shēnghuó zhōng yào zìjué zūnshǒu gōnggòng zhìxù.) \zh{遵守} + loi/règle/ordre.
\end{enumerate}
\textbf{Mémo des collocations :} \zh{举行} + activité ; \zh{履行} + devoir ; \zh{听取} + avis ; \zh{行使} + droit ; \zh{遵守} + loi/règle.
\end{corrige}

\begin{corrige}[Exercice 8 — Remise en ordre et appariements]
\begin{enumerate}
    \item \textbf{Remise en ordre :}
    \begin{itemize}
        \item[a.] \zh{我们刚才举行了班长的选举。} (Wǒmen gāngcái jǔxíng le bānzhǎng de xuǎnjǔ.) « Nous venons d'organiser l'élection du délégué. » Ordre : Sujet + temps (\zh{刚才}) + verbe + \zh{了} + objet déterminé par \zh{的}.
        \item[b.] \zh{公民的权利和义务是不可分割的。} (Gōngmín de quánlì hé yìwù shì bùkě fēngē de.) « Les droits et les devoirs du citoyen sont indissociables. » Structure \zh{是……的} avec adjectif.
        \item[c.] \zh{请把选票投进箱子里面。} (Qǐng bǎ xuǎnpiào tóu jìn xiāngzi lǐmiàn.) « Veuillez mettre le bulletin dans l'urne. » Structure \zh{把} + objet + verbe + complément de direction \zh{进} + lieu.
    \end{itemize}
    \item \textbf{Appariements :} 1--C (\zh{民主} = système où le pouvoir appartient au peuple) ; 2--A (\zh{法律} = règles écrites de la société) ; 3--B (\zh{选举} = choix de représentants par vote) ; 4--D (\zh{义务} = ce qu'on doit faire) ; 5--E (\zh{意见} = point de vue).
\end{enumerate}
\end{corrige}

\begin{corrige}[Exercice 9 — Compréhension écrite type Baccalauréat]
\textbf{Réponses modèles (phrases complètes) :}
\begin{enumerate}
    \item \zh{选举是上周一举行的。} (Xuǎnjǔ shì shàng zhōu yī jǔxíng de.) « L'élection a eu lieu lundi dernier. » (Structure \zh{是……的} pour insister sur la date.)
    \item \zh{参加选举的有四十五位同学，大家都是公民。} (Cānjiā xuǎnjǔ de yǒu sìshíwǔ wèi tóngxué, dàjiā dōu shì gōngmín.) « Quarante-cinq élèves ont participé ; tous sont des citoyens. »
    \item \zh{选举前，三位候选人分别介绍了自己的想法，并回答了同学们的提问。} (Xuǎnjǔ qián, sān wèi hòuxuǎnrén fēnbié jièshào le zìjǐ de xiǎngfǎ, bìng huídá le tóngxuémen de tíwèn.) « Avant le vote, les trois candidats ont présenté leurs idées et répondu aux questions des élèves. »
    \item \zh{投票的过程公开、公平、公正。} (Tóupiào de guòchéng gōngkāi, gōngpíng, gōngzhèng.) « Le processus était transparent, équitable et juste. »
    \item \zh{他表示一定负责听取大家的意见，为同学服务。} (Tā biǎoshì yídìng fùzé tīngqǔ dàjiā de yìjiàn, wèi tóngxué fúwù.) « Il a promis d'écouter les avis de tous et de servir ses camarades. »
\end{enumerate}
\textbf{Barème indicatif (10 pts) :} 2 pts par question (1 pt exactitude du contenu, 1 pt correction de la langue : caractères, ordre des mots, tons du pinyin).
\textbf{Points de vigilance :} répondre par des \emph{phrases complètes} (sujet + verbe) ; ne pas recopier le texte sans adapter ; \zh{位} est le classificateur de politesse pour les personnes.
\end{corrige}

\begin{corrige}[Exercice 10 — Thème et Version : Traduction Bac]
\textbf{1. Thème (propositions de traduction) :}
\begin{itemize}
    \item[a.] \zh{投票是每个喀麦隆公民神圣的权利和法定的义务。} (Tóupiào shì měi ge Kāmàilóng gōngmín shénshèng de quánlì hé fǎdìng de yìwù.) « Sacré » = \zh{神圣} ; « légal » = \zh{法定}. Variante acceptée : \zh{投票既是每个公民神圣的权利，也是法定的义务。} (structure \zh{既……也……}).
    \item[b.] \zh{上周，我们学校以透明、公平的方式举行了班长选举。} (Shàng zhōu, wǒmen xuéxiào yǐ tòumíng, gōngpíng de fāngshì jǔxíng le bānzhǎng xuǎnjǔ.) « de manière … » = \zh{以……的方式} placé avant le verbe ; \zh{了} pour l'action accomplie.
\end{itemize}
\textbf{2. Version (traduction modèle) :}

« La démocratie n'est pas un objet de décoration, mais un moyen de résoudre les problèmes que le peuple doit résoudre. Pour savoir si un pays est démocratique ou non, l'essentiel est de déterminer si le peuple y est véritablement maître du pays. La démocratie populaire de processus complet de la Chine garantit que le peuple jouit de droits et de libertés étendus dans tous les maillons et tous les domaines : élections, concertation, prise de décision, gestion et contrôle. »

\textbf{Barème indicatif (10 pts) :} Thème 5 pts (2,5 pts par phrase : vocabulaire 1 pt, structure 1 pt, caractères 0,5 pt) ; Version 5 pts (fidélité 3 pts, qualité du français 2 pts).
\textbf{Points de vigilance :} \zh{装饰品} = objet décoratif ; \zh{当家作主} = être maître chez soi (expression figée à traduire par « maître du pays ») ; \zh{环节} = maillon, étape. Ne pas traduire mot à mot \zh{全过程人民民主} : accepter « démocratie populaire de processus complet ».
\end{corrige}

\begin{textebox}[Modèle de rédaction — \zh{我们班的选举}]
\zh{上周五，我们班举行了班长选举。} (Shàng zhōu wǔ, wǒmen bān jǔxíng le bānzhǎng xuǎnjǔ.)\\
\zh{三位候选人先介绍了自己的想法。} (Sān wèi hòuxuǎnrén xiān jièshào le zìjǐ de xiǎngfǎ.)\\
\zh{然后，每个同学都行使了投票的权利。} (Ránhòu, měi ge tóngxué dōu xíngshǐ le tóupiào de quánlì.)\\
\zh{我们把选票投进箱子里，过程公平、公正。} (Wǒmen bǎ xuǎnpiào tóu jìn xiāngzi lǐ, guòchéng gōngpíng, gōngzhèng.)\\
\zh{最后，玛丽当选了。} (Zuìhòu, Mǎlì dāngxuǎn le.)\\
\zh{她说：“我一定负责听取大家的意见。”} (Tā shuō: “Wǒ yídìng fùzé tīngqǔ dàjiā de yìjiàn.”)\\
\zh{这次活动让我明白了民主的意义。} (Zhè cì huódòng ràng wǒ míngbai le mínzhǔ de yìyì.)
\end{textebox}

\begin{corrige}[Exercice 11 — Expression écrite guidée : \zh{我们班的选举}]
\textbf{Proposition de rédaction modèle :} voir le texte encadré ci-dessus (environ 95 caractères).

\textbf{Traduction :} « Vendredi dernier, notre classe a organisé l'élection du délégué. Les trois candidats ont d'abord présenté leurs idées. Ensuite, chaque élève a exercé son droit de vote. Nous avons mis nos bulletins dans l'urne ; le processus était équitable et juste. Finalement, Marie a été élue. Elle a dit : "J'écouterai avec responsabilité les avis de tout le monde." Cette activité m'a fait comprendre le sens de la démocratie. »

\textbf{Barème indicatif (10 pts) :} respect de la longueur 80--100 caractères (1 pt) ; au moins 5 mots du glossaire correctement employés (3 pts) ; correction des caractères (2 pts) ; grammaire et ordre des mots (2 pts) ; cohérence du récit en 4 étapes (2 pts).

\textbf{Points de vigilance :}
\begin{itemize}
    \item Mots du glossaire utilisés ici : \zh{举行, 选举, 候选人, 投票, 权利, 公平, 公正, 当选, 负责, 听取, 意见, 民主} (12 mots — largement suffisant).
    \item Ne pas oublier \zh{了} après les verbes d'action accomplie (\zh{举行了, 当选了}).
    \item Structure \zh{把} : \zh{把选票投进箱子里} — le verbe n'est jamais nu.
    \item Éviter de calquer le français : pas de « \zh{是民主的} » pour « c'était démocratique » ; préférer \zh{这是一次民主活动}.
\end{itemize}
\end{corrige}

\newpage

\coursec{Fiche bilan}

\courssub{Fiche bilan — Leçon 21 : Démocratie}

\begin{popfigure}
\begin{tikzpicture}[
  mindmap/.style={concept, execute at begin node=\hskip0pt, font=\small\bfseries},
  level 1 concept/.append style={sibling angle=60, level distance=4.5cm, font=\small},
  level 2 concept/.append style={level distance=3cm, font=\tiny},
  every node/.style={align=center, inner sep=3pt},
  concept color=popblue,
  branch1/.style={concept color=popblue},
  branch2/.style={concept color=poporange},
  branch3/.style={concept color=popgreen},
  branch4/.style={concept color=poppurple},
  branch5/.style={concept color=poppink},
  branch6/.style={concept color=popgold},
]
\node[mindmap, align=center]{Démocratie\\ \zh{民主} (mínzhǔ)}
  child[branch1] {node[mindmap, align=center]{Vocabulaire\\ clé\\ (15 mots)}}
  child[branch2] {node[mindmap, align=center]{Texte support\\ \zh{民主的实践}\\ Lecture guidée}}
  child[branch3] {node[mindmap, align=center]{Radicaux \\ \& Ordre\\ des traits\\ \zh{民 主 选 票}}}
  child[branch4] {node[mindmap, align=center]{Compréhension\\ QCM + Questions\\ ouvertes}}
  child[branch5] {node[mindmap, align=center]{Collocations\\ \zh{行使权利}\\ \zh{参加选举}}}
  child[branch6] {node[mindmap, align=center]{Contexte\\ Cameroun /\\ Chine\\ \zh{全过程人民民主}}};
\end{tikzpicture}
\legende{Carte mentale récapitulative de la leçon 21 — Module 3 : Citoyenneté et ouverture au monde}
\end{popfigure}

\begin{bilan}

\textbf{Lexique clé (chinois $\leftrightarrow$ français)}

\begin{center}
\begin{tabularx}{\linewidth}{|X|X|X|X|}
\hline
\rowcolor{popblueL} \textbf{Caractères} & \textbf{Pinyin} & \textbf{Nature} & \textbf{Traduction} \\ \hline
民主 & mínzhǔ & n. & démocratie \\ \hline
选举 & xuǎnjǔ & v./n. & élection ; élire \\ \hline
投票 & tóupiào & v./n. & vote ; voter \\ \hline
权利 & quánlì & n. & droit (juridique) \\ \hline
义务 & yìwù & n. & devoir, obligation \\ \hline
公民 & gōngmín & n. & citoyen \\ \hline
法律 & fǎlǜ & n. & loi \\ \hline
宪法 & xiànfǎ & n. & constitution \\ \hline
政府 & zhèngfǔ & n. & gouvernement \\ \hline
代表 & dàibiǎo & n./v. & représentant ; représenter \\ \hline
监督 & jiāndū & v. & superviser, contrôler \\ \hline
透明 & tòumíng & adj. & transparent \\ \hline
公平 & gōngpíng & adj. & juste, équitable \\ \hline
参与 & cānyù & v. & participer (à) \\ \hline
决策 & juécè & n./v. & décision ; décider \\ \hline
\end{tabularx}
\end{center}

\textbf{Structures grammaticales à connaître par cœur}

\begin{center}
\begin{tabularx}{\linewidth}{|X|X|X|}
\hline
\rowcolor{poporangeL} \textbf{Structure} & \textbf{Exemple (caractères + pinyin)} & \textbf{Traduction} \\ \hline
通过 + Nom/Verbe, ... & \zh{通过选举，公民行使权利。} (Tōngguò xuǎnjǔ, gōngmín xíngshǐ quánlì.) & Par l'élection, les citoyens exercent leurs droits. \\ \hline
保障 + Objet & \zh{法律保障公民的投票权。} (Fǎlǜ bǎozhàng gōngmín de tóupiàoquán.) & La loi garantit le droit de vote des citoyens. \\ \hline
体现 & \zh{民主体现在决策过程中。} (Mínzhǔ tǐxiàn zài juécè guòchéng zhōng.) & La démocratie se reflète dans le processus de décision. \\ \hline
不仅...而且/也... & \zh{公民不仅有权利，而且有义务。} (Gōngmín bùjǐn yǒu quánlì, érqiě yǒu yìwù.) & Les citoyens ont non seulement des droits, mais aussi des devoirs. \\ \hline
使/让/令 + Quelqu'un + Verbe & \zh{透明的制度让政府更负责。} (Tòumíng de zhìdù ràng zhèngfǔ gèng fùzé.) & Un système transparent rend le gouvernement plus responsable. \\ \hline
\end{tabularx}
\end{center}

\textbf{Méthodes express}

\begin{itemize}
  \item \textbf{Lecture active :} 1) Identifier les mots-clés (titres, mots en gras, radicaux connus). 2) Segmenter les phrases (repérer les virgules chinoises ， et les points 。). 3) Deviner le sens des inconnus par le contexte (collocations, radicaux sémantiques).
  \item \textbf{Compréhension :} Distinguer \textbf{信息显性} (info explicite : qui, quoi, quand, où) et \textbf{信息隐性} (info implicite : pourquoi, opinion de l'auteur, conséquence).
  \item \textbf{Rédaction courte :} Utiliser la structure \textbf{Sujet + 通过 + Moyen + 动词 + Objet} pour expliquer un mécanisme démocratique. Ex: \zh{我们通过投票选出代表。}
  \item \textbf{Oral / Débat :} Préparer 3 arguments avec connecteurs logiques : \zh{首先}... \zh{其次}... \zh{最后}... ; conclure par \zh{总之}...
\end{itemize}

\end{bilan}

\begin{aretenir}[Checklist avant le Bac]
$\square$ Je sais écrire de mémoire les 15 mots du lexique clé (caractères simplifiés, ordre des traits correct). \\
$\square$ Je reconnais et prononce correctement les 4 tons de chaque mot (pinyin sans hésitation). \\
$\square$ Je sais identifier les radicaux \zh{文} (wén), \zh{示} (shì/礻), \zh{手} (shǒu/扌), \zh{言} (yán/讠) dans les mots du thème. \\
$\square$ Je sais transformer \zh{选举} (verbe) en \zh{选举} (nom) et comprendre \zh{投票} comme action et résultat. \\
$\square$ Je maîtrise la structure \zh{通过...} (par le biais de...) pour exprimer le moyen. \\
$\square$ Je sais utiliser \zh{保障} (garantir) avec un objet droit (droit, liberté, sécurité). \\
$\square$ Je distingue \zh{权利} (droit) de \zh{义务} (devoir) et je peux citer un exemple de chaque pour un citoyen camerounais/chinois. \\
$\square$ Je peux répondre en chinois simple à : \zh{什么是民主？} (Qu'est-ce que la démocratie ?). \\
$\square$ Je connais la différence terminologique : \zh{全过程人民民主} (démocratie populaire tout au long du processus) vs \zh{西式民主} (démocratie de style occidental). \\
$\square$ Je sais rédiger un paragraphe de 5-6 phrases sur \zh{我的公民权利和义务} (mes droits et devoirs civiques) avec connecteurs logiques. \\
$\square$ Je relis mon texte pour vérifier : ordre des mots (Sujet - Temps - Lieu - Manière - Verbe - Objet), place des compléments, ponctuation chinoise.
\end{aretenir}

\findecours

\end{document}
